% Options for packages loaded elsewhere
\PassOptionsToPackage{unicode}{hyperref}
\PassOptionsToPackage{hyphens}{url}
%
\documentclass[11pt,oneside]{book}
\usepackage{amsmath,amssymb,mathtools}
\usepackage[margin=1in]{geometry}
\usepackage{iftex}
\ifPDFTeX
  \usepackage[T1]{fontenc}
  \usepackage[utf8]{inputenc}
  \usepackage{textcomp} % provide euro and other symbols
\else % if luatex or xetex
  \usepackage{unicode-math} % this also loads fontspec
  \defaultfontfeatures{Scale=MatchLowercase}
  \defaultfontfeatures[\rmfamily]{Ligatures=TeX,Scale=1}
\fi
\usepackage{lmodern}
\ifPDFTeX\else
  % xetex/luatex font selection
  \setmainfont{DejaVu Serif}
\fi
% Use upquote if available, for straight quotes in verbatim environments
\IfFileExists{upquote.sty}{\usepackage{upquote}}{}
\IfFileExists{microtype.sty}{% use microtype if available
  \usepackage[]{microtype}
  \UseMicrotypeSet[protrusion]{basicmath} % disable protrusion for tt fonts
}{}
\makeatletter
\@ifundefined{KOMAClassName}{% if non-KOMA class
  \IfFileExists{parskip.sty}{%
    \usepackage{parskip}
  }{% else
    \setlength{\parindent}{0pt}
    \setlength{\parskip}{6pt plus 2pt minus 1pt}}
}{% if KOMA class
  \KOMAoptions{parskip=half}}
\makeatother
\usepackage{xcolor}
\usepackage{longtable,booktabs,array}
\usepackage{calc} % for calculating minipage widths
% Correct order of tables after \paragraph or \subparagraph
\usepackage{etoolbox}
\makeatletter
\patchcmd\longtable{\par}{\if@noskipsec\mbox{}\fi\par}{}{}
\makeatother
% Allow footnotes in longtable head/foot
\IfFileExists{footnotehyper.sty}{\usepackage{footnotehyper}}{\usepackage{footnote}}
\makesavenoteenv{longtable}
\setlength{\emergencystretch}{3em} % prevent overfull lines
\providecommand{\tightlist}{%
  \setlength{\itemsep}{0pt}\setlength{\parskip}{0pt}}
\setcounter{secnumdepth}{3}

% Wider number boxes in the table of contents so entries such as 10.10
% do not run into their section titles.
\makeatletter
\renewcommand*\l@section{\@dottedtocline{1}{1.5em}{3.8em}}
\renewcommand*\l@subsection{\@dottedtocline{2}{5.3em}{4.5em}}
\makeatother
\sloppy
\pagestyle{plain}
\ifLuaTeX
  \usepackage{selnolig}  % disable illegal ligatures
\fi
\usepackage{bookmark}
\IfFileExists{xurl.sty}{\usepackage{xurl}}{} % add URL line breaks if available
\urlstyle{same}
\hypersetup{
  hidelinks,
  pdfcreator={LaTeX via pandoc}}

% IRFR convenience macros
\newcommand{\IRFR}{\mathrm{IRFR}}
\newcommand{\dd}{\,\mathrm{d}}
\newcommand{\HN}{$H_N$}
\newcommand{\HD}{$H_D$}
\newcommand{\HR}{$H_R$}


\title{\textbf{IRFR Monograph}\\\large A Boundary-Field Theory of Interest Rates}
\author{}
\date{}

\begin{document}
\frontmatter
\maketitle

{
\setcounter{tocdepth}{2}
\tableofcontents
}
\mainmatter
\chapter{Introduction and Programme}
\section{Motivation}\label{motivation}

Interest-rate modelling has developed through several powerful but
partially distinct traditions. Short-rate models provide intuitive state
variables and tractable dynamics. Heath--Jarrow--Morton (HJM)-type
models provide a broad arbitrage-free representation of forward-rate
evolution. Affine term-structure models impose analytical structure and
calibration discipline. Market models focus directly on traded rates and
market observables.

Each of these traditions has been extraordinarily productive, yet a
conceptual separation remains. In many models, the dynamic consistency
of the term structure is handled in one part of the framework, while the
shape of the curve, the maturity profile of volatility, and the effect
of monetary policy are introduced elsewhere through additional
parametrisation. Humped volatility structures, curve-shape restrictions,
maturity loadings, and policy effects are often specified because they
are useful or empirically plausible, rather than derived from a common
structural mechanism.

The empirical importance of humped volatility makes this separation
especially visible. Ritchken--Chuang-type maturity profiles are useful
precisely because they capture an observed concentration of volatility
away from the shortest and longest maturities. Quadratic term-structure
models and other nonlinear extensions provide one powerful response to
this problem by enriching the state representation. The question pursued
here is different: can some of the same empirically recognisable
maturity structure be explained before adding richer polynomial state
machinery, by first exploiting rolling-maturity geometry and the
boundary at the short end?

The purpose of this monograph is to investigate whether more of this
structure can be generated from fewer principles, while preserving a
clear account of the modelling cost. The aim is not merely parsimony,
but internal consistency: curve shape, volatility structure, maturity
loading, and policy anchoring should be linked through the same
governing field principles, so that every modelling restriction has an
identifiable location in the theory.

The central object is the instantaneous rolling forward rate. Instead of
taking the forward rate at calendar time \(t\) for fixed maturity \(T\) as the
primitive object, the IRFR formulation works in rolling-maturity
coordinates. Writing

\[
x(t,\tau) = F(t,t+\tau), \tau = T - t,
\]

the forward curve becomes a stochastic field over remaining maturity $\tau$.
Calendar time evolves the field, while the maturity coordinate rolls
deterministically toward the short end. The point $\tau=0$ is therefore not
merely a limiting maturity. It is the boundary of the term-structure
domain.

This change of primitive is not just a change of notation. It changes
the geometry of the modelling problem. The curve is no longer viewed
primarily as a family of unrelated forward-rate processes indexed by
maturity. It is viewed as a maturity-domain field transported through
time. Its admissible behaviour is constrained by stochastic evolution,
rolling-maturity consistency, equilibrium moment transport, and boundary
closure at the short end.

This is the sense in which the present work develops a field theory of
interest rates. The phrase is not intended as a metaphor imported from
physics for decorative effect. It refers to the treatment of the term
structure as a stochastic object defined over a continuous maturity
domain, with an interior law, admissible modes, and boundary conditions.
The field is the forward-rate curve in rolling-maturity coordinates. The
boundary is the short end. The boundary law is the rule that closes the
system.

The guiding claim of the monograph is that a substantial part of
term-structure behaviour can be reorganised around this structure. Curve
shape, maturity loading, volatility localisation, and monetary-policy
anchoring need not be introduced as separate modelling interventions.
They can be interpreted as consequences of the same rolling-field
geometry and its boundary laws.

\section{The central principle}\label{the-central-principle}

The organising principle of the monograph is:

Dynamics are determined by invariance; observable shapes are selected by
boundary laws.

This sentence is deliberately schematic, so as to capture the
architecture of the theory.

The first part of the statement concerns the interior of the maturity
domain. Rolling-maturity invariance and equilibrium moment transport
restrict the admissible stochastic dynamics of the IRFR field. The
Kolmogorov forward equation and affine closure then force a particular
relation between drift, diffusion, and maturity loading. The result is
not an arbitrary volatility specification but a structured family of
admissible maturity modes.

The second part of the statement concerns the boundary at $\tau=0$. The
interior dynamics determine what the field is allowed to do, but they do
not by themselves determine which admissible mode is realised. That
selection is made by the boundary law. A Neumann/free\footnote{The label Neumann/free denotes the non-pinned benchmark in which the boundary-active exponential short-end loading is retained. It should not be read as asserting that the operative loading benchmark is the classical homogeneous condition $H'(0)=0$. The operative content is non-pinning: $H_N(0)=A_N$, so the short end remains stochastic.} boundary, a fixed
boundary, and a partially anchored boundary select different maturity
structures, and therefore different observable curve and volatility
shapes.

This separation is central. The theory distinguishes three objects that
are often conflated in term-structure modelling:

\begin{enumerate}
\def\labelenumi{\arabic{enumi}.}
\item
  the stochastic dimension of the model;
\item
  the deterministic maturity geometry of the curve;
\item
  the boundary law that selects admissible modes.
\end{enumerate}

The first object concerns the source of randomness. The second concerns
the shape of the maturity loading through which that randomness is
expressed across the curve. The third concerns the economic regime
imposed at the short end.

This distinction explains how a one-factor stochastic model can still
generate non-trivial term-structure shapes. A model may be one-factor in
randomness while having a richer deterministic maturity geometry. In the
IRFR framework, the stochastic dimension and the maturity-mode structure
are not the same object. The covariance may be rank-one, while the
admissible curve shapes may still include multiple deterministic modes,
repeated-root structures, and boundary-induced humps.

The consequence is a different interpretation of model flexibility.
Flexibility is not obtained by freely specifying more volatility
functions. It arises from the admissible interaction between the
rolling-field interior and the boundary condition at the short end.

\section{The IRFR primitive}\label{the-irfr-primitive}

The modelling primitive is the instantaneous rolling forward rate

\[
x(t,\tau) = F(t,t+\tau), \tau \ge 0
\]

Here \(t\) denotes calendar time and $\tau$ denotes remaining maturity. For a
fixed calendar maturity T, remaining maturity satisfies

\[
\tau = T - t, d\tau/dt = -1
\]

Thus, as calendar time advances, every fixed-maturity contract rolls
toward the short end. The field evolves in two ways at once: it is
shocked stochastically in calendar time, and it is transported
deterministically along the maturity axis.

This transport is not an optional modelling feature. It is a geometric
fact of the rolling-maturity representation. Any admissible drift must
account for it. A pure mean-reversion term is not enough, because the
equilibrium curve itself is being traversed as maturity rolls down. A
transport correction is required to preserve consistency between
stochastic evolution and the movement of the maturity coordinate.

This is the origin of the IRFR affine structure. The drift is not chosen
solely for tractability. It is forced by the requirement that the
rolling field remain compatible with equilibrium moment transport. The
diffusion is then tied to the same maturity loading. Volatility is
therefore not introduced as an independent curve-fitting device; it
inherits its shape from the admissible field structure.

The short end, $\tau=0$, has a special role. In rolling-maturity coordinates
it is a true boundary. Every point of the curve is transported toward
it. It is also the point at which monetary policy most directly enters
the interest-rate system. This makes the boundary mathematically natural
and economically meaningful.

The IRFR primitive therefore turns the term-structure problem into a
boundary-value problem. The interior law governs admissible
rolling-field evolution. The boundary law closes the system.

\section{HJM: representation versus explanation}\label{hjm-representation-versus-explanation}

The relationship with the Heath--Jarrow--Morton (HJM) framework must be
stated carefully.

HJM is a broad and powerful arbitrage-free representation of
forward-rate dynamics. Given a specification of the forward-rate
volatility field, HJM identifies the drift restriction required to
preclude arbitrage. In this sense, HJM describes what arbitrage-free
forward-rate dynamics may be.

The IRFR framework begins from a different question. It asks whether a
particular admissible field structure can be generated from a small set
of structural requirements. The point is not to replace HJM as an
ambient no-arbitrage representation. An admissible IRFR model should be
capable, after translation into standard forward-rate coordinates, of
being represented within a suitable HJM-type formulation.

The distinction is therefore not one of arbitrage-free versus
non-arbitrage-free modelling. It is one of representation versus
explanation.

HJM describes what arbitrage-free forward-rate dynamics may be. IRFR
seeks to explain why a particular admissible field structure should
arise.

In a standard HJM formulation, the volatility structure is typically
specified as an input. The arbitrage-free drift then follows. In the
IRFR formulation, the maturity structure of volatility is constrained by
the rolling-field geometry. The field's drift, diffusion, maturity
loading, and boundary behaviour are not independent modelling choices.
They are linked through the admissibility conditions of the IRFR system.

This is why the HJM embedding is not necessarily structure-preserving.
The same dynamics may be representable in HJM coordinates, but the
reason for their form may become obscure. In the IRFR variables, the
short end is a boundary, the maturity coordinate is the
field\textquotesingle s maturity-domain coordinate, and the admissible
modes arise from the rolling-maturity structure. Translated into a
conventional HJM representation, these features may appear as a
complicated or non-Markovian volatility specification.

Thus HJM should be understood as the ambient arbitrage-free
representation, while IRFR is proposed as a boundary-adapted structural
explanation inside that broader universe.

This distinction will become particularly important in the repeated-root
and boundary-law chapters. A humped volatility loading that appears
exogenous in HJM may arise as an admissible boundary mode in IRFR
coordinates. The object has not changed. The geometry has.

\section{Boundary laws and monetary policy}\label{boundary-laws-and-monetary-policy}

The short end of the yield curve is economically special. Central banks
operate most directly through short-maturity rates: overnight rates,
reserve remuneration, standing facilities, policy corridors, and
short-end forward guidance. They do not directly set the entire forward
curve. Their influence propagates from the short end into the rest of
the maturity domain.

In the IRFR formulation, this observation has a direct mathematical
expression. Monetary policy is represented as a boundary law at $\tau=0$.

A boundary law is not merely an endpoint condition. It is the rule that
determines how the field is closed at the short end. Different closures
correspond to different economic regimes.

A Neumann/free boundary represents the non-pinned benchmark: the short
end remains stochastic and the ordinary exponential short-end loading is
retained. A Dirichlet boundary represents an idealised fixed-policy
regime, in which short-end stochastic loading is suppressed by the
boundary. A Robin boundary represents partial anchoring, in which the
short end is influenced by policy but retains residual market
stochasticity.

These cases should not be read as literal mechanical descriptions of
central-bank operations. They are structural idealisations. Their
purpose is to identify how policy regimes select admissible maturity
modes.

The Dirichlet case is especially important as a limiting benchmark.
Economically, it corresponds to an exaggerated monetary-policy
intervention: the short end is fixed. Mathematically, it imposes a local
boundary condition on the IRFR field. In the repeated-root branch, this
boundary condition generates a linear-exponential humped loading of the
form associated with Ritchken--Chuang-type volatility specifications.

This gives the Dirichlet repeated-root case a special interpretive role.
In standard HJM finite-dimensional coordinates, a stationary humped
volatility specification of this type is usually regarded as
non-Markovian. In the IRFR formulation, the same humped structure is
generated by a Markovian boundary-value field. What appears as
non-Markovian memory in one representation may become a local boundary
law in the natural rolling-maturity geometry.

The claim is precise. The model is not asserted to be one-dimensional
Markov in the ordinary short-rate sense. It is Markovian only in the
field-state sense: an infinite-dimensional boundary-value evolution of
the IRFR field, not a finite-dimensional short-rate or low-dimensional
HJM state. The recovery of Markovian structure occurs because the state
geometry is changed.

The Robin case provides the more realistic policy analogue. It mixes
retained short-end stochastic loading with boundary-induced anchoring.
The short end is partially anchored rather than fixed. The resulting
loading contains both a residual boundary-active component and a
boundary-induced component. This gives a structural interpretation to
humped or shifted volatility profiles: they reflect the mixture of
admissible modes selected by the policy boundary.

In this sense, monetary policy is not modelled as a global curve-wide
force. It is modelled as a boundary law whose effects are transmitted
through the admissible geometry of the field.

\section{Comparison with existing literature}\label{comparison-with-existing-literature}

A useful comparison is with quadratic term-structure models. Quadratic
models enrich the standard affine tradition by allowing bond prices,
yields, or pricing equations to depend quadratically on latent state
variables. They are powerful because they can generate nonlinear yield
dynamics and richer distributional behaviour.

The IRFR framework takes a different route. It does not begin by adding
polynomial state dependence. It first changes the geometric primitive
from calendar maturity to rolling maturity and asks what drift and
loading structures are forced by the resulting boundary-field problem.

In this sense, IRFR is not a substitute for quadratic term-structure
models, but a minimalist structural alternative. It shows how
non-monotone maturity loadings can arise from admissibility and boundary
selection before additional polynomial state machinery is introduced.

\section{Claims and limitations}\label{claims-and-limitations}

The framework is intentionally ambitious, but its claims must be
separated.

First, the monograph establishes admissibility. It constructs a class of
IRFR dynamics satisfying the proposed structural axioms and boundary
closures. This is the mathematical foundation of the work.

Second, it argues for minimality. The claim is not that the framework is
the most general arbitrage-free model of interest rates. It is not. Nor
is the claim that no other model can reproduce the same empirical
shapes. Rather, the claim is that the IRFR construction derives several
objects that are often introduced separately: admissible maturity
loadings, volatility localisation, boundary anchoring, and
policy-induced mode selection. To the extent that these arise from one
mechanism, the model is structurally economical.

Third, it motivates a canonicality conjecture. The repeated appearance
of the same modes, boundary structures, and humped loadings suggests
that the IRFR geometry may be the natural coordinate system for a class
of interest-rate phenomena. This is a conjecture, not a theorem. The
present work does not prove uniqueness of the IRFR construction among
all possible term-structure theories.

Fourth, it proposes an empirical programme. The framework becomes
meaningful only if its structural parameters can be related to observed
yield curves, volatility profiles, and monetary-policy regimes. The
mapping to humped-volatility models provides one bridge, but empirical
support must come from calibration and testing rather than assertion.

Fifth, it distinguishes field-level admissibility from pricing
no-arbitrage. The affine-closure calculation imposes local coherence on
the maturity-domain field: relative stochastic motion across maturities
must be compatible with a common loading structure. This condition is
no-arbitrage-like in motivation, because it rules out arbitrary
inconsistent relative motion along the curve, but it is not itself a
self-financing no-arbitrage pricing theorem. The IRFR value $x(\tau)$
is a rate-field coordinate rather than a strictly positive tradable
price process. A unit-notional exposure to the IRFR has no natural
positive initial value and cannot itself serve as a numeraire. A full
risk-neutral pricing theory therefore requires an additional layer
mapping the IRFR field into tradable instruments, specifying a numeraire
and pricing measure, and deriving the corresponding martingale
restrictions.

The epistemic status of the work can therefore be summarised as follows:

admissibility is proved;

minimality is argued;

canonicality is conjectured;

empirical relevance is to be tested.

This discipline is essential. The purpose of the monograph is not to
declare a final theory of interest rates. It is to develop a coherent
field formulation and then ask whether the resulting structure is
powerful enough to deserve that interpretation.

\section{Roadmap}\label{roadmap}

The monograph proceeds in five parts.

Part I introduces interest rates as a maturity-domain field. It defines
the IRFR primitive, explains rolling-maturity transport, and shows why
the short end becomes a boundary. The goal is to establish the geometric
language of the theory before the technical machinery is introduced.

Part II develops the axiomatic foundations and transport mechanics. It
states the structural axioms, derives the affine drift through
KFE-consistent moment transport, and introduces the common maturity
loading. It also explains the one-factor covariance structure and the
role of signal and noise within a regime.

Part III studies the r-branch and the admissible mode structure. The
ratio between the temporal drift-diffusion scale and the maturity-decay
scale generates distinct-root and repeated-root regimes. At the critical
value r=1, the roots coalesce and the linear-exponential repeated-root
mode appears. This mode is central to the structural interpretation of
humped volatility.

Part IV develops the boundary laws. Axiom 4 is reinterpreted as the
location of economic regime selection. Neumann/free, Dirichlet, and
Robin boundary regimes are studied as the non-pinned benchmark, the
pinned-policy limit, and the partially anchored regime. The Dirichlet
repeated-root case gives the cleanest example of a boundary-selected
volatility hump. The Robin case provides the more realistic model of
partial policy anchoring.

Part V builds the empirical and modelling bridge. It relates the IRFR
boundary modes to Ritchken--Chuang-type humped volatility specifications,
compares the minimalist IRFR route with richer nonlinear frameworks such
as quadratic term-structure models, translates the results into
conventional observable rates, and discusses calibration, testing, and
empirical interpretation. It also identifies open problems, including
multi-factor IRFR fields, stochastic boundary laws, and no-arbitrage
pricing in the presence of a boundary-field short end.

The final principle of the monograph is therefore simple:

Interest rates are maturity-domain fields transported through calendar
time and constrained by boundary laws.

The remaining chapters develop this principle mathematically.

\chapter{The IRFR Primitive and the Field Representation}\label{chapter-2}

\section{Introduction}\label{introduction}

The previous chapter introduced the programme of the monograph: to
formulate interest-rate dynamics as a maturity-domain field theory in
which admissible dynamics arise from rolling-maturity invariance and
observable shapes are selected by boundary laws. This chapter develops
the primitive object on which that programme is built.

The central modelling choice is to work with the instantaneous rolling
forward rate, or IRFR. Instead of describing the forward curve by fixing
a calendar maturity $T$, the IRFR formulation describes the curve by
remaining maturity,

\[
\tau = T - t
\]

The resulting object is

\[
x(t,\tau) = F(t,t+\tau), \tau \ge 0,
\]

where $t$ is calendar time and $\tau$ is time to maturity. The forward curve is
therefore represented as a stochastic field over the half-line {[}0,$\infty$).

This chapter explains why this change of primitive matters. The move
from \((t,T)\) to \((t,\tau)\) does not merely relabel the standard forward-rate
curve. It exposes the transport structure of the term structure. As
calendar time advances, every fixed maturity rolls toward the short end.
The point $\tau=0$ becomes a true boundary of the maturity domain. The model
is therefore not just a stochastic process indexed by maturity, but a
stochastic field transported through calendar time and closed by a
boundary law.

The aim of this chapter is to make that geometry explicit before the
axioms and derivations are introduced. In particular, the chapter shows
how the coordinate transformation turns maturity into a transport
coordinate, separates deterministic roll-down from stochastic evolution,
and makes the short end the natural boundary of the field.

\section{From calendar maturity to rolling
maturity}\label{from-calendar-maturity-to-rolling-maturity}

Let \(F(t,T)\) denote the instantaneous forward rate observed at calendar
time \(t\) for maturity date \(T\). The conventional notation treats \(T\) as the
maturity label. For a fixed \(T\), the process \(F(t,T)\) describes the
evolution of the forward rate associated with that maturity date.

The IRFR formulation uses remaining maturity instead:

\[
\tau = T - t
\]

Equivalently,

\[
T = t + \tau
\]

The instantaneous rolling forward rate is then defined by

\[
x(t,\tau) = F(t,t+\tau)
\]

For each fixed calendar time t, the map

\[
\tau \mapsto x(t,\tau)
\]

is the forward curve expressed as a function of remaining maturity. Thus
\(x(t,\cdot)\) is a curve over the maturity domain, while \(t \mapsto x(t,\cdot)\) describes
the evolution of that curve through calendar time.

The distinction becomes important when a calendar maturity is held
fixed. If \(T\) is fixed, then

\[
\tau(t) = T - t,
\]

and therefore

\[
d\tau/dt = -1
\]

A forward contract with fixed maturity date does not remain at a fixed
point in the maturity domain. It moves deterministically toward the
short end as time passes.

For example, consider a contract that has five years remaining maturity
at time \(t=0\). At that date it sits at \(\tau=5\). One year later, if the
calendar maturity date is unchanged, the same contract has \(tau=4\). It has
moved left along the maturity axis even if no new market shock has
occurred.

This simple identity is the source of the rolling-field geometry.
Calendar time does two things simultaneously: it changes the stochastic
state of the curve and it transports each fixed-maturity point along the
maturity axis.

\section{The transport identity}\label{the-transport-identity}

The rolling-maturity representation makes the transport term explicit.
Since

\[
F(t,T) = x(t,T-t),
\]

differentiating along a fixed calendar maturity \(T\) gives

\[
\frac{d}{dt} F(t,T) = \partial_t  x(t,\tau) - \partial_\tau x(t,\tau), \tau = T - t
\]

Equivalently, the change in the conventional forward rate along a fixed
maturity has two components: the calendar-time change of the rolling
field and the deterministic roll-down along the maturity coordinate.

calendar-time evolution: $\partial_t x(t,\tau)$

maturity transport: $-\partial_\tau x(t,\tau)$

A simple example makes the point concrete. Suppose, momentarily, that
the curve is linear in remaining maturity:

\[
x(t,\tau) = a + b\tau
\]

Then $\partial_\tau x(t,\tau)=b$, so the deterministic roll-down contribution is

\[
-\partial_\tau x(t,\tau) = -b
\]

Thus, even without an explicit calendar-time movement in the field, a
fixed calendar maturity experiences deterministic change as it rolls
along the sloped curve. What appears as drift in fixed-maturity
coordinates may partly be mechanical movement along the curve.

This identity is the first formal indication that the IRFR
representation is not merely a relabelling. It separates stochastic
evolution in calendar time from deterministic transport in remaining
maturity. The transport correction that appears later in the affine
drift is a consequence of this coordinate geometry.

The identity also explains why the short end becomes structurally
important. Since fixed maturities move in the direction of decreasing $\tau$,
all rolling-maturity trajectories are transported toward the boundary
$\tau=0$. The boundary is therefore not appended to the model after the fact;
it is already present in the coordinate transformation.

\section{The forward curve as a maturity-domain
field}\label{the-forward-curve-as-a-maturity-domain-field}

In the IRFR formulation, the forward curve is treated as a field

\[
x : [0,\infty) \times [0,\infty) \to \mathbb{R},
\]

where \(t\) indexes calendar time and $\tau$ indexes remaining maturity. For each
t, $x(t,\tau)$ is a function of $\tau$. For each $\tau$, $x(t,\tau)$ evolves stochastically
through calendar time.

The term ``field'' is used in this precise sense. The model is not a
collection of unrelated scalar processes. It is a maturity-indexed
object whose values across $\tau$ are linked by the rolling structure of the
curve. The maturity coordinate is not simply a label attached to
different instruments. It is the maturity-domain coordinate along which
the curve is transported.

The field viewpoint separates two forms of variation.

First, there is stochastic variation in calendar time. This is the usual
source of risk in an interest-rate model: the curve changes because the
state of the market changes.

Second, there is deterministic variation along the maturity domain. Even
in the absence of a new market shock, a fixed calendar maturity rolls
toward shorter remaining maturity as time passes. This movement must be
reflected in the drift of the field.

A model that ignores this second component risks mixing two distinct
effects: genuine stochastic evolution and deterministic maturity
transport. The IRFR formulation keeps them separate.

This separation is one reason the framework is able to distinguish the
stochastic dimension of the model from the deterministic maturity
geometry of the curve. A model may be one-factor in randomness, driven
by a single Brownian source, while still possessing a non-trivial
maturity-domain structure through its admissible loading functions.

\section{Local rate exposure and the economic meaning of the IRFR
coordinate}\label{local-rate-exposure-and-the-economic-meaning-of-the-irfr-coordinate}

The IRFR primitive is not introduced solely because it gives a cleaner
mathematical representation. It also has a direct economic
interpretation. The local field increment identifies the rate shock at a
particular remaining maturity, and hence the local component of risk to
which a forward-rate exposure at that maturity is sensitive.

This qualification is important. A forward rate is not itself a full
value process. The actual marked-to-market P\&L of a traded instrument
will generally depend on discounting, annuity factors, contract
specification, numeraire choice, and the loading of the instrument on
the forward curve. The IRFR increment should therefore be read as the
local driver of such an exposure, not as a complete valuation formula.

In conventional notation, the forward rate F(t,T) is attached to a fixed
calendar maturity. In rolling-maturity notation, the same rate is
represented as $x(t,\tau)$, where $\tau$ is the remaining maturity. This makes it
possible to locate an infinitesimal forward-rate exposure directly at a
point of the maturity domain.

The local field increment

\[
dx(t,\tau)
\]

therefore represents the instantaneous local rate shock at remaining
maturity $\tau$. This is not merely a visual or notational symmetry. It is
the connection between the stochastic field and the economically
meaningful term-structure exposure. The maturity coordinate describes
where the exposure sits on the curve; the calendar-time increment
identifies the local movement to which that exposure is sensitive.

This observation is one reason the rolling-maturity representation is
structurally useful. The same object that describes the local economics
of a forward-rate exposure also supplies the coordinate system in which
the curve is transported toward the short-end boundary. The IRFR
primitive therefore aligns three elements:

local forward-rate exposure;

rolling maturity transport;

boundary-value geometry.

This alignment is what is meant by the symmetry of the IRFR
representation. The symmetry is not merely aesthetic. It means that the
economically local object and the geometrically local object coincide.
The field coordinate is both the location of a forward-rate exposure and
the maturity-domain coordinate along which the curve evolves.

This is precisely what is harder to see in a standard
Heath--Jarrow--Morton (HJM) representation. HJM provides the broad
arbitrage-free dynamics of the forward curve, but the local economic
exposure, the rolling transport geometry, and the short-end boundary
condition are not organised around the same coordinate system. In the
IRFR formulation they are. This is why the boundary-condition analysis
becomes cleaner: the boundary is not introduced as an additional
restriction on an already-specified volatility field, but appears as the
natural endpoint of the economically meaningful rolling field.

\section{Transport toward the
boundary}\label{transport-toward-the-boundary}

The maturity domain of the IRFR field is the half-line

\[
\tau \in [0,\infty).
\]

The point $\tau=0$ corresponds to the instantaneous short end. Every
fixed-maturity contract moves toward this point as calendar time
advances. Thus the short end is not merely a limiting case of the
forward curve. It is the boundary toward which the rolling field is
transported.

This observation gives the term-structure problem a boundary-value
character. The interior of the maturity domain describes the behaviour
of rates with positive remaining maturity. The short end closes the
system.

The direction of transport is important. Since

\[
d\tau/dt = -1,
\]

calendar time carries the field from larger $\tau$ toward smaller $\tau$. The
boundary at $\tau=0$ is therefore not an arbitrary endpoint. It is the
location at which rolling maturities arrive.

This has two consequences.

First, the drift of the IRFR field must include a transport correction.
If the equilibrium curve depends on $\tau$, then a point rolling down the
curve experiences a deterministic change even if its stochastic
displacement from equilibrium is unchanged. A drift specification that
does not account for this movement is not consistent with the
rolling-maturity geometry.

Second, the boundary at $\tau=0$ requires a closure condition. The field
cannot be fully specified by its interior law alone. The behaviour at
the short end determines which admissible modes are realised.

This is the mathematical origin of the boundary law.

\section{The economic meaning of the short-end
boundary}\label{the-economic-meaning-of-the-short-end-boundary}

The boundary at $\tau=0$ is not only mathematically natural. It is also
economically distinguished.

Central banks operate most directly at the short end of the term
structure. Policy rates, overnight rates, reserve remuneration, standing
facilities, and policy corridors all act first on the shortest
maturities. Longer maturities respond through expectations, risk premia,
market clearing, and the transmission of policy through the curve. But
policy does not usually act by directly fixing every maturity.

For the present chapter, the important point is not yet the
classification of possible policy regimes. That will be developed later.
The important point is simply that $\tau=0$ is the mathematically natural
location at which short-rate policy enters the maturity-domain field.

Thus the IRFR primitive gives monetary policy a natural location without
requiring policy to be modelled as a curve-wide perturbation. Policy
enters as a boundary law at $\tau=0$, and its effects propagate through the
admissible maturity geometry of the field.

\section{Relation to the conventional forward
curve}\label{relation-to-the-conventional-forward-curve}

The IRFR field is not a different market object from the conventional
instantaneous forward curve. It is a different representation of the
same underlying term structure.

Given F(t,T), define

\[
x(t,\tau) = F(t,t+\tau)
\]

Conversely, given $x(t,\tau)$, the conventional forward rate can be recovered
by

\[
F(t,T) = x(t,T-t)
\]

Thus the two representations are formally equivalent at the level of
pointwise curve values. The distinction lies in what the representation
makes structural.

The (t,T) representation is natural for no-arbitrage pricing across
calendar maturities. It treats each maturity date \(T\) as a label. The
(t,$\tau$) representation is natural for studying rolling maturity, field
transport, and boundary closure. It treats remaining maturity as the
domain coordinate.

This is why the IRFR formulation should not be understood as a rejection
of standard forward-rate modelling. It is a change of coordinates
designed to expose a structure that is otherwise difficult to see.

In particular, the boundary at $\tau=0$ is explicit in IRFR coordinates. The
rolling transport identity is explicit. The maturity-domain field is
explicit. These features are present implicitly in any forward-rate
model, but they are not always the organising principles of the theory.

\section{Why the primitive matters}\label{why-the-primitive-matters}

The choice of primitive determines which features of a model are treated
as inputs and which are treated as consequences.

If the volatility function is specified directly as a function of
maturity, then humped volatility can be introduced whenever it is
empirically useful. This may be effective for calibration, but it leaves
open the question of why that functional form should arise.

The IRFR programme asks a different question. Given the rolling-maturity
field and its boundary closure, which maturity loadings are admissible?
Which shapes are selected by boundary conditions? Which volatility
profiles arise from the same structure that determines drift and
transport?

In this sense, the primitive is not neutral. The IRFR primitive is
chosen because it makes the term-structure problem into a field problem
with a boundary. Once the field geometry is fixed, the model has less
freedom to introduce curve shapes independently.

This is the modelling cost discussed in Chapter 1. The framework gains
structural discipline by requiring that curve shape, volatility
localisation, maturity loading, and policy anchoring be traced back to
identifiable locations in the theory. Some restrictions enter through
the interior law. Others enter through the boundary law. The point is to
make these restrictions visible.

The IRFR primitive is therefore the first modelling commitment of the
monograph.

\section{Summary}\label{summary}

The rolling-maturity representation has four immediate consequences.

First, remaining maturity is a transport coordinate: d$\tau$/dt = -1.

Second, the forward curve is a field over the maturity domain: for each
t, $\tau$ $\mapsto$ $x(t,\tau)$ is the curve.

Third, the representation separates stochastic evolution from
deterministic roll-down. A local movement in fixed-maturity coordinates
contains both calendar-time change and transport along the maturity
axis.

Fourth, the short end $\tau=0$ is a boundary. It is mathematically required
for closure and economically distinguished as the natural location of
short-rate policy.

These four consequences explain why the IRFR primitive is not merely a
convenient notation. It is the coordinate system in which local
exposure, roll-down geometry, and boundary analysis are aligned.

\section{Transition to the axioms}\label{transition-to-the-axioms}

The field representation now motivates the axiomatic structure that
follows.

The coordinate transformation has done three things. It has turned
remaining maturity into a transport coordinate. It has turned the short
end into a boundary. And it has separated stochastic evolution in
calendar time from deterministic roll-down along the maturity axis.

The axioms of the next chapter formalise these consequences.

A stochastic model for $x(t,\tau)$ must specify how the field evolves in
calendar time. This is the role of the stochastic evolution axiom.

The model must respect the fact that the maturity coordinate is rolling.
This is the role of rolling-maturity invariance.

The model must preserve consistency between stochastic evolution and the
equilibrium curve as maturities are transported toward the boundary.
This is the role of equilibrium moment transport.

Finally, the model must be closed at $\tau=0$. This is the role of the
boundary closure axiom.

Thus the four axioms of the IRFR framework correspond directly to the
geometry introduced in this chapter:

stochastic evolution;

rolling-maturity invariance;

equilibrium moment transport;

boundary closure.

The next chapter states these axioms explicitly and explains how they
divide the theory into two parts: an interior law determining admissible
dynamics, and a boundary law selecting admissible modes.

\chapter{The IRFR Axioms}\label{chapter-3}

\section{Introduction}\label{introduction-1}

Chapter 2 established the geometric structure of the IRFR field. It
showed that the change from calendar maturity \(T\) to remaining maturity
$\tau$=T-t turns the forward curve into a maturity-domain field, exposes
deterministic rolling transport, and identifies the short end $\tau=0$ as a
boundary.

This chapter formalises that geometry through four axioms. Each axiom
corresponds to one feature of the IRFR representation:

\begin{longtable}[]{@{}
  >{\raggedright\arraybackslash}p{(\columnwidth - 2\tabcolsep) * \real{0.5000}}
  >{\raggedright\arraybackslash}p{(\columnwidth - 2\tabcolsep) * \real{0.5000}}@{}}
\toprule\noalign{}
\begin{minipage}[b]{\linewidth}\raggedright
\textbf{Geometry established in Chapter 2}
\end{minipage} & \begin{minipage}[b]{\linewidth}\raggedright
\textbf{Axiom introduced in Chapter 3}
\end{minipage} \\
\midrule\noalign{}
\endhead
\bottomrule\noalign{}
\endlastfoot
Stochastic field & Axiom 1: stochastic evolution \\
Rolling maturity & Axiom 2: rolling-maturity invariance \\
Transport along equilibrium & Axiom 3: equilibrium moment transport \\
Boundary at $\tau=0$ & Axiom 4: boundary closure \\
\end{longtable}

The axioms are not arbitrary modelling assumptions. They are the formal
counterparts of the geometry already established. Axiom 1 gives the
field local stochastic evolution. Axiom 2 prevents the model from
assigning unrelated dynamics to two contracts with the same remaining
maturity observed at different calendar dates. Axiom 3 requires the
expected dynamics to be consistent with roll-down along the equilibrium
curve. Axiom 4 closes the field at the short-end boundary.

Without Axiom 2, the coordinate $\tau$ would lose its meaning as a
maturity-domain coordinate. It would become only a relabelling of
calendar maturities. The purpose of the axioms is therefore to preserve
the interpretation of $x(t,\tau)$ as a rolling field rather than a collection
of unrelated forward-rate processes.

The first three axioms determine the interior dynamics of the field.
They restrict how the IRFR may evolve away from the boundary. The fourth
axiom determines how the system is closed at $\tau=0$, and therefore where
economic regime enters the model.

Axioms 1--3 determine admissible interior dynamics;

Axiom 4 determines the boundary law.

This division gives mathematical form to the programme stated in Chapter
1: dynamics are determined by invariance, while observable shapes are
selected by boundary laws.

\section{Why axioms?}\label{why-axioms}

The IRFR framework uses axioms to make the modelling cost visible. In
term-structure models, assumptions often enter through several different
choices: the state variable, the factor structure, the drift, the
volatility function, the long-run mean, the policy rule, and the
calibration procedure. These choices may be individually reasonable, but
their structural role is not always explicit.

The axioms assign each restriction to a specific location in the theory.
The stochastic character of the curve enters through Axiom 1. The
rolling geometry of remaining maturity enters through Axiom 2. The
compatibility between stochastic evolution and deterministic roll-down
enters through Axiom 3. The economic regime at the short end enters
through Axiom 4.

This is the sense in which the framework seeks internal consistency
rather than mere parsimony. The aim is not simply to use fewer
assumptions. It is to ensure that curve shape, volatility structure,
maturity loading, and boundary behaviour can be traced back to governing
principles.

The axioms should therefore be read as structural commitments. They
define the class of admissible IRFR fields considered in this monograph.
Later chapters derive their consequences.

\section{Axiom 1: stochastic
evolution}\label{axiom-1-stochastic-evolution}

The first axiom states that the IRFR field evolves stochastically
through calendar time.

For each remaining maturity $\tau>0$, the local evolution of the
field is assumed to admit a diffusion representation,

\[
dx(t,\tau) = -f(x,t,\tau)dt + \sqrt{2} g(x,t,\tau)dW_t,
\]

where $W_t$ is a Brownian motion, f is the IRFR drift-operator
coefficient, g is the IRFR diffusion-scale coefficient, and the actual
SDE drift is -f.

This is the local stochastic evolution axiom.

The axiom does not yet specify the functional form of f or g. It only
states that, locally in calendar time, the IRFR field evolves as a
diffusion. The later axioms restrict the admissible relationship between
the drift operator, the diffusion scale, and the maturity coordinate.

The use of a single Brownian driver reflects the one-factor version of
the theory developed in the first part of the monograph. This does not
mean that the maturity structure is one-dimensional. It means that the
source of randomness is one-dimensional. The way this randomness is
expressed across maturity is governed by the admissible maturity
loading.

In later chapters, the diffusion scale will be written using a
maturity-loading convention, for example through a function H. The
precise variance-rate convention is fixed when the Kolmogorov forward
equation is introduced. For now, the important point is simply that the
field has local stochastic dynamics.

\section{Interpretation of Axiom 1}\label{interpretation-of-axiom-1}

Axiom 1 gives the IRFR field its local risk structure.

At each remaining maturity $\tau$, the local field increment

\[
dx(t,\tau)
\]

represents the local rate shock to which an infinitesimal forward-rate
exposure at remaining maturity $\tau$ is sensitive. The full mark-to-market
P\&L of a traded instrument will generally depend on discounting,
annuity effects, contract-specific loading, and numeraire choice. The
IRFR increment is therefore not itself a complete valuation object. It
identifies the local field component driving that exposure.

This is the economic role of Axiom 1. It connects the stochastic field
to local forward-rate risk without requiring the forward rate itself to
be treated as a full value process.

The actual SDE drift, -f, describes the expected local movement of the
field. The diffusion-scale term g describes local uncertainty, with $\sqrt{}$2g
as the conventional volatility coefficient. The Brownian driver
represents the common stochastic shock. The maturity dependence of the
response to that shock is not yet specified; it will be constrained by
the later admissibility conditions.

This is deliberate. The IRFR framework does not begin by choosing a
volatility curve. It begins by specifying that the field evolves
stochastically, then uses rolling invariance, moment transport, and
boundary closure to determine which structures are admissible.

\section{Axiom 2: rolling-maturity
invariance}\label{axiom-2-rolling-maturity-invariance}

The second axiom states that the structural form of the IRFR dynamics is
invariant under the rolling-maturity representation.

The remaining maturity coordinate is

\[
\tau = T - t
\]

For a fixed calendar maturity T,

\[
d\tau/dt = -1
\]

Thus, as calendar time advances, a fixed-maturity point moves along the
maturity axis. A model of the IRFR field must be consistent with this
movement.

Axiom 2 imposes that, when expressed in rolling-maturity coordinates,
the structural form of the local dynamics depends on remaining maturity
$\tau$, rather than on the calendar maturity date \(T\) itself. Equivalently,
within a given regime, the same remaining maturity is governed by the
same structural law across calendar time.

For example, without Axiom 2, the model could assign different
structural dynamics to two contracts that both have $\tau$=5, merely because
they are observed at different calendar dates. This would break the
interpretation of $\tau$ as a maturity-domain coordinate rather than a
calendar label.

A one-year remaining maturity should be governed by the one-year
structural law, regardless of the calendar date on which it is observed,
except insofar as the state or regime of the field has changed. This is
the invariance principle underlying the IRFR construction.

\section{Interpretation of Axiom 2}\label{interpretation-of-axiom-2}

Rolling-maturity invariance is the geometric heart of the model.

Without it, the forward curve could still be written as $x(t,\tau)$, but the
notation would impose no discipline. The model could assign arbitrary
dynamics to each point of the maturity domain. Axiom 2 prevents this by
giving the maturity coordinate intrinsic meaning.

The point $\tau$=5 represents five years of remaining maturity, not a
particular calendar maturity date. If a contract has five years
remaining today, and another contract has five years remaining at a
later calendar date, the same structural law applies to both, within the
same regime.

This is what allows the model to speak meaningfully about maturity
loading. If the response of the curve to a stochastic shock depends on
$\tau$, then that dependence is not merely a label attached to a particular
maturity date. It is part of the field's deterministic maturity
geometry.

Axiom 2 is therefore the bridge between the economic object and the
field representation. It is also the first step toward boundary
analysis. Once $\tau$ is the maturity-domain coordinate, $\tau=0$ is the boundary
of that domain.

\section{Axiom 3: equilibrium moment
transport}\label{axiom-3-equilibrium-moment-transport}

The third axiom imposes consistency between stochastic evolution and
deterministic roll-down.

Let

\[
m(\tau)
\]

denote an equilibrium mean curve for the IRFR field. Since m is a
function of remaining maturity, a fixed calendar maturity rolling
through time also moves along this equilibrium curve.

For a fixed T, remaining maturity is

\[
\tau = T - t
\]

The equilibrium value associated with that fixed maturity is therefore

\[
m(T-t)
\]

Differentiating with respect to calendar time gives

\[
\frac{d}{dt} m(T-t) = -m'(\tau)
\]

The sign is important. As calendar time advances, the maturity
coordinate decreases. A fixed-maturity point rolls toward the short end.

Axiom 3 requires the first moment of the stochastic dynamics to be
compatible with this deterministic roll-down. Informally,

\[
expected stochastic evolution must be consistent with transport along m(\tau)
\]

This condition is what later forces the drift structure to contain two
components: a displacement term measuring deviation from the equilibrium
curve and a transport correction reflecting movement along the maturity
axis.

The detailed derivation is deferred to the next chapter. For now, the
point is structural. If the field has an equilibrium curve, and
maturities roll toward the boundary, then the expected motion of the
field must respect the fact that the equilibrium curve is itself being
traversed.

\section{Running example: a sloped equilibrium
curve}\label{running-example-a-sloped-equilibrium-curve}

A simple running example makes the role of transport visible.

Suppose the equilibrium curve is locally linear,

\[
m(\tau) = a + b\tau
\]

For a fixed calendar maturity T, the equilibrium value along the rolling
maturity is

\[
m(T-t) = a + b(T-t)
\]

Differentiating gives

\[
\frac{d}{dt} m(T-t) = -b
\]

Thus, even if the field is exactly on the equilibrium curve and no new
stochastic shock occurs, the value associated with a fixed maturity
changes mechanically as the contract rolls along the curve.

If b\textgreater0, the equilibrium curve is upward sloping in $\tau$, and the
fixed-maturity point rolls downward as time passes. If b\textless0, the
roll-down effect has the opposite sign.

The flat case b=0 is useful as a limiting benchmark: there is then no
roll-down correction. But precisely because the flat case hides the
transport effect, the sloped case reveals why the correction is needed.

A model that omits this deterministic roll-down will misidentify part of
the curve's mechanical movement as stochastic drift. Axiom 3 prevents
this by requiring the expected dynamics to respect the transport of the
equilibrium curve.

This example is deliberately simple, but it captures the essential
point. In a rolling-maturity field, the drift must be compatible not
only with stochastic mean reversion, but also with motion along the
maturity axis.

\section{Interpretation of Axiom 3}\label{interpretation-of-axiom-3}

Axiom 3 is where the coordinate geometry begins to constrain the
dynamics.

A standard local diffusion can specify drift and volatility at each
state. But an IRFR field is not merely a collection of local diffusions.
A point associated with a fixed calendar maturity is moving through the
maturity domain. Its expected evolution must therefore include both
stochastic displacement and deterministic roll-down.

Suppose the field is exactly on the equilibrium curve. As calendar time
passes, a fixed-maturity point should remain consistent with the
equilibrium curve as its remaining maturity declines. If the drift does
not account for this movement, the model creates an inconsistency
between the probabilistic evolution of the rate and the deterministic
geometry of remaining maturity.

This implies that the affine drift-operator structure derived later must
contain a transport correction. The correction is not an arbitrary
adjustment. It is required by the rolling-maturity geometry.

Axiom 3 therefore connects the probabilistic dynamics of the field with
its deterministic maturity geometry.

\section{Axiom 4: boundary closure}\label{axiom-4-boundary-closure}

The fourth axiom states that the IRFR field must be closed at the
short-end boundary,

\[
\tau = 0
\]

The maturity domain is the half-line {[}0,$\infty$). The point $\tau=0$ is the
instantaneous short end. Since fixed-maturity points are transported
toward this boundary, the interior dynamics alone do not fully determine
the field. A boundary condition is required.

Axiom 4 is the boundary closure axiom.

At this stage, the axiom should be understood abstractly. It does not
prescribe a single universal boundary condition. It states that an
admissible IRFR model must specify how the field behaves at the short
end.

The distinction is crucial:

the need for boundary closure is structural;

the form of boundary closure is regime-dependent.

In the original non-policy benchmark, the boundary may be interpreted as
a natural or positivity-preserving closure. In the monetary-policy
extension, different policy regimes correspond to different boundary
laws. These cases will be studied later. The present chapter only
identifies the boundary as the location at which closure enters.

Thus Axiom 4 is the point at which economic regime enters the model.

\section{Interpretation of Axiom
4}\label{interpretation-of-axiom-4}

Axiom 4 completes the field problem.

Axioms 1-\/-3 determine the admissible interior dynamics. They specify
that the field evolves stochastically, respects rolling-maturity
invariance, and transports moments consistently along the equilibrium
curve. But they do not by themselves determine how the field is closed
at the short end.

The boundary condition performs this selection. It is the closure rule
for the maturity-domain field, not merely another local modelling
assumption. Different closures can select different admissible modes,
and therefore different observable curve and volatility shapes.

Economically, this is natural. The short end is where monetary policy
acts most directly. Mathematically, it is the boundary of the field. The
boundary law therefore provides the link between policy regime and
term-structure shape.

The later boundary chapters will study Neumann/free closure, Dirichlet
fixed-boundary closure, and Robin partial-anchoring closure. The present
chapter does not analyse these cases. It only establishes where they
enter the theory.

\section{Interior law versus boundary
law}\label{interior-law-versus-boundary-law}

The four axioms separate the model into two components.

The first component is the interior law, governed by Axioms 1-\/-3. It
determines how the field evolves away from the boundary, how the
maturity coordinate constrains the dynamics, and how the drift must
remain compatible with equilibrium transport.

The second component is the boundary law, governed by Axiom 4. It
determines how the field is closed at the short end and which admissible
behaviour is selected.

This separation allows the model to distinguish between the set of
dynamically admissible structures and the economic regime that selects
from them.

Axioms 1-\/-3: admissible interior dynamics

Axiom 4: boundary selection

This is the technical version of the monograph's organising principle:
invariance determines the admissible dynamics, while boundary laws
select observable shapes.

\section{Summary and transition}\label{summary-and-transition}

The IRFR framework is built from four structural axioms. Axiom 1 gives
the field local stochastic evolution. Axiom 2 imposes rolling-maturity
invariance. Axiom 3 requires equilibrium moment transport. Axiom 4
closes the field at the short-end boundary.

Together they organise the model as a boundary-value problem for a
stochastic maturity-domain field. The first three axioms determine the
interior law: they restrict the admissible drift, diffusion, and
maturity loading. The fourth axiom determines the boundary law: it
selects which admissible behaviour is realised at the short end.

The next chapter turns the axioms into equations. Axiom 1 provides the
local diffusion form. Axiom 2 imposes rolling-maturity invariance. Axiom
3 requires consistency of first moments under deterministic roll-down.
Together, these conditions force the affine drift-operator structure and
the common maturity loading that underlie the IRFR model.

The Kolmogorov forward equation provides the mechanism for this
derivation. It links the local drift and diffusion of the field to the
evolution of the probability density and its moments. When combined with
rolling-maturity transport, it yields the admissibility conditions that
govern the interior law.

The next chapter therefore studies the Kolmogorov forward equation,
affine closure, and moment transport. It is there that the axioms stated
here become calculable restrictions on the form of the IRFR dynamics.

\chapter{The Kolmogorov Forward Equation and Affine Closure}\label{chapter-4}

\section{Introduction}\label{introduction-2}

Chapter 3 stated the four axioms of the IRFR framework. The first three
axioms determine the interior law of the field: local stochastic
evolution, rolling-maturity invariance, and equilibrium moment
transport. The fourth axiom closes the field at the short-end boundary.

This chapter turns the interior axioms into equations. Its purpose is to
show how the local stochastic specification of the IRFR field becomes a
calculable restriction on drift, diffusion, and maturity loading.

The required tool is the Kolmogorov forward equation. The KFE links the
local drift and diffusion of a stochastic process to the evolution of
its probability density and moments. In the present setting, it is the
mechanism through which the axioms become admissibility conditions.

The main result is that the IRFR drift operator cannot be chosen freely.
Once the field is required to evolve as a diffusion, respect rolling
maturity, and transport its equilibrium moments consistently, the drift
structure must combine two components: a displacement term relative to
the equilibrium curve and a transport correction generated by
deterministic roll-down along the maturity axis.

The IRFR field is

\[
x(t,\tau)=F(t,t+\tau), \tau\ge0
\]

For each remaining maturity $\tau$, the local stochastic evolution is written
as

\[
dx(t,\tau) = -f(x,t,\tau)dt + \sqrt{2} g(x,t,\tau)dW_t,
\]

where f is the IRFR drift operator coefficient and g is the IRFR
diffusion-scale coefficient. The actual SDE drift is -f. In conventional
SDE notation, the volatility coefficient would be $\sigma$=$\sqrt{}$2g. This chapter
focuses on the interior law. Boundary closure is not yet imposed.
Boundary conditions return later, once the admissible interior dynamics
have been derived.

\section{From axioms to a density
equation}\label{from-axioms-to-a-density-equation}

Axiom 1 states that, at each remaining maturity $\tau$, the IRFR field
evolves locally as a diffusion,

\[
dx(t,\tau) = -f(x,t,\tau)dt + \sqrt{2} g(x,t,\tau)dW_t
\]

For a fixed $\tau$, let

\[
p(x,t;\tau)
\]

denote the probability density of the IRFR value $x(t,\tau)$. The Kolmogorov
forward equation describes how probability mass flows over the state
variable under drift and diffusion.

With the convention

\[
dx = -f(x,t,\tau)dt + \sqrt{2} g(x,t,\tau)dW_t,
\]

the density satisfies

\[
\frac{\partial p}{\partial t} = \partial/\partial x {[} f(x,t,\tau)p {]} + \partial^2/\partial x^2 {[} g^2(x,t,\tau)p {]}
\]

It is useful to write the instantaneous variance rate as

\[
v(x,t,\tau)=2g^2(x,t,\tau)
\]

Then the KFE becomes

\[
\frac{\partial p}{\partial t} = \partial/\partial x {[} fp {]} + 1/2 \partial^2/\partial x^2 {[} vp {]}
\]

This equation is local in the state variable x, but its coefficients may
depend on the maturity coordinate $\tau$. The role of Axioms 2 and 3 is to
restrict that $\tau$-dependence. The KFE is therefore the bridge between
stochastic evolution and moment transport.

\section{Moment equations}\label{moment-equations}

The first moment of the IRFR density is

\[
M_1(t,\tau) = E{[}x(t,\tau){]} = \int x p(x,t;\tau) dx
\]

Differentiating under the integral sign and using the KFE gives

\[
\frac{\partial M_1}{\partial t} = \int x (\frac{\partial p}{\partial t}) dx
\]

Substituting the KFE and integrating by parts, assuming boundary terms
vanish under the admissibility conditions, yields

\[
\frac{\partial M_1}{\partial t} = -E{[}f(x,t,\tau){]}
\]

The integration by parts step moves derivatives from the density onto
the test function x. Since the second derivative of x is zero, the
diffusion contribution drops out of the first-moment equation, leaving
the actual drift contribution, -E{[}f{]}, in expectation.

Thus the expected drift of the SDE is -E{[}f{]}. This is the first point
at which the axioms impose structure: Axiom 3 requires this moment
evolution to be consistent with deterministic roll-down along the
equilibrium curve.

Let the equilibrium mean curve be

\[
m(\tau)
\]

If the field is in equilibrium, then

\[
M_1(t,\tau)=m(\tau)
\]

For a fixed calendar maturity T, the remaining maturity is $\tau$=T-t. The
equilibrium value associated with that fixed maturity is therefore

\[
m(T-t)
\]

As shown in Chapter 3,

\[
\frac{d}{dt} m(T-t) = -m'(\tau)
\]

Equivalently, the calendar-time evolution of the rolling field must
include a transport correction associated with the derivative of m with
respect to $\tau$. This is the moment-transport requirement.

\section{The transport correction}\label{the-transport-correction}

To see why a transport correction is unavoidable, return to the running
example

\[
m(\tau)=\alpha+\beta\tau
\]

For a fixed calendar maturity T,

\[
m(T-t)=\alpha+\beta(T-t),
\]

and hence

\[
\frac{d}{dt} m(T-t) = -\beta
\]

Since

$m'(\tau)$=$\beta$,

the roll-down contribution is

-$m'(\tau)$=-$\beta$.

Thus, even when the field lies exactly on the equilibrium curve, the
fixed-maturity value changes mechanically as remaining maturity
declines.

If the drift were specified only as a displacement from equilibrium, it
would vanish on the equilibrium curve. That would be inconsistent with
roll-down whenever $m'(\tau)$ is non-zero. The drift must therefore include a
term that accounts for movement along the equilibrium curve.

The precise sign of this term depends on whether one refers to the IRFR
operator coefficient f or to the actual SDE drift -f. The fixed-maturity
roll-down contribution is -$m'(\tau)$, while the compensating term in the
fixed-$\tau$ rolling-field drift appears with the opposite sign. Since the
actual drift is -f, the corresponding term inside f appears with a minus
sign. In the convention used here, the admissible IRFR drift operator
takes the form

\[
f(x,t,\tau) = a^2{[}x-m(\tau){]} - m'(\tau)
\]

The term -$m'(\tau)$ inside f produces the +$m'(\tau)$ transport correction in the
actual SDE drift -f. It is not an arbitrary adjustment. It is required
by the coordinate geometry.

\section{Affine closure of the drift
operator}\label{affine-closure-of-the-drift-operator}

The next question is why the drift operator should be affine in x.

The first-moment equation is

\[
\frac{\partial M_1}{\partial t} = -E{[}f(x,t,\tau){]}
\]

Closure of the first moment requires E{[}f(x,t,$\tau$){]} to depend only on
$M_1(t,\tau)$, rather than on higher moments of the density. This holds for
all admissible densities when f is at most affine in x. A nonlinear
drift operator would generally cause the first-moment equation to depend
on higher moments, breaking first-moment closure.

Accordingly, consider the affine form

\[
f(x,t,\tau)=A(\tau)x+B(\tau),
\]

where $A(\tau)$ and $B(\tau)$ are functions of remaining maturity. Then

\[
E{[}f(x,t,\tau){]} = A(\tau)M_1(t,\tau)+B(\tau)
\]

At equilibrium, $M_1(t,\tau)$=$m(\tau)$. The moment-transport condition requires
the actual expected drift, -E{[}f{]}, to match the required transport of
the equilibrium curve. This imposes a relation between $A(\tau)$, $B(\tau)$, and
$m(\tau)$.

Choosing the displacement scale as $a^2$, the affine drift can be written
as

\[
f(x,t,\tau) = a^2{[}x-m(\tau){]} - m'(\tau)
\]

The sign convention used here is discussed explicitly in Section 4.7.

Equivalently,

\[
f(x,t,\tau) = a^2x - a^2m(\tau) - m'(\tau)
\]

Thus, in the affine representation,

$A(\tau)=a^2$,

and

$B(\tau)=-a^2m(\tau)-m'(\tau)$.

The affine drift operator is therefore not introduced merely for
tractability. It is the closure form compatible with first-moment
transport under the IRFR geometry.

\section{Interpretation of the affine
drift}\label{interpretation-of-the-affine-drift}

The IRFR drift operator

\[
f(x,t,\tau) = a^2{[}x-m(\tau){]} - m'(\tau)
\]

contains two distinct terms. When inserted into the local evolution
equation, the actual SDE drift is

\[
-f(x,t,\tau) = -a^2{[}x-m(\tau){]} + m'(\tau)
\]

This displayed form makes explicit that the actual SDE drift contains
the standard mean-reverting displacement term together with the fixed-$\tau$
transport correction.

The first term,

$a^2${[}x-$m(\tau)${]},

measures displacement from the equilibrium curve. Through the actual
drift -f, it produces the mean-reverting contribution $-a^2[x-m(\tau)]$.

The second term,

-$m'(\tau)$,

is the transport component inside f. Through the actual drift -f, it
produces the +$m'(\tau)$ transport correction required by fixed-$\tau$
rolling-field dynamics.

These two terms have different origins. The displacement term is
probabilistic: it describes the local tendency of the stochastic field
relative to equilibrium. The transport term is geometric: it is forced
by the rolling coordinate.

The running example makes the distinction visible. If

\[
m(\tau)=\alpha+\beta\tau,
\]

then

$m'(\tau)$=$\beta$.

The transport component is constant. It vanishes only when $\beta=0$, that is,
when the equilibrium curve is flat. In a sloped equilibrium curve, the
correction is unavoidable.

This is the first major payoff of the IRFR formulation. The drift
operator is not simply chosen. It is decomposed into a stochastic
displacement component and a deterministic transport component, each
with a clear structural origin.

\section{Sign conventions and mean
reversion}\label{sign-conventions-and-mean-reversion}

The affine drift operator written above uses the convention

\[
f(x,t,\tau) = a^2{[}x-m(\tau){]} - m'(\tau)
\]

The local stochastic evolution is written with -f as the actual drift,

\[
dx(t,\tau) = -f(x,t,\tau)dt + \sqrt{2}g(x,t,\tau)dW_t
\]

Thus f should not be read directly as the SDE drift coefficient. The
actual drift is -f. Under this convention, the displacement component
$a^2${[}x-$m(\tau)${]} inside f produces the standard mean-reverting
contribution $-a^2[x-m(\tau)]$ in the SDE.

We retain the IRFR convention because it is the convention used in the
admissibility algebra and in the later boundary-condition analysis. It
also allows the structural comparison between f and g to be written
directly. The essential point is the consistency of the sign convention
across the local evolution equation, the KFE, the moment-transport
condition, the equilibrium density, and the boundary laws.

For this reason, the affine form should be read as part of the complete
IRFR operator convention. Later chapters retain this convention
throughout. Whenever the formula is translated into a standard SDE
notation, the corresponding drift coefficient is -f.

\section{Variance and common maturity
loading}\label{variance-and-common-maturity-loading}

While the drift is determined by first-moment transport, the diffusion
structure enters through second-moment consistency.

Axiom 1 allows the local diffusion scale g to depend on the maturity
coordinate. Axiom 2 prevents this dependence from being arbitrary across
calendar maturities. Axiom 3 links the diffusion structure to the same
maturity geometry that governs the drift.

The one-factor specification means that all maturities share a common
Brownian driver,

$dW_t$.

Therefore the covariance between two maturity points is determined by
the product of their actual diffusion coefficients. If the
diffusion-scale coefficient has the form

\[
g(\tau)=aH(\tau),
\]

so that the conventional volatility coefficient is $\sqrt{}$2$g(\tau)$=$\sqrt{2}aH(\tau)$, then
the instantaneous covariance between maturities $\tau_1$ and $\tau_2$ is

\[
\operatorname{Cov}(dx(t,\tau_1), dx(t,\tau_2)) = 2a^2H(\tau_1)H(\tau_2)dt
\]

The normalisation with $\sqrt{}$2 in the SDE is chosen for later algebraic
convenience. The economically relevant object is the maturity loading
$H(\tau)$, which determines how the common stochastic shock is expressed
across the curve.

This covariance kernel is rank one. The stochastic dimension is one, but
the deterministic maturity loading $H(\tau)$ can be non-trivial.

This distinction will matter later. It is what allows a one-factor
stochastic model to generate structured maturity profiles without
introducing additional Brownian factors.

\section{Moment closure and
admissibility}\label{moment-closure-and-admissibility}

The KFE generates equations for all moments of the density. The first
moment gives the drift condition. The second moment introduces the
variance-rate structure.

Let

V(t,$\tau$)=Var{[}$x(t,\tau)${]}.

Under the diffusion specification, the evolution of the second moment
depends on both the actual drift -f and the variance rate

\[
v(x,t,\tau)=2g^2(x,t,\tau)
\]

For the IRFR model to remain admissible, the variance structure must
close consistently with the same maturity loading that appears in the
drift.

This is where coefficient matching enters. Coefficient matching requires
that the moment equations generated by the KFE hold for all admissible
states. This imposes algebraic restrictions on the drift operator and
variance coefficients. In practice, this means equating the $\tau$-dependent
coefficients generated by the moment equations with the proposed
functional form of the maturity loading.

Conceptually, coefficient matching asks whether the proposed drift
operator and diffusion scale can satisfy the KFE moment equations
simultaneously. If they can, the model is admissible. If not, the chosen
maturity loading is incompatible with the IRFR axioms.

Thus admissibility is not imposed by assertion. It is tested through the
KFE. The affine drift operator and the common maturity loading are the
first outputs of that test.

\section{Relation to HJM}\label{relation-to-hjm}

From the HJM perspective, the preceding construction can be read as
identifying a restricted admissible subclass of forward-rate dynamics.

The distinction is that the diffusion or volatility field has not been
freely specified and then paired with the corresponding arbitrage-free
drift. Instead, the admissibility conditions restrict the drift operator
and maturity loading together. The KFE and moment-transport equations
act as a consistency filter.

This is the technical version of the representation-versus-explanation
distinction introduced in Chapter 1. An IRFR model may be representable
in HJM coordinates, but the reason for its form is clearest in
rolling-maturity field coordinates.

The HJM representation describes what arbitrage-free forward-rate
dynamics may be. The IRFR derivation explains why a particular
boundary-adapted field structure is admissible.

\section{Summary}\label{summary-1}

This chapter has turned the first three IRFR axioms into the beginning
of a calculable interior law.

Axiom 1 gives the local diffusion. Axiom 2 imposes rolling-maturity
invariance. Axiom 3 requires equilibrium moment transport. Together,
these conditions force the drift operator to take an affine
transport-compatible form,

\[
f(x,t,\tau) = a^2{[}x-m(\tau){]} - m'(\tau)
\]

The first term measures displacement from equilibrium and generates mean
reversion through the actual drift -f. The second term is the transport
component inside f, generating the required +$m'(\tau)$ correction in the
actual drift.

The chapter has also introduced the role of the common maturity loading
in the diffusion structure. In the one-factor case, all maturities share
the same Brownian shock, but the deterministic loading determines how
that shock is expressed across the curve.

The main message is that the interior law is not freely chosen. It is
constrained by the interaction between stochastic evolution, rolling
maturity, and moment transport.

\section{Transition to admissible
modes}\label{transition-to-admissible-modes}

The next step is to determine which maturity loadings are admissible.

The affine drift operator identifies the structure of local mean
evolution, but the full IRFR model also requires the diffusion and
variance structure to satisfy the KFE consistently. This leads to
coefficient matching and to the emergence of admissible maturity modes.

The next chapter develops this admissibility calculation. It shows how
the KFE restricts the maturity loading, how the r-branch arises, and why
the repeated-root case later plays a central role in the boundary-law
interpretation of humped volatility.

\chapter{Admissible Maturity Modes and the r-Branch}\label{chapter-5}

\section{Introduction}

Chapter 4 derived the affine drift operator of the IRFR field. Under the
convention

\[
dx(t,\tau) = -f(x,t,\tau)\,dt + \sqrt{2} g(x,t,\tau)\,dW_t,
\]

the Kolmogorov forward equation and first-moment transport imply the
drift-operator form

\[
f(x,t,\tau) = a^2{[}x - m(\tau){]} - m'(\tau)
\]

The actual SDE drift is therefore

\[
-f(x,t,\tau) = -a^2{[}x - m(\tau){]} + m'(\tau),
\]

which contains standard mean reversion toward the equilibrium curve
together with the fixed-$\tau$ transport correction.

Chapter 4 also introduced the diffusion-scale coefficient g, with
instantaneous variance rate

\[
v(x,t,\tau) = 2g^2(x,t,\tau)
\]

The task of the present chapter is to determine which maturity loadings
are admissible. The drift structure alone does not complete the model.
The diffusion structure must also satisfy the Kolmogorov forward
equation consistently with the same maturity-domain geometry. This
requirement leads to coefficient matching and to a family of admissible
maturity modes.

The key result is that the admissible loading is not arbitrary. It is
governed by a second-order maturity equation whose solutions divide into
branches. These branches are indexed by a dimensionless parameter r,
which measures the relationship between the temporal drift-diffusion
scale and the maturity-decay scale.

The chapter develops three ideas:

\begin{center}\emph{coefficient matching restricts the maturity loading;}\end{center}
\begin{center}\emph{the resulting maturity equation generates an \(r\)-branch;}\end{center}
\emph{the repeated-root case r = 1 produces the linear-exponential
mode.}

The repeated-root case will later become central to the boundary-law
interpretation of humped volatility.

\section{From affine closure to admissibility}

The affine drift operator derived in Chapter 4 is

\[
f(x,t,\tau) = a^2{[}x - m(\tau){]} - m'(\tau)
\]

This form is required by first-moment closure and equilibrium transport.
It ensures that the expected dynamics of the field are compatible with
deterministic roll-down along the equilibrium curve.

But admissibility requires more than first-moment consistency. The full
probability density must evolve consistently under the KFE. In
particular, the second moment and variance rate must close in a way
compatible with the same maturity loading.

The diffusion-scale coefficient g determines the local variance rate,

\[
v(x,t,\tau) = 2g^2(x,t,\tau)
\]

In the one-factor model, all maturities share the same Brownian shock,
so the covariance between maturity points is generated by the product of
their maturity loadings. If

\[
g(\tau) = aH(\tau),
\]

then the covariance kernel is

\[
\operatorname{Cov}(dx(t,\tau_1), dx(t,\tau_2)) = 2a^2H(\tau_1)H(\tau_2)\,dt
\]

Thus the stochastic dimension is one, but the deterministic maturity
loading $H(\tau)$ may have non-trivial shape.

The question is therefore:

\[
which functions H(\tau) are compatible with the IRFR axioms?
\]
This is the admissibility problem.

\section{Coefficient matching}\label{coefficient-matching}

Coefficient matching is the algebraic test for admissibility.

The KFE generates moment equations. The affine drift operator determines
how the first moment evolves. The diffusion-scale coefficient determines
how the second moment and variance rate evolve. For a proposed maturity
loading to be admissible, these equations must be mutually consistent
for all admissible states in the model class.

In practice, this means that the $\tau$-dependent coefficients generated by
the moment equations must match the proposed functional form of the
maturity loading. If the terms generated by the KFE cannot be expressed
in the same loading basis, the proposed structure is not closed under
the dynamics.

This is the point at which the IRFR framework becomes restrictive. One
cannot choose an arbitrary maturity loading and then regard it as
admissible. The loading must survive the coefficient-matching test.

The result is a maturity-domain equation for the admissible loading
$H(\tau)$. In canonical form, this equation may be written as

\[
H''(\tau) + \lambda H'(\tau) + \mu H(\tau) = 0,
\]

where $\lambda$ and $\mu$ are constants determined by the drift-diffusion scale and
the maturity-domain normalisation.

The exact normalisation will be fixed below through the r-parameter. At
this stage, the important structural point is that admissible maturity
loadings are not arbitrary functions. They are modes of a second-order
linear equation on the maturity domain.

\section{The characteristic
equation}\label{the-characteristic-equation}

The admissibility equation admits exponential trial solutions of the
form

\[
H(\tau) = e^{-k\tau}
\]

Substituting this into

\[
H''(\tau) + \lambda H'(\tau) + \mu H(\tau) = 0
\]

gives

\[
k^2 - \lambda k + \mu = 0
\]

Thus the admissible decay rates are the roots of the characteristic
equation.

In the generic case, the two roots are distinct:

\[
k_1 \neq k_2
\]

The admissible loading is then

\[
H(\tau) = C_1e^{-k_1\tau} + C_2e^{-k_2\tau}
\]

For admissible decaying modes on the half-line, the relevant roots have
positive real part. In the real distinct-root case, this means

k$_1$ \textgreater{} 0, k$_2$ \textgreater{} 0.

The constants $C_1$ and C$_2$ are not determined by the interior equation
alone. They are selected by normalisation and boundary closure. This is
the first point at which the distinction between interior law and
boundary law becomes concrete:

\begin{center}\emph{the interior law determines the admissible modes;}\end{center}
\begin{center}\emph{the boundary law selects their realised combination.}\end{center}
\section{The r-parameter}\label{the-r-parameter}

It is useful to parameterise the characteristic equation by a
dimensionless quantity r. The role of r is to measure the relative
strength of the temporal drift-diffusion scale and the maturity-domain
decay scale.

Equivalently, r controls the discriminant of the characteristic
equation. In the normalisation used in the remainder of the monograph, r
is chosen so that the discriminant is proportional to r-1. Thus

\[
r \neq 1
\]

corresponds to distinct roots, while

\[
r = 1
\]
is precisely the repeated-root case.

This is the key mathematical role of r. It is not introduced as an
additional modelling factor. It is a dimensionless way of describing
which branch of the admissible maturity equation is realised.

The three relevant cases are therefore:

\emph{r \textgreater{} 1: distinct real roots under the chosen
normalisation;}

\begin{center}\emph{\(r=1\): repeated root.}\end{center}
\emph{r \textless{} 1: a non-generic branch whose admissibility depends
on the chosen real/complex root convention and boundary domain.}

The detailed interpretation of the r \textless{} 1 branch will therefore
depend on the exact normalisation and boundary regime. The repeated-root
case r = 1, however, has an invariant significance: it is the point at
which the two admissible decay modes coalesce.

\section{Distinct-root branch}\label{distinct-root-branch}

When the characteristic roots are distinct, the admissible loading is

\[
H(\tau) = C_1e^{-k_1\tau} + C_2e^{-k_2\tau}, k_1 \neq k_2
\]

This branch represents a mixture of two maturity-decay modes. The
stochastic dimension of the model remains one, because all maturities
are driven by the same Brownian shock. But the deterministic maturity
geometry may contain more than one decay scale.

This distinction is essential. A model can be one-factor in randomness
while still having a structured maturity-domain response. The covariance
kernel remains rank one, but the loading through which the common shock
is expressed across the curve can have non-trivial shape.

The distinct-root branch is therefore the generic admissible case.
Boundary conditions and normalisation determine which linear combination
of the two modes is realised.

\section{Repeated-root branch}\label{repeated-root-branch}

The repeated-root case occurs when

\[
k_1 = k_2 = k
\]

At this point, the two-exponential representation degenerates. The
second independent solution is no longer a distinct exponential.
Instead, the admissible loading takes the form

\[
H(\tau) = (C_1 + C_2\tau)e^{-k\tau}
\]

For admissible decaying modes on the half-line, we take

\[
k >{} 0
\]

The second term,

\[
\tau e^{-k\tau},
\]

is the important one. It is not a simple exponential decay. It vanishes
at the short end, rises to an interior maximum, and then decays to zero
as ($\tau\to\infty$).

Indeed, for

\[
h(\tau) = \tau e^{-k\tau},
\]

we have

\[
h'(\tau) = e^{-k\tau}(1 - k\tau)
\]

The maximum occurs at

\[
\tau = 1/k
\]

Thus the repeated-root branch naturally generates a humped maturity
loading.

This is the first appearance of the humped structure as an admissibility
result. It has not been imposed as a volatility shape. It has appeared
because the two admissible exponential modes have coalesced.

\section{The r = 1 branch}\label{the-r-1-branch}

Under the normalisation used in this monograph, the repeated-root case
is the r = 1 branch.

This branch has special importance because it generates the
linear-exponential loading

\[
H(\tau) \propto \tau e^{-k\tau}
\]

Under the parameter convention used later, the repeated-root decay rate
will be written as

\[
k = a^2,
\]

so that the humped component becomes

\[
H(\tau) \propto \tau e^{-a^2\tau}
\]

At this stage, this is an interior admissibility statement, not yet a
boundary-condition statement. The r = 1 branch identifies a humped mode
that is allowed by the interior law. Whether that mode is realised
depends on boundary closure.

This distinction is central:

\begin{center}\emph{interior admissibility produces the repeated-root mode;}\end{center}
\begin{center}\emph{boundary closure selects its economic realisation.}\end{center}
The r = 1 branch is therefore the bridge between the
coefficient-matching algebra and the later boundary-law interpretation
of humped volatility.

\section{Relation to humped volatility}

Humped volatility profiles are familiar in interest-rate modelling. In
many models, such a hump is specified directly because it fits observed
term-structure volatility.

The IRFR interpretation is different. The repeated-root mode shows that
a humped maturity profile can arise from the admissibility structure of
the field itself.

The key object is

\[
h(\tau) = \tau e^{-k\tau}
\]

This profile has three features:

\[
h(0) = 0;
\]

\[
h(\tau) has an interior maximum;
\]

\[
h(\tau) \to 0 as \tau \to \infty
\]

These are precisely the qualitative features associated with humped
forward-rate volatility loadings: little or no instantaneous short-end
volatility under a pinned boundary, maximal volatility at an
intermediate maturity, and decaying sensitivity at long maturities.

This interpretation should be stated carefully. The existence of a
humped admissible mode does not by itself prove that observed volatility
surfaces must take this form. It shows that such a shape is structurally
available within the IRFR field and can be selected by appropriate
boundary closure.

Empirical relevance remains a calibration question. Structural
admissibility is the mathematical result.

\section{Interior modes and boundary selection}

The analysis of this chapter reinforces the division introduced earlier.

The interior law determines the admissible maturity modes. These modes
arise from the KFE, moment closure, and coefficient matching. They are
mathematical consequences of the first three axioms.

The boundary law selects from these modes. It determines which
admissible combination is realised at the short end and how the field is
closed.

For the distinct-root branch, the boundary condition selects the
constants in the two-exponential combination,

\[
H(\tau) = C_1e^{-k_1\tau} + C_2e^{-k_2\tau}
\]

For the repeated-root branch, the boundary condition selects the
constants in

\[
H(\tau) = (C_1 + C_2\tau)e^{-k\tau}
\]

The Dirichlet condition, studied later, is especially important because
it naturally suppresses the constant term and leaves the humped
linear-exponential component.

Thus the boundary law does not merely impose an endpoint value. It
selects the maturity structure of risk.

\section{Relation to Ritchken--Chuang}

The repeated-root branch provides the structural bridge to the
Ritchken--Chuang humped volatility form.

In the conventional HJM setting, a humped volatility specification may
be introduced as a reduced-form maturity function. In the IRFR
formulation, a structurally similar humped profile appears as the
repeated-root mode of the admissibility equation.

The same qualitative object is therefore reinterpreted as an admissible
boundary-field mode rather than a freely imposed volatility shape.

This does not place the IRFR model outside HJM. The resulting dynamics
should still admit an HJM representation under translation. The point is
that the HJM representation is not structure-preserving: it does not
reveal why the humped form arises.

The IRFR interpretation is therefore:

\begin{center}\emph{humped volatility as boundary-selected admissible mode;}\end{center}
rather than

\begin{center}\emph{humped volatility as exogenous volatility input.}\end{center}
The full monetary-policy interpretation will require the
boundary-condition analysis developed later.

\section{Summary}

This chapter has extended the affine-closure analysis of Chapter 4 to
the maturity loading.

The KFE and moment equations do not allow arbitrary maturity loadings.
Coefficient matching imposes a second-order maturity equation. The
solutions of that equation define the admissible modes of the IRFR
field.

In the generic case, the characteristic roots are distinct, producing a
two-exponential branch,

\[
H(\tau) = C_1e^{-k_1\tau} + C_2e^{-k_2\tau}
\]

At the repeated-root value, identified under the later normalisation
with r = 1, the admissible loading becomes

\[
H(\tau) = (C_1 + C_2\tau)e^{-k\tau}
\]

The term

\[
\tau e^{-k\tau}
\]

is humped. It is therefore the structural origin of the humped
volatility branch that later connects to Ritchken--Chuang-type
specifications and to the Dirichlet boundary interpretation.

The main message is that humped maturity loading is not inserted by
hand. It arises as a repeated-root admissible mode of the IRFR field.

\section{Transition to boundary closure}

The next chapter turns from interior admissibility to boundary
selection.

The present chapter identified which maturity modes are allowed by the
interior law. It did not yet determine which mode is realised. That
selection requires Axiom 4: boundary closure at $\tau=0$.

The next chapter therefore studies boundary conditions as economic
regimes. It explains how natural, Dirichlet, and Robin closures select
different admissible combinations of the maturity modes derived here,
and why the Dirichlet repeated-root case plays a special role in the
monetary-policy interpretation of humped volatility.

\chapter{Boundary Closure and Economic Regimes}\label{chapter-6}

\section{Introduction}\label{introduction-3}

Chapter 5 derived the admissible maturity modes of the IRFR field. The
interior law, generated by Axioms 1-3, restricts the maturity loading to
a family of admissible modes. In the generic case, these modes are
combinations of distinct exponentials. At the repeated-root point,
identified with the r = 1 branch under the normalisation used in this
monograph, the admissible loading contains the linear-exponential
component

\[
\tau \exp(-a^2\tau)
\]
This humped component is an interior admissibility result. It is allowed
by the dynamics, but it is not yet selected.

The present chapter turns to selection.

The IRFR field is defined on the half-line

\[
\tau \in [0,\infty),
\]
and all fixed-maturity points are transported toward the short end,

\[
\tau = 0
\]
Thus the field is not fully determined by its interior law. A boundary
condition is required at the short end. This is the content of Axiom 4.

The key distinction is:

\begin{center}\emph{Axioms 1--3 determine admissible modes;}\end{center}
\begin{center}\emph{Axiom 4 selects the realised boundary regime.}\end{center}
Boundary closure is therefore not a technical afterthought. It is the
mechanism by which economic regime enters the IRFR model.

\section{Why boundary closure is necessary}\label{why-boundary-closure-is-necessary}

The maturity domain of the IRFR field is the half-line. Unlike a model
defined on the whole real line, a half-line field has an endpoint. That
endpoint is not arbitrary. It is the instantaneous short end of the term
structure.

The rolling-maturity identity

\[
d\tau/dt = -1
\]
means that every fixed calendar maturity moves toward this endpoint as
calendar time advances. The short end is therefore the location at which
the rolling field must be closed.

The interior dynamics derived in the previous chapters specify how the
field evolves for positive remaining maturity. They determine the
admissible forms of drift, diffusion, and maturity loading away from the
boundary. But they do not, by themselves, determine what happens at $\tau$ =
0.

This is a standard feature of boundary-value problems. An interior
differential equation determines a space of possible solutions. A
boundary condition selects the particular solution, or the particular
family of solutions, compatible with the endpoint constraint.

In the IRFR setting, this mathematical fact has direct economic meaning.
The short end is where policy, funding conditions, and overnight-rate
mechanisms enter most directly. Boundary closure is therefore both a
mathematical necessity and an economic modelling choice.

\section{Interior admissibility versus boundary selection}\label{interior-admissibility-versus-boundary-selection}

The distinction between admissibility and selection is central.

Interior admissibility answers the question:

\begin{center}\emph{Which maturity loadings are dynamically allowed?}\end{center}
Boundary selection answers the question:

\emph{Which admissible loading is realised under a given short-end
regime?}

Chapter 5 answered the first question. It showed that admissible
loadings arise as solutions of a second-order maturity equation. In the
distinct-root case,

\[
H(\tau) = C_1 \exp(-k_1\tau) + C_2 \exp(-k_2\tau)
\]
In the repeated-root case,

\[
H(\tau) = (C_1 + C_2\tau) \exp(-k\tau)
\]
The constants $C_1$ and C$_2$ are not fixed by the interior equation alone.
They encode the remaining freedom after admissibility has been imposed.
Boundary closure fixes, constrains, or relates these constants.

Thus Axiom 4 does not create the admissible modes. It selects among
them.

This point prevents a common misunderstanding. The humped repeated-root
loading is not imposed by the boundary condition from nothing. It first
appears as an admissible interior mode. The boundary condition
determines whether that mode is selected, suppressed, or mixed with
another admissible component.

\section{The short end as an economic boundary}\label{the-short-end-as-an-economic-boundary}

The short end of the yield curve has a special economic role.

Central banks operate most directly through short-maturity rates:
overnight rates, reserve remuneration, standing facilities, policy
corridors, and short-end guidance. These mechanisms do not usually set
the entire forward curve directly. Instead, they act at or near the
short end, and their effects propagate through the maturity structure.

In the IRFR formulation, this observation has a natural mathematical
expression. Monetary policy enters through the boundary condition at

\[
\tau = 0
\]
This does not mean that every monetary-policy regime literally fixes the
instantaneous short rate. Different policy regimes correspond to
different forms of boundary closure. Some regimes leave the short end
relatively free. Others pin it tightly. More realistic regimes partially
anchor the short end while allowing residual stochastic variation.

The boundary condition is therefore the mathematical location of the
policy regime.

This is the sense in which monetary policy is represented as a boundary
law. It is not introduced as a global curve-wide force. It is introduced
as a closure rule at the endpoint of the maturity-domain field.

\section{Three canonical boundary regimes}\label{three-canonical-boundary-regimes}

The monograph focuses on three canonical boundary regimes:

\[
Neumann/free closure;
\]
\begin{center}\emph{Dirichlet fixed-boundary closure;}\end{center}
\begin{center}\emph{Robin partial-anchoring closure.}\end{center}
These regimes should be read as structural idealisations, not literal
institutional descriptions. They are not claimed to exhaust all possible
IRFR boundary laws. They provide the three canonical cases studied in
this monograph: a non-pinned benchmark, a pinned boundary, and a
partially anchored boundary.

\begin{longtable}[]{@{}
  >{\raggedright\arraybackslash}p{(\columnwidth - 6\tabcolsep) * \real{0.2500}}
  >{\raggedright\arraybackslash}p{(\columnwidth - 6\tabcolsep) * \real{0.2500}}
  >{\raggedright\arraybackslash}p{(\columnwidth - 6\tabcolsep) * \real{0.2500}}
  >{\raggedright\arraybackslash}p{(\columnwidth - 6\tabcolsep) * \real{0.2500}}@{}}
\toprule\noalign{}
\begin{minipage}[b]{\linewidth}\raggedright
\textbf{Boundary regime}
\end{minipage} & \begin{minipage}[b]{\linewidth}\raggedright
\textbf{Representative loading / condition}
\end{minipage} & \begin{minipage}[b]{\linewidth}\raggedright
\textbf{Economic interpretation}
\end{minipage} & \begin{minipage}[b]{\linewidth}\raggedright
\textbf{Mode effect}
\end{minipage} \\
\midrule\noalign{}
\endhead
\bottomrule\noalign{}
\endlastfoot
Neumann/free & $H_N(\tau)=A_N\exp(-a_N^2\tau)$ & non-pinned benchmark &
retains boundary-active exponential loading \\
Dirichlet & $H(0)=0$ & pinned short end & suppresses $C_1$ in the
repeated-root case \\
Robin & $\alpha H(0)+\beta H'(0)=\gamma$ & partial anchoring & mixes admissible
components \\
\end{longtable}

Neumann/free closure\footnote{The term Neumann/free denotes the non-pinned boundary benchmark in which the boundary-active exponential loading is retained. The word ``Neumann'' is used in this boundary-regime sense only: it should not be read as imposing the classical homogeneous Neumann condition $H'(0)=0$ on the stochastic loading. The operative condition in the present monograph is non-pinning rather than zero derivative.} represents the benchmark in which no explicit
monetary-policy intervention is imposed through the boundary law. It
should not be read as the absence of monetary institutions or policy
influence in the economy. Rather, it is the case in which the short end
remains stochastic and is closed without being externally pinned.

A Dirichlet boundary represents an idealised fixed-policy regime. The
short end is pinned at a prescribed value. This is an exaggerated policy
intervention, but it is mathematically important because it reveals how
a fixed short-end boundary selects admissible modes.

A Robin boundary represents partial anchoring. It links the value of the
loading and its derivative at the boundary. Economically, this
corresponds to a short end influenced by policy but not perfectly fixed.

The rest of this chapter explains how these canonical closures act on
the admissible mode space.

\section{Neumann/free closure}\label{naturalfree-closure}

For the loading-level analysis developed in this monograph, the natural
benchmark is represented by the free exponential short-end stochastic
mode

\[
H_N(\tau)=A_N e^{-a_N^2\tau}
\]

This is the retained exponential loading for the non-pinned benchmark. The terminology follows the convention fixed in Chapter 6.

This loading satisfies

\[
H_N(0) = A_N
\]
Thus, unless $A_N=0$, the short end retains stochastic loading. The
boundary is closed, but it is not pinned.

This choice should be understood carefully. It is not a theorem that
every possible natural IRFR closure must take this exact exponential
form. A full state-space formulation could express Neumann/free closure
through a flux, regularity, or positivity-preserving endpoint condition.
Rather, the exponential free-boundary loading is the benchmark used here
to define the non-pinned regime against which Dirichlet and Robin
closure are compared.

Economically, this Neumann/free benchmark represents the case of no
explicit monetary-policy intervention through the boundary law. It is
not a claim that monetary policy is absent in an institutional sense. It
means that the model imposes no additional policy-pinning condition at
$\tau=0$. The short end remains part of the stochastic field rather than
being externally pinned.

This regime is therefore a non-pinned benchmark closure. It allows
stochastic exposure to enter immediately at the short end, rather than
suppressing the boundary-active component.

\section{Dirichlet closure}\label{dirichlet-closure}

A Dirichlet boundary condition fixes the value of the field or loading
at the boundary. In schematic form,

\[
H(0)=0,
\]
or, in terms of the short-end field,

\begin{center}\(x(t,0)=\text{prescribed boundary value}.\)\end{center}
The exact form depends on whether the boundary condition is imposed on
the field level, the displacement from equilibrium, or the maturity
loading. But the structural meaning is the same: the boundary value is
fixed.

Economically, the Dirichlet case represents an idealised pinned
short-end regime. It is not intended as a literal description of all
monetary-policy operations. Rather, it is a limiting case in which the
short end is treated as fixed by policy.

This case is especially important for the repeated-root branch.

Recall that the repeated-root admissible loading has the form

\[
H(\tau) = (C_1 + C_2\tau) \exp(-k\tau)
\]
At the boundary,

\[
H(0)=C_1
\]
A Dirichlet condition of the form

\[
H(0)=0
\]
therefore sets

\[
C_1=0
\]
The realised loading becomes

\[
H(\tau) = C_2\tau \exp(-k\tau)
\]
Thus the Dirichlet boundary selects the humped linear-exponential
component.

This is the first explicit example of boundary selection. The
repeated-root mode is admissible by the interior law. The Dirichlet
boundary selects the humped component by suppressing the constant
exponential term.

\section{The Dirichlet repeated-root branch}\label{the-dirichlet-repeated-root-branch}

The Dirichlet repeated-root case is central to the monograph.

Under the r=1 repeated-root branch, the admissible loading is

\[
H(\tau) = (C_1 + C_2\tau) \exp(-a^2\tau)
\]
Imposing the Dirichlet condition

\[
H(0)=0
\]
gives

\[
C_1=0,
\]
and hence

\[
H(\tau) = C_2\tau \exp(-a^2\tau)
\]
This loading is humped. It vanishes at the boundary, rises to an
interior maximum, and decays at long maturities.

The maximum occurs at

\[
\tau=1/a^2
\]
This is the structural source of the humped volatility branch.

In a conventional HJM representation, a humped volatility loading of
this type may appear as an exogenous maturity function. In the IRFR
formulation, it is selected by the combination of interior admissibility
and boundary closure:

\begin{center}\emph{\(r=1\) repeated-root admissibility;}\end{center}
\begin{center}\emph{Dirichlet short-end boundary closure.}\end{center}
The humped profile is therefore not introduced by hand. It is a
boundary-selected admissible mode.

\section{Markovianity and representation}\label{markovianity-and-representation}

The Dirichlet repeated-root branch also clarifies the Markovianity claim
made in Chapter 1.

The claim is not finite-dimensional Markovianity in the ordinary
short-rate sense. It is Markovianity of the boundary-value field
evolution.

This distinction matters. In standard finite-dimensional HJM
coordinates, a stationary humped volatility profile of this type is
often treated as non-Markovian. The hump appears as an externally
specified maturity structure that cannot generally be represented by a
low-dimensional Markov state without additional machinery.

In the IRFR formulation, the same humped structure appears differently.
It is a boundary-selected mode of a rolling maturity field. The state is
the field, the short end is the boundary, and the humped loading is
generated by the boundary-value structure.

The representation matters.

What appears as non-Markovian structure in one coordinate system may
become a local boundary condition in another. This is why the IRFR
geometry is not merely notational. It changes which features of the
model are primitive and which are derived.

\section{Robin closure}\label{robin-closure}

A Robin boundary condition mixes the value of the field and its
derivative at the boundary. In schematic form,

\[
\alpha H(0)+\beta H'(0)=\gamma,
\]
where $\alpha$, $\beta$, and $\gamma$ encode the boundary regime.

Economically, this represents partial anchoring. The short end is
influenced by policy, but not perfectly pinned. There remains residual
stochastic movement at the boundary.

The Robin case is therefore the natural bridge between the Neumann/free
benchmark and the fixed boundary. It allows a continuum of regimes
between a free short end and a fully pinned short end.

In the mode language, the Robin condition selects a mixture of
admissible modes. For the repeated-root branch,

\[
H(\tau) = (C_1 + C_2\tau) \exp(-k\tau),
\]
the boundary value and derivative are

\[
H(0)=C_1,
\]
and

\[
H'(0)=C_2-kC_1
\]
A Robin condition therefore imposes a linear relation between $C_1$ and C$_2$.
Unlike the Dirichlet case, it need not suppress the constant term
completely. Instead, it selects a mixed loading containing both
exponential and humped components.

This is economically important. A realistic policy regime may partially
anchor the short end without eliminating all short-end variation. Robin
closure captures this intermediate regime.

\section{Boundary law as policy transmission}\label{boundary-law-as-policy-transmission}

The boundary law determines how short-end policy enters the admissible
maturity structure.

In a conventional reduced-form model, policy effects may be introduced
through shifts in the short rate, changes in factor dynamics, or changes
in volatility specification. In the IRFR model, policy enters through
the boundary condition.

This has two consequences.

First, policy is local at the point of intervention. It enters at the
short-end boundary rather than as a global curve-wide forcing term.

Second, the effect of policy is transmitted through the admissible
maturity modes of the field. The boundary does not directly prescribe
the entire curve. It selects a mode or combination of modes, and the
interior law propagates that selection across maturity.

This is the boundary-law interpretation of monetary policy.

The short end is where policy acts most directly. The maturity field is
the structure through which that action is transmitted.

\section{Relation to the r-branch}\label{relation-to-the-r-branch}

Boundary closure interacts with the r-branch developed in Chapter 5.

The r-branch determines the admissible interior mode structure. The
boundary condition determines the realised combination.

For $r\neq1$, the admissible loading is a combination of distinct exponential
modes. Boundary closure selects the constants in that combination.

For r=1, the repeated-root loading is

\[
H(\tau) = (C_1 + C_2\tau) \exp(-a^2\tau)
\]
The Neumann/free benchmark retains the boundary-active exponential
stochastic mode,

\[
H_N(\tau) = A_N \exp(-a_N^2\tau)
\]
The Dirichlet condition selects the humped component

\[
H(\tau) \propto \tau \exp(-a^2\tau)
\]
The Robin condition selects a mixture of the constant exponential and
humped components.

Thus the r-branch determines the available mode structure, while the
boundary law determines the economic realisation of that structure.

\section{Summary}\label{summary-2}

This chapter has introduced boundary closure as the selection mechanism
for the IRFR field.

The interior law, derived from Axioms 1-3, determines the admissible
maturity modes. Axiom 4 determines how the field is closed at the short
end. This closure selects the realised combination of admissible modes.

The main boundary regimes are Neumann/free closure, Dirichlet
fixed-boundary closure, and Robin partial-anchoring closure.

The Neumann/free benchmark retains the boundary-active exponential
loading

\[
H_N(\tau) = A_N \exp(-a_N^2\tau)
\]
The Dirichlet repeated-root case is especially important. Under r=1, the
Dirichlet condition suppresses the constant exponential term and leaves
the humped loading

\[
H(\tau) \propto \tau \exp(-a^2\tau)
\]
The main message is:

\begin{center}\emph{interior admissibility creates the mode space;}\end{center}
\begin{center}\emph{boundary closure selects the realised mode.}\end{center}
\section{Transition to detailed boundary cases}\label{transition-to-detailed-boundary-cases}

The next chapters study the boundary regimes in detail.

The Neumann/free case provides the non-pinned benchmark closure. The
Dirichlet case provides the idealised pinned-policy limit and selects
the repeated-root humped branch. The Robin case provides the more
realistic partially anchored policy regime.

These boundary cases will show how the same admissible interior law can
generate different observable maturity loadings depending on how the
short end is closed.

\chapter{Neumann/free Boundary Closure and the Free Benchmark}\label{chapter-7}

\section{Introduction}\label{introduction-4}

Chapter 6 introduced boundary closure as the selection mechanism for the
IRFR field. The interior law determines the admissible maturity modes,
while the boundary law determines how the short end is closed.

This chapter studies the benchmark case: natural, or free, boundary
closure.

The natural boundary is the regime in which the short end is not
externally pinned. The field is allowed to meet the boundary without
imposing a fixed short-end value. In this sense, Neumann/free closure
represents the non-pinned benchmark against which the Dirichlet and
Robin regimes are compared.

Neumann/free closure leaves the short end stochastic;

Dirichlet closure pins the short end;

Robin closure partially anchors the short end.

The purpose of the present chapter is to establish the first case
clearly. Natural closure provides the baseline regime in which the short
end remains active and stochastic. It does not select the pure humped
repeated-root branch in the same way as the Dirichlet condition.
Instead, it retains the ordinary boundary-active stochastic mode.

This benchmark is important for two reasons. First, it separates
interior admissibility from policy selection. A humped repeated-root
mode may be admissible, but it is not automatically realised. Its
selection depends on how the short end is closed. Second, it clarifies
what monetary policy changes. A pinned or partially anchored policy
regime is meaningful only relative to a benchmark in which the short end
remains free.

\section{What ``natural'' means}\label{what-natural-means}

For the loading-level analysis developed in this monograph, natural
boundary closure is represented by the Neumann/free benchmark in which
the short end remains part of the stochastic field.

\[
H_N(\tau) = A_N e^{-a_N^2\tau}
\]

This loading satisfies

\[
H_N(0) = A_N
\]

Thus, unless $A_N=0$, the short end retains stochastic loading. The
boundary is closed, but it is not pinned.

This choice should be understood carefully. It is not a theorem that
every possible natural IRFR closure must take this exact exponential
form. A full state-space formulation could express Neumann/free closure
through a flux, regularity, or positivity-preserving endpoint condition.
Rather, the exponential free-boundary loading is the benchmark used here
to define the non-pinned regime against which Dirichlet and Robin
closure are compared.

Economically, this natural benchmark represents the case of no explicit
monetary-policy intervention through the boundary law. It is not a claim
that monetary policy is absent in an institutional sense. It means that
the model imposes no additional policy-pinning condition at $\tau=0$. The
short end remains part of the stochastic field rather than being
externally pinned.

\section{The short end under natural
closure}\label{the-short-end-under-natural-closure}

The short end $\tau=0$ is economically distinguished, but Neumann/free closure
does not treat it as fixed.

Instead, the short end remains part of the stochastic field. It may move
with the same underlying Brownian source as the rest of the curve,
subject to the admissibility restrictions of the IRFR model.

This matters because the short end is where the difference between free
and policy regimes becomes visible. Under a natural boundary, the model
does not impose a fixed policy value at the boundary.

\[
x(t,0) = fixed policy value
\]

is not imposed. Nor is the loading-level condition

\[
H(0) = 0
\]

imposed. The short end is closed, but not pinned. This distinction will
be important in the next chapter. Dirichlet closure suppresses the
boundary-active stochastic component. Natural closure, by contrast,
retains it.

\section{Natural closure and probability
flux}\label{natural-closure-and-probability-flux}

One way to understand Neumann/free closure is through flux.

The KFE derived in Chapter 4 can be written in terms of probability flow
over the state variable. Under the convention

\[
dx=-f\,dt+\sqrt{2}\,g\,dW_t,
\]

the KFE is

\[
\frac{\partial p}{\partial t} = \partial(fp)/\partial x + \partial^2(g^2p)/\partial x^2
\]

This may be written as a conservation equation,

\[
\frac{\partial p}{\partial t} = -\frac{\partial J}{\partial x},
\]

where J is the probability flux in the state variable. With the above
sign convention,

\[
J = -fp - \partial(g^2p)/\partial x
\]

A natural boundary condition is often associated with vanishing boundary
flux. In the state variable, this means that probability does not leak
out of the admissible state domain.

This flux is a flux in the state variable x, not along the maturity
coordinate $\tau$. The analogous maturity-domain idea is that the field meets
the short-end boundary through an internally compatible endpoint
condition rather than through an externally imposed boundary value.

The precise flux or endpoint condition depends on the state domain and
on the positivity or admissibility constraints imposed on the IRFR
density. For the present chapter, the important point is conceptual:
Neumann/free closure is a compatibility condition, not a pinned-value
condition. It allows the field to remain stochastic at the boundary
while preserving admissibility.

\section{Natural closure and the free stochastic
mode}\label{natural-closure-and-the-free-stochastic-mode}

The Neumann/free benchmark retains the ordinary exponential stochastic
mode. In the r=1 benchmark, the loading is

\[
H_N(\tau) = A_N e^{-a_N^2\tau}
\]

The associated instantaneous variance rate is

\[
\partial \operatorname{Var}(x^N_{t,t+\tau})/\partial t = 2a_N^2 A_N^2 e^{-2a_N^2\tau}
\]

This profile is monotone in maturity and non-zero at the short end
whenever $A_N\neq0$. It therefore represents a regime in which short-end
stochastic exposure is retained.

This is the important distinction from Dirichlet closure. Natural
closure does not suppress the boundary-active exponential component. It
allows stochastic exposure to enter immediately at the short end.

\section{Benchmark interpretation}\label{benchmark-interpretation}

The Neumann/free benchmark is not the pure humped repeated-root branch.
Its loading is

\[
H_N(\tau) = A_N e^{-a_N^2\tau}
\]

By contrast, the Dirichlet repeated-root loading is

\[
H_D^{(1)}(\tau) = C_D \tau e^{-a_D^2\tau}
\]

The two regimes therefore differ sharply at the boundary:

\[
H_N(0) = A_N,
\]

\[
H_D^{(1)}(0) = 0
\]

Natural closure retains short-end loading. Dirichlet closure suppresses
it.

This makes the economic distinction transparent. Under Neumann/free closure,
the short end remains stochastic. Under Dirichlet closure, the short-end
stochastic loading is removed by the pinned boundary.

\section{Natural closure versus Dirichlet
selection}\label{natural-closure-versus-dirichlet-selection}

The comparison with Dirichlet closure is instructive.

The Neumann/free benchmark has loading

\[
H_N(\tau) = A_N e^{-a_N^2\tau}
\]

The Dirichlet repeated-root branch has loading

\[
H_D^{(1)}(\tau) = C_D \tau e^{-a_D^2\tau}
\]

The former is boundary-active, while the latter is boundary-suppressed:

$H_N(0)=A_N,\quad H_D^{(1)}(0)=0$.

\begin{longtable}[]{@{}
  >{\raggedright\arraybackslash}p{(\columnwidth - 6\tabcolsep) * \real{0.2500}}
  >{\raggedright\arraybackslash}p{(\columnwidth - 6\tabcolsep) * \real{0.2500}}
  >{\raggedright\arraybackslash}p{(\columnwidth - 6\tabcolsep) * \real{0.2500}}
  >{\raggedright\arraybackslash}p{(\columnwidth - 6\tabcolsep) * \real{0.2500}}@{}}
\toprule\noalign{}
\begin{minipage}[b]{\linewidth}\raggedright
\textbf{Boundary regime}
\end{minipage} & \begin{minipage}[b]{\linewidth}\raggedright
\textbf{Loading}
\end{minipage} & \begin{minipage}[b]{\linewidth}\raggedright
\textbf{Short-end loading}
\end{minipage} & \begin{minipage}[b]{\linewidth}\raggedright
\textbf{Economic meaning}
\end{minipage} \\
\midrule\noalign{}
\endhead
\bottomrule\noalign{}
\endlastfoot
Neumann/free & $A_N e^{-a_N^2\tau}$ & retained & non-pinned
benchmark \\
Dirichlet & $C_D \tau e^{-a_D^2\tau}$ & suppressed & pinned boundary \\
\end{longtable}

This is the point of boundary law. The short-end regime changes the
realised loading, and therefore changes how stochastic sensitivity is
distributed across maturity.

The difference has economic meaning. Dirichlet closure removes short-end
stochastic loading and represents an idealised pinned policy boundary.
Natural closure preserves short-end variation and represents the
non-pinned benchmark.

\section{Natural closure and the benchmark
economy}\label{natural-closure-and-the-benchmark-economy}

The natural boundary is best interpreted as a benchmark economy rather
than as a literal absence of all policy.

No real interest-rate system is free of institutional structure.
Overnight markets, collateral systems, reserve regimes, and central-bank
operating frameworks always matter. Nevertheless, the natural boundary
is useful because it describes the limiting case in which the short end
is not pinned by an externally prescribed value.

It is the regime in which the field is closed by internal admissibility
rather than by policy anchoring. This makes it the appropriate baseline
for the monetary-policy extensions.

Neumann/free closure: free short end;

Dirichlet closure: pinned short end;

Robin closure: partially anchored short end.

The benchmark therefore defines what policy intervention changes. It
changes the boundary law.

\section{Consequences for volatility
shape}\label{consequences-for-volatility-shape}

Natural closure does not imply a flat volatility structure. The
Neumann/free benchmark has a non-trivial maturity profile, but it is the
boundary-active exponential profile rather than the pinned-boundary
hump.

$H_N(\tau)=A_Ne^{-a_N^2\tau}$.

The pure humped Dirichlet profile

\[
\tau e^{-a^2\tau}
\]

is specifically associated with suppressing the boundary-active
component and selecting the repeated-root boundary mode.

Natural closure can produce decay and maturity dependence. But it does
not express the same economic content as a pinned short end. The
short-end loading remains present.

This matters for empirical interpretation. Observed volatility near the
short end contains information about the boundary law, not merely about
the interior dynamics.

\section{Relation to monetary
policy}\label{relation-to-monetary-policy}

The natural boundary provides the reference point for monetary-policy
interpretation.

In a free-boundary benchmark, the short end remains stochastic. Policy
is not represented as fixing the boundary value. The field evolves
according to the interior admissibility structure and the natural
closure condition.

When policy becomes active in the boundary-law sense, the closure
changes. A Dirichlet boundary represents the limiting case in which the
short end is pinned. A Robin boundary represents the more realistic case
in which the short end is partially anchored.

Thus monetary policy is not modelled as a new interior force. It is
modelled as a change in the boundary law.

interior dynamics describe admissible propagation;

boundary law describes regime selection.

The natural boundary is the non-pinned benchmark from which policy
anchoring is measured.

\section{Summary}\label{summary-3}

This chapter has studied natural boundary closure as the benchmark
regime of the IRFR field.

Natural closure is required because the maturity domain is the half-line
and the short end is a boundary. But unlike Dirichlet closure, it does
not prescribe a fixed boundary value. It allows the short end to remain
stochastic.

In the loading-level benchmark used here, Neumann/free closure retains
the ordinary exponential short-end stochastic mode,

$H_N(\tau)=A_Ne^{-a_N^2\tau}$.

This contrasts with the Dirichlet repeated-root loading,

$H_D^{(1)}(\tau)=C_D\tau e^{-a_D^2\tau}$,

which suppresses short-end loading.

Neumann/free closure closes the field without pinning the short end.

\section{Transition to Dirichlet
closure}\label{transition-to-dirichlet-closure}

The next chapter studies the Dirichlet boundary in detail.

Dirichlet closure fixes the boundary value at the short end. In the
repeated-root r=1 branch, this suppresses the boundary-active stochastic
component and selects the humped loading

\[
H(\tau) \propto \tau e^{-a^2\tau}
\]

This case represents the idealised pinned-policy limit. It is
mathematically simple, economically extreme, and central to the
connection between IRFR boundary law and humped volatility.

\chapter{Dirichlet Boundary Closure and the r = 1 Humped Branch}\label{chapter-8}

\section{Introduction}\label{introduction-5}

Chapter 7 studied Neumann/free boundary closure as the benchmark regime
of the IRFR field. Natural closure closes the maturity-domain field
without pinning the short end. The short end remains stochastic, and the
boundary-active exponential loading is retained.

This chapter studies the opposite limiting case: Dirichlet boundary
closure.

A Dirichlet boundary fixes the value of the field, or of the relevant
displacement or loading, at the short-end boundary. In schematic loading
form,

\[
H(0)=0,
\]

or, at the field level,

\begin{center}\(x(t,0)=\text{prescribed boundary value}.\)\end{center}

The exact object to which the condition is applied depends on the
formulation. The structural meaning is the same: the boundary value is
fixed.

Economically, the Dirichlet case represents an idealised pinned-policy
regime. It is not intended as a literal description of all central-bank
operations. It is a limiting case in which the short end is treated as
fixed by policy. Precisely because it is extreme, it is analytically
revealing.

The central result of the chapter is that, in the repeated-root r=1
branch, Dirichlet closure selects the humped linear-exponential loading

\[
H(\tau) \propto \tau e^{-a^2\tau}
\]

The humped loading is therefore not an exogenous volatility input. It is
the boundary-selected survivor of the repeated-root admissible mode
space.

\section{Dirichlet closure as fixed-boundary
selection}\label{dirichlet-closure-as-fixed-boundary-selection}

The IRFR maturity domain is the half-line

\[
\tau \in [0,\infty).
\]

The endpoint $\tau=0$ is the instantaneous short end. A Dirichlet condition
fixes the value of the relevant boundary object at that endpoint.

In the simplest loading formulation, the condition is

\[
H(0)=0
\]

This says that the maturity loading vanishes at the boundary. In
economic terms, the short-end stochastic loading is suppressed. The
boundary is not merely closed; it is pinned.

At the field level, a Dirichlet condition may instead be written as

\[
x(t,0)=x_0(t),
\]

where $x_0(t)$ is a prescribed short-end boundary value. In the most
idealised fixed-policy case, $x_0(t)$ may be constant over the relevant
horizon. More generally, it may be a deterministic policy path.

In the loading-level analysis below, we use the corresponding
homogeneous boundary condition H(0)=0, which represents the suppression
of short-end stochastic loading.

For the mode-selection argument, the essential feature is that the
boundary value is prescribed externally. This distinguishes Dirichlet
closure from Neumann/free closure, which retains the boundary-active
stochastic component.

\section{Economic interpretation: the pinned short
end}\label{economic-interpretation-the-pinned-short-end}

The Dirichlet case should be interpreted as an idealised pinned
short-end regime.

Central banks act most directly at the short end. In a policy-corridor,
administered-rate, or fixed-policy-rate regime, the shortest maturities
may be strongly anchored relative to the rest of the curve. The
Dirichlet condition abstracts this anchoring into a fixed boundary
value.

This abstraction is deliberately sharp. It suppresses short-end
stochastic loading and asks what maturity structure remains admissible
in the interior. The boundary condition therefore acts as a selection
rule on the admissible maturity modes.

This is why the Dirichlet case is central to the monetary-policy
extension. It provides the cleanest mathematical representation of
extreme short-end policy anchoring.

\section{Action on the repeated-root mode
space}\label{action-on-the-repeated-root-mode-space}

The repeated-root branch has admissible loading

\[
H(\tau)=(C_1+C_2\tau)e^{-k\tau}
\]

Under the normalisation used for the r=1 branch,

\[
k=a^2,
\]

so

\[
H(\tau)=(C_1+C_2\tau)e^{-a^2\tau}
\]

The two independent components are

$e^{-a^2\tau}$

and

\[
\tau e^{-a^2\tau}
\]

The first component is non-zero at the boundary, while the second
component vanishes at the boundary:

$e^{-a^2\cdot 0}=1$, $\tau e^{-a^2\tau}|_{\tau=0}=0$.

Thus the two repeated-root components behave differently at the short
end. The Dirichlet condition distinguishes precisely between them.

A Dirichlet condition of the form H(0)=0 suppresses the constant
exponential component and leaves the linear-exponential component. This
is the algebraic mechanism by which the Dirichlet boundary selects the
humped branch.

\section{The selection calculation}\label{the-selection-calculation}

Start with the repeated-root admissible loading

\[
H(\tau)=(C_1+C_2\tau)e^{-a^2\tau}
\]

Evaluating at the boundary gives

\[
H(0)=C_1
\]

The Dirichlet condition H(0)=0 therefore implies

$C_1=0$.

The realised loading is

\[
H(\tau)=C_2\tau e^{-a^2\tau}
\]

Up to normalisation,

\[
H(\tau)\propto\tau e^{-a^2\tau}
\]

The humped term was already present in the interior admissible mode
space. The Dirichlet boundary did not create it from nothing. It
selected it by eliminating the boundary-active constant component.

\section{Shape of the selected
loading}\label{shape-of-the-selected-loading}

Ignoring the normalisation constant, define

\[
h(\tau)=\tau e^{-a^2\tau}
\]

Then $h(0)=0$ and $h(\tau)\to0$ as $\tau\to\infty$. The derivative is

$h'(\tau)=e^{-a^2\tau}(1-a^2\tau)$.

The maximum occurs when

$1-a^2\tau=0$,

so

\[
\tau_{\max}=1/a^2
\]

Thus the selected loading vanishes at the short end, reaches an interior
maximum, and decays at long maturity. This is the qualitative structure
of a humped forward-rate volatility loading.

The location of the maximum is controlled by $a^2$. A larger $a^2$ shifts the
hump toward shorter maturities. A smaller $a^2$ shifts it toward longer
maturities. Thus $a^2$ controls not only the strength of local mean
reversion in the affine drift operator, but also the maturity scale of
the selected Dirichlet mode.

\section{Positivity and the lower admissible
boundary}\label{positivity-and-the-lower-admissible-boundary}

The Dirichlet repeated-root branch also has an important admissibility
property. The selected loading

$H_D^{(1)}(\tau)=C_D\tau e^{-a_D^2\tau}$

is non-negative on the half-line whenever $C_D\geq0$. It satisfies

\[
H_D^{(1)}(0)=0, H_D^{(1)}(\tau)\ge0 for \tau\ge0, H_D^{(1)}(\tau)\to0 as \tau\to\infty
\]

Thus the boundary-selected loading does not force a sign change across
maturity. In the Dirichlet regime, the short-end boundary can therefore
be interpreted not only as a pinned policy value, but also as a lower
admissible boundary for the rolling field.

Write the field relative to a minimum admissible level:

\[
x(t,\tau)=x_{\min}+y(t,\tau), y(t,\tau)\ge0
\]

The Dirichlet condition is then imposed on the displacement,

\[
y(t,0)=0
\]

For each fixed maturity $\tau>0$, the shifted Dirichlet dynamics
remain in the positive-state multiplicative KFE class considered in
Appendix B. In schematic form,

\[
dy_t=\kappa_D(\tau)(\theta_D(\tau)-y_t)dt+\sqrt{2} \eta_D(\tau)y_t\,dW_t,
\]

with positive parameter functions. The stationary KFE then admits a
shifted inverse-gamma density on y\textgreater0, equivalently on
x\textgreater $x_{\min}$. The same Feller boundary classification as in the
positive-state benchmark shows that the lower boundary y=0 is
inaccessible for each fixed $\tau>0$, provided the process starts
in the positive state domain and the parameter conditions are satisfied.

The lower-boundary interpretation is therefore stronger than a static
sign condition on the loading. It is supported by the shifted
inverse-gamma stationary density and by maturity-by-maturity Feller
inaccessibility of the lower state boundary. The detailed density and
boundary-classification calculation is given in Appendix B.

The remaining qualification is field-level regularity. To pass from
fixed-$\tau$ positivity to the full maturity-field statement

\[
y(t,\tau)\ge0 for all \tau\ge0,
\]

one also requires sufficient regularity of the field in the maturity
coordinate. Under such regularity, the Dirichlet boundary gives y(t,0)=0
while the shifted process remains positive for each fixed
$\tau>0$.

This should be read as an admissibility statement about the IRFR field
under the chosen boundary normalisation. It is not a claim that nominal
market rates can never be negative in every institutional setting.

\section{Comparison with Neumann/free
closure}\label{comparison-with-naturalfree-closure}

The contrast with Neumann/free closure is immediate.

The Neumann/free benchmark retains the boundary-active exponential
stochastic mode:

$H_N(\tau)=A_Ne^{-a_N^2\tau}$.

The short-end loading remains present:

$H_N(0)=A_N$.

Under Dirichlet closure, the selected repeated-root loading is

$H_D^{(1)}(\tau)=C_D\tau e^{-a_D^2\tau}$.

The short-end loading is suppressed:

$H_D^{(1)}(0)=0$.

The difference is not merely a difference in shape. It is a difference
in boundary selection.

\begin{longtable}[]{@{}
  >{\raggedright\arraybackslash}p{(\columnwidth - 6\tabcolsep) * \real{0.2500}}
  >{\raggedright\arraybackslash}p{(\columnwidth - 6\tabcolsep) * \real{0.2500}}
  >{\raggedright\arraybackslash}p{(\columnwidth - 6\tabcolsep) * \real{0.2500}}
  >{\raggedright\arraybackslash}p{(\columnwidth - 6\tabcolsep) * \real{0.2500}}@{}}
\toprule\noalign{}
\begin{minipage}[b]{\linewidth}\raggedright
\textbf{Boundary regime}
\end{minipage} & \begin{minipage}[b]{\linewidth}\raggedright
\textbf{Loading}
\end{minipage} & \begin{minipage}[b]{\linewidth}\raggedright
\textbf{Short-end loading}
\end{minipage} & \begin{minipage}[b]{\linewidth}\raggedright
\textbf{Economic meaning}
\end{minipage} \\
\midrule\noalign{}
\endhead
\bottomrule\noalign{}
\endlastfoot
Neumann/free & $A_N e^{-a_N^2\tau}$ & retained & non-pinned
benchmark \\
Dirichlet & $C_D \tau e^{-a_D^2\tau}$ & suppressed & pinned boundary \\
\end{longtable}

This is why the Dirichlet boundary has a natural interpretation as a
pinned-policy limit.

\section{Markovianity of the Dirichlet humped
branch}\label{markovianity-of-the-dirichlet-humped-branch}

The Dirichlet repeated-root branch clarifies the Markovianity issue.

In a standard HJM representation, a stationary humped volatility loading
of the form

\[
\tau e^{-a^2\tau}
\]

may appear as an exogenous maturity-dependent volatility specification.
In finite-dimensional HJM realisation theory, such a specification is
often treated as difficult to represent by a low-dimensional Markov
state.

The IRFR interpretation is different. In IRFR coordinates, the humped
loading is not an arbitrary maturity function. It is a boundary-selected
admissible mode of a rolling maturity field. The state is the field, the
short end is the boundary, and the boundary condition is local.

Thus the claim is not that the model is one-dimensional Markov in the
ordinary short-rate sense. It is not. The claim is that the Dirichlet
humped branch is Markovian as a boundary-value evolution of the IRFR
field.

The Markovianity claim should be understood in the field-state sense.
The state at time \(t\) is not a single short rate, nor a finite vector of
HJM factors, but the rolling field

$X_t=\{x(t,\tau):\tau\geq0\}$,

together with the boundary regime defining the domain of admissible
evolution. With respect to the filtration generated by the driving
Brownian motion and the current field state, the future law of
$X_{t+s}$ is determined by $X_t$, the interior operator, and the
boundary condition. In this sense the Dirichlet repeated-root branch is
Markovian as an infinite-dimensional boundary-value field evolution.
This is weaker than finite-dimensional Markovianity, and it is the only
Markovianity claim being made here.

\section{Relation to Ritchken--Chuang-type
volatility}\label{relation-to-ritchken-chuang-type-volatility}

The selected Dirichlet loading

\[
H(\tau)\propto\tau e^{-a^2\tau}
\]

has the same qualitative form as the humped maturity loading associated
with Ritchken--Chuang-type volatility specifications.

In the conventional use of such specifications, the humped form is
introduced directly as a maturity function. The IRFR interpretation is
different. The same qualitative form appears as the result of boundary
selection: r=1 repeated-root admissibility together with Dirichlet
boundary closure.

This does not place the IRFR model outside HJM. The resulting dynamics
should still admit an HJM representation after translation into standard
forward-rate coordinates. The distinction is that the HJM representation
does not reveal why the humped form arises.

In IRFR coordinates, the humped volatility profile is not an exogenous
input. It is a boundary-selected admissible mode.

\section{Why the Dirichlet case is
extreme}\label{why-the-dirichlet-case-is-extreme}

The Dirichlet boundary is analytically clean, but economically extreme.

A real central bank rarely fixes the short end in a perfectly
deterministic way over all relevant horizons. Even in strongly anchored
regimes, there may be residual uncertainty about future policy,
liquidity conditions, corridor implementation, and market
microstructure.

The Dirichlet case should therefore be interpreted as a limiting regime.
It identifies what happens when short-end stochastic loading is fully
suppressed.

This extreme case is valuable because it reveals the mechanism clearly.
Once the short-end loading is forced to vanish, the repeated-root branch
leaves the humped component as the surviving admissible mode.

More realistic regimes will generally be mixed. They will allow some
short-end variation while still anchoring the boundary. This is the role
of Robin closure, studied later.

\section{Boundary law and policy
transmission}\label{boundary-law-and-policy-transmission}

The Dirichlet case provides the sharpest illustration of policy as
boundary law.

The policy action is local: it pins the short-end boundary. The
consequence is non-local: the selected admissible mode determines how
stochastic sensitivity is distributed across the maturity domain.

The boundary condition does not directly prescribe the entire curve. It
selects a mode. The interior law then propagates that selection across
maturity.

policy enters locally at $\tau=0$;

the admissible field transmits the effect across $\tau>0$.

In the Dirichlet repeated-root case, this transmission takes the form of
a humped maturity loading. The short end is suppressed, intermediate
maturities carry the greatest sensitivity, and long maturities decay.

\section{Summary}\label{summary-4}

This chapter has studied Dirichlet boundary closure as the idealised
pinned short-end regime.

The repeated-root r=1 admissible loading is

\[
H(\tau)=(C_1+C_2\tau)e^{-a^2\tau}
\]

Dirichlet closure imposes

\[
H(0)=0
\]

Since H(0)=$C_1$, the boundary condition gives $C_1=0$. The selected
loading is therefore

\[
H(\tau)=C_2\tau e^{-a^2\tau}
\]

This loading is humped, with maximum at

\[
\tau=1/a^2
\]

The Dirichlet r=1 branch also admits a lower-boundary interpretation. In
shifted coordinates, the branch supports a shifted inverse-gamma
stationary density for fixed maturities $\tau>0$, and the lower
boundary is inaccessible maturity-by-maturity under the Feller
conditions. Full field-level lower-boundary preservation requires the
corresponding maturity-regularity assumption.

Dirichlet closure selects the humped repeated-root mode.

\section{Transition to Robin
closure}\label{transition-to-robin-closure}

The next chapter studies Robin boundary closure.

Robin closure relaxes the extreme assumption of a fully pinned short
end. Instead of fixing the boundary value outright, it allows a residual
boundary-active component while adding a boundary-induced repeated-root
component:

$H_R(\tau)=(C_0+C_1\tau)e^{-a_R^2\tau}$.

Economically, this represents partial anchoring. The short end is
influenced by policy but not perfectly fixed. In the repeated-root
branch, Robin closure therefore provides the more realistic bridge
between the Neumann/free benchmark and the idealised Dirichlet limit.

\chapter{Robin Boundary Closure and Partial Policy Anchoring}\label{chapter-9}

\section{Introduction}\label{introduction-6}

Chapter 8 studied Dirichlet boundary closure as the idealised
pinned-policy limit. In the repeated-root r=1 branch, the Dirichlet
condition suppresses short-end stochastic loading and selects the humped
loading

\[
H(\tau)\propto\tau e^{-a^2\tau}
\]

This is the cleanest example of boundary selection: the interior
equation creates the repeated-root mode space, while the boundary
condition selects the component compatible with a pinned short end.

The Dirichlet case is analytically revealing, but economically extreme.
A real policy regime rarely fixes the short end perfectly. Even when the
policy rate is strongly anchored, the market short end may still reflect
liquidity conditions, corridor mechanics, expectations of future policy,
balance-sheet constraints, and residual stochastic variation.

This chapter studies the intermediate case: Robin boundary closure.

Robin closure retains residual short-end stochastic loading while
allowing a boundary-induced repeated-root component. Algebraically, in
the repeated-root branch this gives a mixed loading of the form

$H_R(\tau)=(C_0+C_1\tau)e^{-a_R^2\tau}$.

Economically, Robin closure represents partial policy anchoring. The
short end is influenced by policy but not completely fixed.

The repeated-root r=1 branch is not studied here because empirical
humped volatility proves that this branch is the correct one. Rather, it
is studied because it is a mathematically distinguished branch of the
admissibility equation: the two exponential modes coalesce, producing
the linear-exponential component. The fact that this component resembles
humped volatility forms used in interest-rate modelling makes the branch
economically interpretable and worthy of detailed study. It does not, by
itself, establish that r=1 is empirically preferred.

The role of the Robin analysis is therefore conditional and precise. It
shows that, within the r=1 repeated-root mode space, partial boundary
anchoring selects a family of mixtures between a boundary-active
exponential component and a humped linear-exponential component.

\section{Robin closure as a mixed boundary
condition}\label{robin-closure-as-a-mixed-boundary-condition}

A Robin boundary condition imposes a linear relation between the value
of the loading and its derivative at the boundary. In homogeneous form,

$\alpha H(0)+\beta H'(0)=0$.

More generally,

$\alpha H(0)+\beta H'(0)=\gamma$.

Economically, the Robin case represents partial anchoring. The short end
is not fully pinned, but neither is it left entirely unconstrained by
the boundary law.

It is useful to distinguish two ideas. First, the Neumann/free benchmark
retains the ordinary exponential short-end stochastic mode,

$H_N(\tau)=A_Ne^{-a_N^2\tau}$.

Second, the Dirichlet repeated-root boundary suppresses short-end
loading and selects

$H_D^{(1)}(\tau)=C_D\tau e^{-a_D^2\tau}$.

Robin closure lies between these cases economically. It retains residual
short-end loading while also allowing a boundary-induced repeated-root
component. Algebraically, this gives

$H_R(\tau)=(C_0+C_1\tau)e^{-a_R^2\tau}$.

Thus Robin closure is best understood as the partially anchored
repeated-root regime. It does not simply equal the Neumann/free
benchmark, and it does not fully collapse to the Dirichlet branch unless
the residual short-end component is suppressed.

\section{Economic interpretation: partial
anchoring}\label{economic-interpretation-partial-anchoring}

The economic interpretation of Robin closure is partial policy
anchoring.

Under Neumann/free closure, the short end remains part of the stochastic
field. It is closed, but not pinned. Under Dirichlet closure, the short
end is pinned and short-end stochastic loading is suppressed.

Robin closure lies between these regimes. It allows the short end to be
influenced by policy while retaining residual stochastic movement.

This is closer to how many real monetary-policy regimes operate. A
central bank may strongly influence overnight or short-maturity rates,
but the market short end may still move because of liquidity premia,
expectations, collateral conditions, implementation frictions, and
uncertainty about future policy.

The Robin condition captures this by imposing not a fixed value, but a
relation between the boundary level and the boundary slope.
Economically, this says that the short end is anchored relative to the
way the curve leaves the boundary.

the boundary is constrained, but not fixed.

\section{Action on the repeated-root mode
space}\label{action-on-the-repeated-root-mode-space-1}

The repeated-root Robin loading is

$H_R(\tau)=(C_0+C_1\tau)e^{-a_R^2\tau}$.

The two components are

$C_0e^{-a_R^2\tau}$

and

$C_1\tau e^{-a_R^2\tau}$.

The first is boundary-active:

$H_R(0)=C_0$.

The second is boundary-induced and vanishes at the short end.

Differentiating,

$H_R'(\tau)=[C_1-a_R^2(C_0+C_1\tau)]e^{-a_R^2\tau}$,

so

$H_R'(0)=C_1-a_R^2C_0$.

A homogeneous Robin condition may be written as

$H_R'(0)=-\rho H_R(0)$,

where

$\rho=-H_R'(0)/H_R(0)$

whenever $H_R(0)\neq0$. Substituting the boundary value and derivative gives

\[
\rho = a_R^2 - C_1/C_0
\]

Equivalently,

\[
C_1/C_0 = a_R^2 - \rho
\]

Thus the Robin loading can be written as

$H_R(\tau)=C_0[1+(a_R^2-\rho)\tau]e^{-a_R^2\tau}$.

Writing k=$a_R^2$, this is

$H_R(\tau)=C_0[1+(k-\rho)\tau]e^{-k\tau}$.

The parameter $\rho$ measures the negative boundary slope of the loading
relative to its short-end level. It therefore controls the mixture of
boundary-active exponential loading and boundary-induced humped loading.

\section{Limiting cases and reference
regimes}\label{limiting-cases-and-reference-regimes}

The Robin repeated-root loading is

$H_R(\tau)=(C_0+C_1\tau)e^{-a_R^2\tau}$.

This form contains two components:

$e^{-a_R^2\tau}$, $\tau e^{-a_R^2\tau}$.

If $C_1=0$, the loading is purely exponential:

$H_R(\tau)=C_0e^{-a_R^2\tau}$.

This is the boundary-active profile associated with retained short-end
stochastic loading. It aligns with the Neumann/free benchmark.

If $C_0=0$, the loading becomes

$H_R(\tau)=C_1\tau e^{-a_R^2\tau}$.

This is the Dirichlet repeated-root humped loading.

Thus Robin closure should be read as a mixed repeated-root family. It
interpolates economically between retained short-end stochastic loading
and pinned-boundary suppression by varying the relative weights of $C_0$
and $C_1$.

This is an economic interpolation, not a claim that all three cases are
obtained as finite parameter values of one homogeneous Robin formula.

\begin{longtable}[]{@{}
  >{\raggedright\arraybackslash}p{(\columnwidth - 6\tabcolsep) * \real{0.2500}}
  >{\raggedright\arraybackslash}p{(\columnwidth - 6\tabcolsep) * \real{0.2500}}
  >{\raggedright\arraybackslash}p{(\columnwidth - 6\tabcolsep) * \real{0.2500}}
  >{\raggedright\arraybackslash}p{(\columnwidth - 6\tabcolsep) * \real{0.2500}}@{}}
\toprule\noalign{}
\begin{minipage}[b]{\linewidth}\raggedright
\textbf{Regime}
\end{minipage} & \begin{minipage}[b]{\linewidth}\raggedright
\textbf{Loading}
\end{minipage} & \begin{minipage}[b]{\linewidth}\raggedright
\textbf{Short-end loading}
\end{minipage} & \begin{minipage}[b]{\linewidth}\raggedright
\textbf{Interpretation}
\end{minipage} \\
\midrule\noalign{}
\endhead
\bottomrule\noalign{}
\endlastfoot
Neumann/free & $A_N e^{-a_N^2\tau}$ & retained & non-pinned
benchmark \\
Robin & $(C_0+C_1\tau)e^{-a_R^2\tau}$ & partially retained & partial
anchoring \\
Dirichlet & $C_D \tau e^{-a_D^2\tau}$ & suppressed & pinned boundary \\
\end{longtable}

\section{The anchoring parameter}\label{the-anchoring-parameter}

The Robin anchoring parameter is

$\rho=-H_R'(0)/H_R(0)$.

For

$H_R(\tau)=(C_0+C_1\tau)e^{-a_R^2\tau}$,

we have

$H_R(0)=C_0$, $H_R'(0)=C_1-a_R^2C_0$.

Therefore,

$\rho=a_R^2-C_1/C_0$.

Equivalently,

$C_1$=($a_R^2$-$\rho$)$C_0$.

This gives

$H_R(\tau)=C_0[1+(a_R^2-\rho)\tau]e^{-a_R^2\tau}$.

The parameter $\rho$ controls how the loading leaves the boundary. If
$\rho$=$a_R^2$, then $C_1=0$, and the loading is pure exponential. If
$\rho$\textless $a_R^2$, then $C_1$\textgreater0, and the loading contains a
positive boundary-induced linear-exponential component. If
$\rho$\textgreater $a_R^2$, then $C_1$\textless0, and the loading may eventually
change sign.

\section{Shape of the Robin-selected
loading}\label{shape-of-the-robin-selected-loading}

Write

\[
k=a_R^2, q=k-\rho
\]

Then

$H_R(\tau)=C_0(1+q\tau)e^{-k\tau}$.

The derivative is

$H_R'(\tau)=C_0[(q-k)-kq\tau]e^{-k\tau}$.

An interior stationary point occurs when

(q-k)-kq$\tau=0$.

If q$\neq$0, this gives

\[
\tau_*=(q-k)/(kq) = -\rho/{[}k(k-\rho){]}
\]

The shape depends on the anchoring parameter. If $\rho$=$a_R^2$, then q=0 and
the loading is pure exponential. If $\rho$\textless $a_R^2$, then
q\textgreater0 and the loading remains globally non-negative for
$C_0>0$. If $\rho$\textgreater $a_R^2$, then q\textless0 and the
loading eventually changes sign.

The Robin family can therefore produce boundary-active decay,
shoulder-like profiles, or partially humped profiles. Unlike the
Dirichlet branch, it need not vanish at the short end.

In the pure Dirichlet branch,

$H_D^{(1)}(\tau)\propto\tau e^{-k\tau}$,

and the maximum is necessarily at

\[
\tau=1/k
\]

In the Robin case, the curve may retain short-end loading and may or may
not display a pronounced interior hump, depending on the mixture
selected by the boundary condition.

\section{Positivity and admissibility under Robin
closure}\label{positivity-and-admissibility-under-robin-closure}

The Dirichlet r=1 branch has the cleanest lower-boundary interpretation.
Its selected loading

$H_D^{(1)}(\tau)=C_D\tau e^{-k\tau}$

is non-negative on the half-line for $C_D\geq0$ and vanishes at the
boundary. In the shifted positive-state formulation, Appendix B shows
that the Dirichlet branch admits a shifted inverse-gamma stationary
density for each fixed $\tau>0$ and that the lower boundary is
inaccessible maturity-by-maturity under the Feller conditions.

Robin closure is more delicate because it retains residual short-end
loading. The selected loading

$H_R(\tau)=C_0(1+q\tau)e^{-k\tau}$

is non-negative on the half-line if the prefactor and the linear term
preserve sign. If $C_0>0$, the necessary and sufficient
condition for non-negativity on the entire half-line is

$q\geq0$.

Equivalently,

$k-\rho\geq0$, $\rho\leq k$.

If $\rho>k$, the linear term eventually becomes negative, and the
loading changes sign at sufficiently large maturity. Such a branch may
still have a mathematical interpretation, but it is not compatible with
a simple non-negative lower-boundary interpretation on the whole
half-line.

When the Robin loading is sign-compatible and the shifted positive-state
KFE structure is retained, Robin inherits the same maturity-by-maturity
inverse-gamma and Feller logic as the Dirichlet case. A separate density
derivation is not required. The additional Robin requirement is the
global sign condition on the loading.

Thus Robin admissibility has two layers: the local shifted-state
diffusion must remain in the positive-state Feller class, and the
maturity loading must remain sign-compatible on the domain under
consideration. Full field-level lower-boundary preservation again
requires the corresponding maturity-regularity assumption.

\section{Robin closure and monetary-policy
regimes}\label{robin-closure-and-monetary-policy-regimes}

Robin closure gives a natural language for partially anchored
monetary-policy regimes.

In a fully pinned regime, the short end is treated as fixed. This is the
Dirichlet idealisation. In a free-boundary benchmark, the short end
remains stochastic. This is Neumann/free closure. Between these extremes
lies the more realistic case in which policy anchors the short end
imperfectly.

The Robin condition captures this intermediate regime. The boundary
value $H_R(0)=C_0$ represents the amount of short-end loading retained.
The derivative $H_R'(0)$ represents how the maturity loading leaves the
boundary. The Robin relation constrains these two quantities relative to
one another.

Economically, this means that policy does not eliminate short-end
stochastic variation entirely. Instead, it controls how short-end
variation is transmitted into the rest of the curve.

The anchoring parameter $\rho$ therefore acts as a transmission coefficient.
It determines how much short-end stochastic variation enters the
interior through the admissible mode mixture.

This is exactly the role of a partially anchored policy regime. Policy
influences the boundary, but the market still supplies residual
stochastic movement and slope at the short end.

\section{Comparison with natural and Dirichlet
regimes}\label{comparison-with-natural-and-dirichlet-regimes}

The three canonical closures now have clearer interpretations.

Neumann/free closure retains the ordinary exponential short-end
stochastic mode:

$H_N(\tau)=A_Ne^{-a_N^2\tau}$.

Dirichlet closure suppresses short-end loading and selects the
repeated-root humped mode:

$H_D^{(1)}(\tau)=C_D\tau e^{-a_D^2\tau}$.

Robin closure lies between these regimes by retaining a residual
short-end component while adding a boundary-induced repeated-root
component:

$H_R(\tau)=(C_0+C_1\tau)e^{-a_R^2\tau}$.

Neumann/free $\to$ short-end stochastic loading retained;

\begin{center}\emph{Dirichlet \(\to\) short-end stochastic loading suppressed;}\end{center}

Robin $\to$ short-end stochastic loading partially retained.

The same boundary-field logic therefore produces different observable
maturity loadings depending on how the short end is closed.

\section{Relation to humped
volatility}\label{relation-to-humped-volatility}

Robin closure modifies the interpretation of humped volatility.

In the Dirichlet case, the humped loading is selected cleanly:

$H_D^{(1)}(\tau)\propto\tau e^{-k\tau}$.

In the Robin case, the selected loading is generally mixed:

$H_R(\tau)=(C_0+C_1\tau)e^{-k\tau}$.

This may display hump-like behaviour, depending on the anchoring
parameter and coefficient signs. But the hump is no longer forced to
vanish at the short end. Instead, the profile may include a residual
short-end component.

This is economically plausible. In partially anchored regimes, short
maturities may retain some volatility, while intermediate maturities may
still display heightened sensitivity. Robin closure therefore provides a
structural route to shifted or partially humped volatility profiles.

This observation should be read carefully. The resemblance between
Robin-selected profiles and empirical humped-volatility specifications
motivates the study of the r=1 branch, but it does not validate that
branch by itself. Empirical relevance remains a calibration and testing
question.

Dirichlet: pure pinned-boundary hump;

Robin: mixed anchored-boundary profile.

\section{Boundary law as a regime
parameter}\label{boundary-law-as-a-regime-parameter}

The Robin boundary condition suggests a broader modelling
interpretation.

The boundary law itself can be regarded as a regime parameter. Changing
the boundary coefficients, or equivalently changing the anchoring ratio
$\rho$, changes the realised maturity loading without changing the interior
admissible equation.

This is significant. It means that policy regimes can be modelled as
changes in boundary law rather than changes in the interior dynamics.

The interior law describes the admissible propagation of stochastic
shocks through maturity. The boundary law determines how the short end
is closed and therefore which admissible profile is realised.

In empirical implementation, $\rho$ becomes an estimable parameter linking
observed maturity loadings to implicit policy anchoring strength.
Different monetary-policy environments may therefore correspond not to
different interior dynamics, but to different boundary-law selections
within the same admissible branch.

observed maturity profile = interior admissible structure + boundary-law
selection via $\rho$.

Thus changes in observed term-structure shape need not reflect changes
in the underlying interior dynamics. They may instead reflect changes in
the boundary regime.

\section{Summary}\label{summary-5}

This chapter has studied Robin boundary closure as the partially
anchored regime of the IRFR field.

The repeated-root Robin loading is

$H_R(\tau)=(C_0+C_1\tau)e^{-a_R^2\tau}$.

The boundary value and slope are

$H_R(0)=C_0$, $H_R'(0)=C_1-a_R^2C_0$.

The anchoring parameter satisfies

\[
\rho = a_R^2 - C_1/C_0,
\]

or equivalently

$H_R(\tau)=C_0[1+(a_R^2-\rho)\tau]e^{-a_R^2\tau}$.

This loading retains short-end stochastic loading unless $C_0=0$. It
therefore represents partial policy anchoring.

For $\rho\leq a_R^2$, the selected loading can preserve a non-negative
lower-boundary interpretation on the half-line, assuming
$C_0>0$. Under the shifted positive-state KFE structure, the
Robin case then inherits the fixed-maturity inverse-gamma and Feller
logic described for Dirichlet in Appendix B. For $\rho$\textgreater $a_R^2$,
the loading changes sign at sufficiently long maturity and requires a
different interpretation.

Robin closure selects a mixed repeated-root loading.

This mixed loading lies between the Neumann/free benchmark and the
Dirichlet pinned-boundary limit. It provides a structural representation
of policy regimes in which the short end is influenced by policy but not
perfectly fixed.

\section{Transition}\label{transition}

The next step is to connect the boundary-law taxonomy to observable
term-structure behaviour.

The Neumann/free, Dirichlet, and Robin cases show that the same broad
boundary-field logic can generate different maturity loadings depending
on how the short end is closed. The next part of the monograph relates
these boundary-selected loadings to observable volatility profiles,
Ritchken--Chuang-type specifications, and empirical policy regimes.

The purpose is no longer only to derive admissible modes, but to ask how
those modes can be identified, calibrated, and tested in market data.

\chapter{Boundary-Selected Loadings and Observable Volatility}\label{chapter-10}

\section{Introduction}\label{introduction-7}

The previous chapters developed the boundary-law structure of the IRFR
field.

Chapter 7 studied Neumann/free boundary closure as the non-pinned
benchmark. Chapter 8 studied Dirichlet closure as the idealised
pinned-policy limit. Chapter 9 studied Robin closure as the partially
anchored regime.

Together, these chapters established the central boundary-law principle:

\begin{center}\emph{the interior law determines the admissible mode space;}\end{center}
\begin{center}\emph{the boundary law selects the realised maturity loading.}\end{center}
The present chapter translates that principle into observable
term-structure behaviour. The object of interest is the maturity loading
$H(\tau)$, which determines how the common stochastic shock is expressed
across the curve. Under the one-factor IRFR convention,

\[
g(\tau) = aH(\tau),
\]

and the local stochastic evolution is written as

\[
dx(t,\tau)=-f(x,t,\tau)dt+\sqrt{2} g(x,t,\tau)dW_t
\]

When $g(\tau)$ = a$H(\tau)$, the instantaneous covariance between maturity points
is

\[
\operatorname{Cov}(dx(t,\tau_1),dx(t,\tau_2))=2a^2H(\tau_1)H(\tau_2)dt
\]

Thus $H(\tau)$ is the deterministic maturity profile through which the one
stochastic factor is transmitted across the curve.

The purpose of this chapter is to interpret the boundary-selected forms
of $H(\tau)$ as observable volatility loadings.

\begin{center}\emph{observable volatility shape is not imposed directly;}\end{center}
\emph{it is selected by boundary closure from the admissible mode
space.}

Equivalently, the volatility loading is the observable footprint of the
boundary law.

\section{From diffusion scale to volatility
loading}\label{from-diffusion-scale-to-volatility-loading}

In the IRFR convention, the diffusion-scale coefficient is g, while the
conventional SDE volatility coefficient is $\sqrt{}$2g.

For the one-factor maturity-loading representation,

\[
g(\tau)=aH(\tau)
\]

Therefore the conventional instantaneous volatility loading is

\[
\sqrt{2}aH(\tau)
\]

The economically relevant maturity shape is $H(\tau)$. The constants a and $\sqrt{}$2
set scale and normalisation. The function $H(\tau)$ determines where along
the curve stochastic sensitivity is concentrated.

This distinction matters. A one-factor model has only one Brownian
source, but it can still generate non-trivial volatility shapes because
the deterministic maturity loading can vary across $\tau$.

Thus stochastic dimension and observed maturity profile are different
objects. The covariance-rate kernel makes this explicit:

\[
K(\tau_1,\tau_2) = (1/dt) \operatorname{Cov}(dx(t,\tau_1), dx(t,\tau_2)) = 2a^2 H(\tau_1)H(\tau_2)
\]

The kernel is rank one, but its maturity dependence is shaped by H. The
curve can therefore display structured maturity sensitivity even in a
one-factor stochastic model.

In empirical covariance terms, the one-factor IRFR kernel has a single
non-zero volatility direction. The corresponding maturity eigenfunction
is proportional to $H(\tau)$. Thus, within the one-factor restriction, the
boundary-selected loading may be compared directly with the leading
empirical volatility mode of forward-rate changes.

\section{Boundary-selected loading
families}\label{boundary-selected-loading-families}

The boundary chapters produced three canonical loading families. They
should not be presented as three finite parameter settings of a single
homogeneous Robin formula. Rather, they are the observable loading forms
associated with the three boundary regimes developed in Chapters 7-9.

The Neumann/free benchmark retains the ordinary exponential short-end
stochastic loading:

\[
H_N(\tau) = A_N \exp(-a_N^2\tau)
\]

The short-end loading is retained:

\[
H_N(0) = A_N
\]

Dirichlet closure suppresses short-end stochastic loading. In the
repeated-root branch it selects the boundary-induced humped loading:

\[
H_D^{(1)}(\tau) = C_D \tau \exp(-a_D^2\tau)
\]

The short-end loading is suppressed:

\[
H_D^{(1)}(0) = 0
\]

Robin closure represents partial anchoring. It retains a residual
short-end component while allowing a boundary-induced repeated-root
component:

\[
H_R(\tau) = (C_0 + C_1\tau) \exp(-a_R^2\tau)
\]

When $C_0\neq0$, this can also be written in anchoring-ratio form as

\[
H_R(\tau) = C_0{[}1 + (a_R^2 - \rho)\tau{]} \exp(-a_R^2\tau),
\]

where

\[
\rho = -H_R'(0) / H_R(0)
\]

The three regimes may therefore be summarised as follows.

\begin{longtable}[]{@{}
  >{\raggedright\arraybackslash}p{(\columnwidth - 6\tabcolsep) * \real{0.2500}}
  >{\raggedright\arraybackslash}p{(\columnwidth - 6\tabcolsep) * \real{0.2500}}
  >{\raggedright\arraybackslash}p{(\columnwidth - 6\tabcolsep) * \real{0.2500}}
  >{\raggedright\arraybackslash}p{(\columnwidth - 6\tabcolsep) * \real{0.2500}}@{}}
\toprule\noalign{}
\begin{minipage}[b]{\linewidth}\raggedright
\textbf{Boundary regime}
\end{minipage} & \begin{minipage}[b]{\linewidth}\raggedright
\textbf{Loading}
\end{minipage} & \begin{minipage}[b]{\linewidth}\raggedright
\textbf{Short-end loading}
\end{minipage} & \begin{minipage}[b]{\linewidth}\raggedright
\textbf{Interpretation}
\end{minipage} \\
\midrule\noalign{}
\endhead
\bottomrule\noalign{}
\endlastfoot
Neumann/free & $A_N\exp(-a_N^2\tau)$ & retained & non-pinned benchmark \\
Dirichlet & $C_D\tau\exp(-a_D^2\tau)$ & suppressed & pinned boundary \\
Robin & $(C_0+C_1\tau)\exp(-a_R^2\tau)$ & partially retained & partial
anchoring \\
\end{longtable}

These forms are not unrelated volatility parametrisations. They are
boundary-selected realisations of the IRFR maturity-domain structure.

\section{The Neumann/free exponential
benchmark}\label{the-naturalfree-exponential-benchmark}

Neumann/free closure provides the non-pinned benchmark. It retains the
boundary-active exponential mode

\[
H_N(\tau) = A_N \exp(-a_N^2\tau)
\]

This loading is non-zero at the short end whenever $A_N\neq0$. Thus
short-end stochastic exposure is retained.

The associated instantaneous variance rate is

\[
\partial_t  \operatorname{Var}(x^N_{t,t+\tau}) = 2a_N^2 A_N^2 \exp(-2a_N^2\tau)
\]

This profile is monotone in maturity. Its role is not to generate the
pinned-boundary hump. Its role is to provide the free benchmark against
which policy anchoring is measured.

Economically, the Neumann/free benchmark means that no additional
policy-pinning condition is imposed at $\tau=0$. It is not a claim that
monetary policy is absent institutionally. It means that the boundary
law does not suppress the short-end stochastic loading.

\section{The Dirichlet humped
loading}\label{the-dirichlet-humped-loading}

The Dirichlet case gives the cleanest humped volatility profile.

The selected repeated-root loading is

\[
H_D^{(1)}(\tau) = C_D \tau \exp(-a_D^2\tau)
\]

Ignoring the scale constant, define

\[
h_D(\tau) = \tau \exp(-a_D^2\tau)
\]

Then

\[
h_D(0) = 0, h_D(\tau) \to 0 as \tau \to \infty
\]

The derivative is

\[
h_D'(\tau) = \exp(-a_D^2\tau)(1 - a_D^2\tau)
\]

Thus

\[
h_D'(\tau) = 0 \Leftrightarrow \tau = 1/a_D^2
\]

The maximum occurs at

\[
\tau_{\max} = 1/a_D^2
\]

This is the standard qualitative humped shape: low short-end volatility,
increasing sensitivity at intermediate maturities, and decay at long
maturities.

The interpretive distinction is essential. In the IRFR framework, this
shape is not specified as an empirical convenience. It is the
boundary-selected survivor of the repeated-root mode space under a
pinned short-end condition.

This is why the Dirichlet branch provides the cleanest structural route
to humped volatility specifications.

\section{The Robin mixed loading}\label{the-robin-mixed-loading}

The Robin case generalises the Dirichlet profile by retaining some
boundary-active component.

The Robin-selected repeated-root loading is

\[
H_R(\tau) = (C_0 + C_1\tau) \exp(-a_R^2\tau)
\]

The two components are

$C_0\exp(-a_R^2\tau)$

and

$C_1\tau\exp(-a_R^2\tau)$.

The first is boundary-active. The second is boundary-induced and
vanishes at the short end.

At the boundary,

\[
H_R(0) = C_0
\]

Differentiating gives

\[
H_R'(0) = C_1 - a_R^2 C_0
\]

Thus, whenever $C_0\neq0$,

\[
\rho = -H_R'(0)/H_R(0) = a_R^2 - C_1/C_0
\]

Equivalently,

\[
C_1/C_0 = a_R^2 - \rho
\]

Therefore the Robin loading can be written as

\[
H_R(\tau) = C_0{[}1 + (a_R^2 - \rho)\tau{]} \exp(-a_R^2\tau)
\]

Writing k=$a_R^2$, this becomes

$H_R(\tau)=C_0[1+(k-\rho)\tau]e^{-k\tau}$.

This formula shows exactly how Robin closure mixes residual short-end
stochastic loading with the boundary-induced linear-exponential
component.

\section{Shape diagnostics of the Robin
family}\label{shape-diagnostics-of-the-robin-family}

The Robin family has useful short-end diagnostics. Write

\[
k = a_R^2, q = k - \rho
\]

Then

\[
H_R(\tau) = C_0(1 + q\tau) \exp(-k\tau)
\]

The boundary value is

$H_R(0)=C_0$,

and the boundary slope is

$H_R'(0)=-C_0\rho$.

Thus, whenever $H_R(0)\neq0$,

\[
\rho = -H_R'(0) / H_R(0)
\]

The Robin parameter therefore has a direct observable interpretation: it
is the negative boundary slope of the volatility loading relative to its
short-end level. Low $\rho$ corresponds to weak anchoring. Larger positive $\rho$
corresponds to stronger suppression of the short-end direction.

The same family also gives a simple turning-point diagnostic. Since

\[
H_R'(\tau) = C_0{[}(q - k) - kq\tau{]} \exp(-k\tau),
\]

an interior critical point satisfies

\[
(q - k) - kq\tau = 0
\]

If q$\neq$0, then

\[
\tau_* = (q - k)/(kq) = -\rho/{[}k(k - \rho){]}
\]

The profile has an interior turning point only when this expression is
positive and lies in the relevant maturity domain. Thus $\rho$ controls not
only the short-end value and boundary slope, but also whether the
profile is monotone, shoulder-like, or humped.

This diagnostic is useful for empirical interpretation because it links
the estimated boundary behaviour of the volatility loading to the shape
of the full maturity profile.

\section{Sign compatibility and lower-boundary
interpretation}\label{sign-compatibility-and-lower-boundary-interpretation}

The Dirichlet repeated-root loading

\[
H_D^{(1)}(\tau) = C_D \tau \exp(-a_D^2\tau)
\]

is non-negative on the half-line when $C_D\geq0$. It also vanishes at the
short end.

For Robin closure, sign compatibility is more delicate. In $\rho$-form,

\[
H_R(\tau) = C_0{[}1 + (a_R^2 - \rho)\tau{]} \exp(-a_R^2\tau)
\]

For $C_0$\textgreater0, non-negativity on the full half-line requires

\[
\rho \le a_R^2
\]

If $\rho$\textgreater $a_R^2$, the linear factor eventually changes sign. Such
a branch may remain mathematically meaningful, but it no longer supports
a simple globally non-negative lower-boundary interpretation on the full
half-line.

This sign condition is a loading-level admissibility condition. The
stronger state-space statement is developed in Appendix B and summarised
in Chapter 12: under the shifted positive-state KFE structure, the
Dirichlet branch admits a shifted inverse-gamma stationary density for
each fixed $\tau>0$, and the lower boundary is inaccessible
maturity-by-maturity under the Feller conditions. Robin branches inherit
the same local lower-boundary logic when their loading is
sign-compatible.

\section{Comparison with Ritchken--Chuang-type
forms}\label{comparison-with-ritchken-chuang-type-forms}

Humped volatility functions of the form

\[
\tau \exp(-k\tau)
\]

or related linear-exponential variants are familiar in interest-rate
modelling. Ritchken--Chuang-type specifications use such maturity
profiles because they capture the empirical observation that volatility
is often low at the very short end, higher at intermediate maturities,
and lower again at long maturities.

The IRFR interpretation is different in emphasis.

The resemblance to these specifications does not prove that the r=1
branch is the empirically correct branch. Nor does it show that the IRFR
model is uniquely determined by observed humped volatility.

Rather, the point is structural:

\emph{a familiar empirical volatility shape appears as an admissible
boundary-selected mode.}

In conventional modelling, the humped shape is often introduced as a
volatility input. In the IRFR framework, the same qualitative shape
appears from repeated-root admissibility together with boundary closure.

This provides a structural explanation for a class of volatility shapes
already known to be useful. It does not replace empirical calibration or
testing. It motivates them.

The appropriate claim is therefore modest but meaningful: IRFR gives a
structural route to familiar humped volatility forms.

\section{Quadratic models and boundary-selected
loadings}\label{quadratic-models-and-boundary-selected-loadings}

Quadratic term-structure models provide an important benchmark for
nonlinear interest-rate modelling. Their strategy is to enrich the state
representation: yields and bond prices are allowed to depend
quadratically on latent factors, and the resulting pricing equations
inherit the corresponding polynomial or Riccati-type structure.

The IRFR framework follows a different logic. It does not first
introduce quadratic state dependence. Instead, it starts from the
rolling-maturity field and asks which maturity loadings are admissible
under the interior law and which of those loadings are selected by the
boundary condition.

The contrast can be stated simply:

\begin{center}\emph{quadratic models generate shape through nonlinear state mapping;}\end{center}
\emph{IRFR generates shape through admissible modes and boundary
selection.}

This distinction matters for the interpretation of humped volatility. In
a quadratic or polynomial framework, curvature is produced by the chosen
state mapping. In the IRFR framework, the humped component

\[
\tau \exp(-a^2\tau)
\]

appears as the coalescing-root mode of the maturity-domain admissibility
problem and is then selected by the boundary law. The hump is therefore
not inserted as a polynomial factor or an exogenous volatility curve. It
is the observable survivor of the interior mode space after boundary
closure has acted.

This does not make the IRFR framework more general than quadratic
term-structure models. It makes a different claim. Before adding
polynomial state dependence, one should first account for the geometry
of rolling maturity and the selection role of the short-end boundary.

\section{Representation in HJM
coordinates}\label{representation-in-hjm-coordinates}

The boundary-selected IRFR loadings can be translated into conventional
forward-rate coordinates through

\[
x(t,\tau) = F(t,t+\tau)
\]

Equivalently,

\[
F(t,T) = x(t,T-t)
\]

A maturity loading expressed in $\tau$-coordinates becomes a loading in
time-to-maturity coordinates in the HJM representation.

For example, the Dirichlet r=1 loading

\[
H(\tau) \propto \tau \exp(-a^2\tau)
\]

corresponds to

\[
H(T-t) \propto (T-t) exp{[}-a^2(T-t){]}
\]

From an HJM perspective, this may appear as an externally specified
time-to-maturity volatility function. The HJM drift restriction then
gives the corresponding arbitrage-free drift.

From the IRFR perspective, however, the same function has a different
origin. It is not first specified as a volatility input. It is selected
by the rolling-field boundary law.

This is the representation-versus-explanation distinction introduced
earlier. HJM provides the ambient no-arbitrage representation. IRFR
explains how rolling-maturity geometry and boundary closure may select
particular admissible loadings.

\section{Empirical interpretation of boundary
regimes}\label{empirical-interpretation-of-boundary-regimes}

The boundary-law framework gives an interpretation of observed
volatility profiles in terms of regimes.

A profile with substantial short-end loading is naturally associated
with a free or weakly anchored boundary. This corresponds to
Neumann/free closure or low-anchoring Robin closure.

A profile with suppressed short-end loading and an interior maximum is
naturally associated with a pinned or strongly anchored boundary. The
limiting case is Dirichlet closure.

A profile with residual short-end loading and intermediate-maturity
curvature is naturally associated with partial anchoring. This
corresponds to Robin closure.

Thus empirical volatility shape may contain information about the
short-end boundary regime.

This does not mean that observed volatility alone identifies the
boundary law uniquely. Different model specifications may produce
similar shapes. Multiple factors, stochastic volatility, measurement
noise, and market segmentation can all affect observed term-structure
volatility.

The claim is more limited:

\emph{boundary-selected IRFR modes provide a structural interpretation
of volatility shape.}

They suggest a way to connect maturity profiles to policy anchoring
without treating the volatility function as purely exogenous.

\section{The role of the r=1
branch}\label{the-role-of-the-r1-branch}

The repeated-root r=1 branch plays a special role because it produces
the linear-exponential component

\[
\tau \exp(-k\tau)
\]

This component is mathematically distinguished: it appears when two
admissible exponential modes coalesce. It is also economically
interpretable because it gives a natural humped maturity profile.

However, the branch should not be overinterpreted. The empirical
usefulness of humped volatility forms does not prove that r=1 is the
correct branch. It only shows that the repeated-root branch is a
meaningful case to study.

The proper interpretation is conditional:

\emph{if the admissible mode structure is near the repeated-root
branch,}

\emph{then boundary selection gives a structural route to humped or
partially humped volatility profiles.}

Whether observed markets are best described by this branch is an
empirical question. It must be tested through calibration, stability,
and comparison with alternative branches.

The theoretical contribution is to identify why the r=1 branch is
structurally special and how its boundary-selected loadings relate to
familiar empirical forms.

\section{Summary}\label{summary-6}

This chapter has translated the boundary-selected loading structure into
observable volatility terms.

The key object is the maturity loading $H(\tau)$. In the one-factor IRFR
model,

\[
g(\tau) = aH(\tau),
\]

and the instantaneous covariance-rate kernel is

\[
K(\tau_1,\tau_2) = 2a^2H(\tau_1)H(\tau_2)
\]

Thus $H(\tau)$ determines the maturity shape of volatility.

The Neumann/free benchmark retains the ordinary exponential short-end
stochastic loading:

\[
H_N(\tau) = A_N \exp(-a_N^2\tau)
\]

The Dirichlet repeated-root branch suppresses short-end loading and
selects the humped profile:

\[
H_D^{(1)}(\tau) = C_D \tau \exp(-a_D^2\tau)
\]

The Robin repeated-root branch gives the partially anchored mixed
profile:

\[
H_R(\tau) = (C_0 + C_1\tau) \exp(-a_R^2\tau)
\]

These are not arbitrary volatility functions. They are boundary-selected
maturity loadings. Their empirical relevance must be tested, but their
structural origin is clear.

\emph{observable volatility shape is the footprint of boundary
selection.}

\chapter{Calibration, Identification, and Empirical Discipline}\label{chapter-11}

\section{Introduction}\label{introduction-8}

The preceding chapters derived the IRFR framework as a boundary-value
theory of interest-rate dynamics.

The early chapters introduced the instantaneous rolling forward rate as
a maturity-domain field. The middle chapters derived the affine drift
operator, admissible maturity modes, and the r-branch. The boundary
chapters showed how Neumann/free, Dirichlet, and Robin closures select
different realised maturity loadings. Chapter 10 then translated those
boundary-selected loadings into observable volatility profiles.

The present chapter turns from structural derivation to empirical
discipline.

The purpose is not to claim that the IRFR framework has already been
empirically validated. It has not. The purpose is to state how the
framework should be taken to data, what must be estimated, what can and
cannot be inferred from volatility shape, and how boundary regimes might
be tested.

Estimation and identification are distinct tasks. Estimation concerns
recovering the empirical maturity loading $H(\tau)$, and possibly its scale,
from data. Identification concerns mapping features of that loading ---
its decay, short-end value, short-end slope, hump location, and factor
structure --- to structural parameters such as k, $\rho$, and the boundary
regime.

There is also a preliminary data-construction step. The IRFR field is
not observed directly. Market data are typically available as bond
prices, swap rates, futures, deposits, OIS curves, or option-implied
quantities. Recovering $x(t,\tau)$ and estimating changes in $H(\tau)$ therefore
requires curve construction, interpolation, smoothing, and a choice of
instruments and discounting conventions. These choices are not
innocuous. They affect the estimated maturity loading and, consequently,
any inferred boundary parameter. Empirical IRFR calibration should
therefore treat curve construction as part of the modelling procedure
rather than as a neutral preprocessing step.

The guiding principle is:

empirical fit is evidence, not proof of canonicality.

A humped volatility profile may be consistent with the repeated-root r =
1 branch, but it does not prove that the repeated-root branch is the
true branch. A good calibration may support the usefulness of the
framework, but it does not establish uniqueness. A boundary-selected
loading may explain a volatility shape, but the same shape may be
approximated by other models.

This chapter therefore sets out ten principles for calibration,
identification, and empirical testing. These principles are not
additional axioms. They are empirical safeguards. They are intended to
prevent the structural claims of the monograph from being overstated
when the model is confronted with data.

The ten principles are:

\begin{enumerate}
\item Estimate the maturity loading, not only the volatility level.
\item Separate stochastic dimension from maturity geometry.
\item Identify the decay scale.
\item Estimate boundary anchoring from short-end behaviour.
\item Do not infer $r=1$ from hump resemblance alone.
\item Distinguish scale from shape.
\item Use principal components carefully.
\item Treat boundary regimes as estimable hypotheses.
\item Check positivity and admissibility separately.
\item Treat empirical fit as evidence, not proof.
\end{enumerate}

Together, these principles define the empirical programme of the IRFR
framework.

\section{Principle 1: estimate the loading, not only the volatility
level}\label{principle-1-estimate-the-loading-not-only-the-volatility-level}

The primary empirical object in the one-factor IRFR model is the
maturity loading

\[
H(\tau)
\]

Under the convention used in the previous chapters,

\[
g(\tau)=aH(\tau),
\]

and the conventional volatility coefficient is

\[
\sqrt{2} aH(\tau)
\]

The overall volatility level matters, but it is not the main structural
object. The structural object is the shape of $H(\tau)$: where the curve is
sensitive to the common shock, how that sensitivity changes across
maturity, and how the loading behaves at the short end.

A calibration that matches only the total variance level is therefore
insufficient. It may reproduce the average magnitude of rate changes
while missing the maturity-domain geometry.

The IRFR framework makes a stronger claim. It says that the shape of the
maturity loading contains information about the boundary law. A loading
with substantial short-end value suggests a different boundary regime
from a loading that vanishes at the short end. A loading with a humped
component suggests a different admissible-mode structure from a purely
monotone exponential decay.

Thus calibration must estimate not only how much volatility there is,
but where that volatility lives along the maturity domain.

The empirical question is therefore:

what is the estimated shape of $H(\tau)$?

Only after this shape is estimated can one ask which IRFR boundary
regime it resembles.

\section{Principle 2: separate stochastic dimension from maturity
geometry}\label{principle-2-separate-stochastic-dimension-from-maturity-geometry}

A one-factor stochastic model does not imply a flat or trivial maturity
profile.

In the one-factor IRFR model, all maturities share a common Brownian
shock,

$dW_t$.

The covariance kernel has the rank-one form

\[
K(\tau_1,\tau_2)=2a^2H(\tau_1)H(\tau_2)
\]

The stochastic dimension is one. But the deterministic maturity geometry
is encoded in $H(\tau)$, and this function may be non-trivial.

This distinction is essential. A model may be one-factor in randomness
while still having structured maturity sensitivity. The curve may react
to one shock, but the loading of that shock across maturities can be
humped, shouldered, monotone, or boundary-active.

Empirically, this means that the number of stochastic factors and the
shape of the volatility loading should not be conflated.

A one-factor PCA approximation, for example, may identify a dominant
empirical volatility direction. That direction is not merely a scalar
factor. It is a function over maturity. In the one-factor IRFR setting,
that function is the empirical counterpart of $H(\tau)$.

Under the one-factor IRFR restriction, the leading empirical
eigenfunction must coincide, up to scale and estimation error, with the
maturity loading $H(\tau)$. This is a testable restriction, not merely an
interpretation.

Thus the empirical task is not only to ask how many factors are needed.
It is also to ask whether the leading maturity loading has a shape
compatible with an IRFR boundary-selected mode.

\section{Principle 3: identify the decay
scale}\label{principle-3-identify-the-decay-scale}

The admissible IRFR loadings contain a maturity-decay scale.

In the repeated-root branch, the generic loading has the form

\[
H(\tau)=(C_1+C_2\tau)e^{-k\tau}
\]

In the $r=1$ normalisation used in the previous chapters,

\[
k=a^2
\]

This relation is part of the structural normalisation of the
repeated-root branch. In empirical calibration, one may impose this
relation or treat the decay parameter and volatility scale separately,
depending on the chosen normalisation and estimation strategy.

The Dirichlet $r=1$ loading is

$H_D(\tau)=C_D\tau e^{-a^2\tau}$.

Ignoring scale, define

$h_D(\tau)=\tau e^{-a^2\tau}$.

Then

$h_D'(\tau)=e^{-a^2\tau}(1-a^2\tau)$.

The maximum occurs when

$1-a^2\tau=0$,

so

\[
\tau_{\max}=1/a^2
\]

Equivalently,

$a^2$=1/$\tau_{\max}$.

This gives a simple empirical diagnostic. If an estimated volatility
loading has a clear interior maximum, the location of that maximum gives
an estimate of the maturity decay scale under the Dirichlet
repeated-root interpretation.

The same idea can be applied more generally. In Robin closure, the
location and even the existence of an interior turning point depend on
the anchoring parameter. But the decay scale k remains a central object.
It controls the long-maturity decay of the loading.

Thus calibration should identify the decay scale explicitly rather than
burying it inside an arbitrary volatility parametrisation.

\section{Principle 4: estimate boundary anchoring from short-end
behaviour}\label{principle-4-estimate-boundary-anchoring-from-short-end-behaviour}

Robin closure introduces a boundary anchoring parameter.

The homogeneous Robin condition is

$\alpha H(0)+\beta H'(0)=0$.

When $\beta\neq0$, define $\rho=\alpha/\beta$.

Then the Robin condition may be written as

$H'(0)=-\rho H(0)$.

Thus, whenever $H(0)\neq0$,

$\rho=-H'(0)/H(0)$.

This is a direct empirical interpretation. The anchoring parameter is
the negative boundary slope of the volatility loading relative to its
short-end level.

A low value of $\rho$ corresponds to weak anchoring: the loading leaves the
boundary relatively freely. A high value of $\rho$ corresponds to stronger
suppression of the short-end direction.

For the repeated-root Robin loading,

$H_R(\tau)=(C_0+C_1\tau)e^{-a_R^2\tau}$,

we have

$H_R(0)=C_0$,

and

$H_R'(0)=C_1-a_R^2C_0$.

Thus, whenever $C_0\neq0$,

$\rho=a_R^2-C_1/C_0$.

Equivalently,

$C_1/C_0=a_R^2-\rho$.

The boundary parameter can therefore be inferred from the estimated
short-end level and slope of the loading, provided these quantities are
observable or can be reliably extrapolated.

In practice, however, both H(0) and $H'(0)$ are extrapolated objects. The
instantaneous short end is rarely observed directly as a clean maturity
point, and the slope at the boundary is sensitive to curve construction,
smoothing, interpolation, and the chosen parametric form. Estimates of $\rho$
should therefore be treated as structural diagnostics rather than
directly observed quantities.

This is one of the strongest empirical implications of the Robin
formulation, but it is also one of the most delicate. Boundary anchoring
is not merely a qualitative policy label. It becomes a parameter that
can, in principle, be estimated from the shape of the maturity loading
near the short end.

\section{Principle 5: do not infer r = 1 from hump resemblance
alone}\label{principle-5-do-not-infer-r-1-from-hump-resemblance-alone}

The repeated-root branch is structurally important because it generates
the linear-exponential component

\[
\tau e^{-k\tau}
\]

This component naturally produces a humped maturity profile. That
profile resembles familiar humped volatility specifications used in
interest-rate modelling.

This resemblance is useful, but it is not proof.

A hump identifies at most a functional feature, not a root structure.
Many models can generate humped or hump-like maturity profiles. Only
within the IRFR admissibility calculation does the linear-exponential
form arise from a repeated root of the maturity-domain operator.

A humped empirical volatility profile may therefore be consistent with
the r = 1 repeated-root branch. It may motivate testing that branch. It
may make the branch economically interpretable. But it does not
establish that r = 1 is the correct branch.

Distinct-root branches must also be considered. They generate sums of
exponentials without the linear-exponential repeated-root term,
typically producing monotone or double-exponential decay rather than the
pure pinned-boundary hump. Depending on coefficients and boundary
conditions, such alternatives may approximate some observed shapes.

The correct empirical statement is conditional:

if the market is well described by the repeated-root IRFR branch,

then boundary selection provides a structural explanation for humped or
partially humped volatility profiles.

The empirical task is to test this conditional claim. One should compare
the repeated-root branch against distinct-root alternatives, richer
factor structures, and conventional volatility specifications.

The r = 1 branch should therefore be treated as a distinguished
candidate, not as an empirically established fact.

\section{Principle 6: distinguish scale from
shape}\label{principle-6-distinguish-scale-from-shape}

The observed magnitude of volatility and the maturity shape of
volatility are different empirical objects.

In the one-factor IRFR convention, the conventional volatility
coefficient is

\[
\sqrt{2} aH(\tau)
\]

Multiplying $H(\tau)$ by a constant changes the volatility scale, but not the
underlying shape after normalisation. Conversely, changing k or $\rho$
changes the shape of the loading, even if the total volatility level is
rescaled to match observed variance.

A useful empirical decomposition is therefore:

\begin{quote}
\item  volatility level;

\item  maturity decay;

\item  boundary anchoring;

\item  factor dimension.
\end{quote}

These are distinct dimensions of calibration.

The volatility level determines the overall size of rate shocks. The
maturity-decay scale determines how quickly the loading decays at long
maturities. The anchoring parameter determines how the loading behaves
near the short end. The factor dimension determines how many independent
stochastic sources are needed.

Without separating these objects, calibration can become misleading. A
model may fit variance levels while misidentifying the boundary regime.
Or it may fit a humped profile by adjusting scale and decay without
capturing the correct short-end behaviour.

Thus calibration should normalise loadings before comparing shapes, and
then estimate scale separately.

\section{Principle 7: use principal components
carefully}\label{principle-7-use-principal-components-carefully}

Empirical term-structure changes are often analysed using principal
components.

In the one-factor IRFR setting, the covariance kernel is

\[
K(\tau_1,\tau_2)=2a^2H(\tau_1)H(\tau_2)
\]

This kernel has rank one. Its single non-zero empirical direction is
proportional to $H(\tau)$.

Thus, within the one-factor restriction, the boundary-selected loading
may be compared directly with the leading empirical volatility mode of
forward-rate changes.

This provides a natural empirical procedure:

\begin{quote}
\item  estimate the leading maturity loading;

\item  normalise its scale;

\item  compare its shape with IRFR boundary-selected forms.
\end{quote}

However, PCA must be used carefully. Empirical yield-curve changes often
require more than one factor. Level, slope, and curvature components may
all be present. Measurement choices, curve construction, sampling
frequency, and overlapping maturities can affect the estimated
eigenfunctions.

Therefore, a mismatch between the one-factor IRFR loading and the
leading empirical component does not automatically invalidate the IRFR
framework. It may indicate the need for a multi-factor extension, a
different sampling horizon, or a more careful mapping from observed
rates to the IRFR field.

Likewise, a good match does not prove the theory. It supports the
interpretation that the leading empirical maturity mode may be
boundary-selected in the IRFR sense.

The empirical use of PCA should therefore be diagnostic, not definitive.

\section{Principle 8: treat boundary regimes as estimable
hypotheses}\label{principle-8-treat-boundary-regimes-as-estimable-hypotheses}

Neumann/free, Dirichlet, and Robin closures should not be assumed merely
because they are mathematically available.

They should be treated as competing empirical hypotheses.

The Neumann/free benchmark predicts retained short-end stochastic
loading through the ordinary exponential profile

$H_N(\tau)=A_Ne^{-a_N^2\tau}$.

A Dirichlet regime predicts suppressed short-end loading and, in the
repeated-root branch, the pure humped component

$H_D^{(1)}(\tau)=C_D\tau e^{-a_D^2\tau}$.

A Robin regime predicts a mixed repeated-root profile with partial
short-end retention,

$H_R(\tau)=(C_0+C_1\tau)e^{-a_R^2\tau}$.

These regimes imply different observable shapes near the short end.
Neumann/free closure retains boundary-active stochastic exposure.
Dirichlet closure suppresses it. Robin closure partially retains it
while allowing a boundary-induced linear-exponential component.

Thus the empirical question is not only:

does the model fit volatility?

It is also:

which boundary law best explains the estimated maturity loading?

This suggests a regime-testing procedure. Estimate the maturity loading,
normalise it, and compare the fit of Neumann/free, Dirichlet, and Robin
forms. Then assess whether the inferred boundary regime is stable across
time and consistent with the corresponding monetary-policy environment.

The economic interpretation matters. Dirichlet closure corresponds to
strong short-end anchoring or lower-bound enforcement. Robin closure
corresponds to partial policy anchoring. Neumann/free closure
corresponds to the non-pinned short-end benchmark. A useful empirical
result should connect the estimated boundary regime to the institutional
regime it is meant to represent.

The result should be interpreted probabilistically and comparatively.
The data may favour one boundary regime over another, but the conclusion
remains model-dependent.

Boundary laws are structural hypotheses. They must earn their empirical
interpretation through calibration and testing.

\section{Principle 9: check positivity and admissibility
separately}\label{principle-9-check-positivity-and-admissibility-separately}

The previous chapters introduced lower-boundary interpretations for
Dirichlet and sign-compatible Robin closures.

In the Dirichlet $r=1$ case, the selected loading $H_D^{(1)}(\tau)=C_D\tau e^{-a_D^2\tau}$

is non-negative on the half-line when $C_D\geq0$ and vanishes at $\tau=0$.
This supports a lower-boundary interpretation when the field is written
relative to an admissible minimum level,

\[
x(t,\tau)=x_{\min}+y(t,\tau), y(t,\tau)\ge0
\]

This lower-boundary interpretation is stronger than a static sign
observation. Under the shifted multiplicative KFE structure, the
Dirichlet branch admits a shifted inverse-gamma stationary density for
each fixed $\tau>0$. The same Feller classification as the
positive-state benchmark makes the lower boundary y = 0 inaccessible
maturity-by-maturity, provided the process starts in the positive state
domain and the parameter conditions are satisfied.

The remaining qualification is field-level regularity. To lift the
fixed-maturity result to the full maturity-field statement

\[
y(t,\tau)\ge0 for all \tau\ge0,
\]

one must assume or prove sufficient regularity of the field in the
maturity coordinate $\tau$. Under such regularity, fixed-$\tau$ Feller
inaccessibility supports preservation of the lower-boundary state domain
across the maturity field.

In the Robin case, the repeated-root loading can be written as

$H_R(\tau)=C_0[1+(a_R^2-\rho)\tau]e^{-a_R^2\tau}$.

For $C_0>0$, non-negativity on the full half-line requires

$\rho\leq a_R^2$.

When this sign condition is satisfied and the shifted positive-state KFE
structure is retained, Robin inherits the same maturity-by-maturity
inverse-gamma and Feller logic. If the Robin loading changes sign, the
simple lower-boundary interpretation no longer applies globally on the
half-line.

Calibration should therefore separate three levels of admissibility:

\begin{quote}
\item  loading-level sign compatibility;

\item  fixed-maturity state-domain preservation under the shifted KFE and
Feller conditions;

\item  field-level preservation across all maturities, which requires
maturity regularity.
\end{quote}

This distinction prevents overclaiming while also recognising that the
Dirichlet lower-boundary construction is supported by more than the sign
of the loading alone.

\section{Principle 10: empirical fit is evidence, not
proof}\label{principle-10-empirical-fit-is-evidence-not-proof}

The final principle is epistemic.

A successful calibration would be important. If the IRFR
boundary-selected loadings fit observed volatility profiles well, if the
estimated boundary regimes correspond sensibly to monetary-policy
environments, and if the parameters are stable and interpretable, then
the framework would gain empirical support.

But empirical fit is not proof of uniqueness or canonicality.

Other models may fit the same data. Other mechanisms may generate
similar volatility shapes. The IRFR framework may be useful without
being uniquely correct.

This is why the claims of the monograph must remain separated:

\begin{quote}
\item  admissibility is derived;

\item  minimality is argued;

\item  canonicality is conjectured;

\item  empirical relevance is tested.
\end{quote}

Empirical work can support the last two claims. It cannot replace the
mathematical derivation of admissibility, and it cannot by itself prove
canonicality.

This discipline is essential if the IRFR framework is to be taken
seriously.

\section{Minimal one-factor IRFR
calibration}\label{minimal-one-factor-irfr-calibration}

The preceding principles suggest a minimal one-factor calibration
programme.

\begin{enumerate}
\item Estimate the empirical maturity loading $H(\tau)$.
\item Normalise the loading to separate shape from scale.
\item Estimate $k$ from long-maturity decay or hump location.
\item Estimate $\rho$ from short-end slope where $H(0)\neq0$.
\item Compare Neumann/free, Dirichlet, and Robin boundary-regime fits.
\item Check loading sign compatibility and, where relevant, lower-boundary admissibility.
\item Test stability across time and policy regimes.
\end{enumerate}

This is the smallest empirical implementation that still respects the
structure of the theory.

It estimates a maturity loading, separates shape from level, identifies
the decay and boundary parameters, checks admissibility diagnostics, and
tests the boundary regime. It does not yet require a full multi-factor
extension, stochastic volatility, or option-pricing calibration.

The minimal model is therefore not the final empirical model. It is the
baseline test of the structural hypothesis.

\section{A practical calibration
sequence}\label{a-practical-calibration-sequence}

A more detailed empirical workflow proceeds as follows.

First, construct forward-rate or rolling-forward-rate changes from
market data. The construction should be consistent with the IRFR
maturity coordinate $\tau$.

Second, estimate the empirical covariance structure across maturity.
This may be done through historical changes, option-implied
volatilities, or a combination of both.

Third, extract the leading maturity loading, for example through PCA or
a parametric volatility fit.

Fourth, normalise the loading so that shape and scale are separated.

Fifth, estimate the decay scale k, or $a^2$ under the repeated-root
convention. If a Dirichlet-type hump is present, the hump location gives
the diagnostic $a^2\approx1/\tau_{\max}$.

Sixth, estimate the Robin anchoring parameter where applicable:
$\rho=-H'(0)/H(0)$.

Seventh, compare boundary-regime fits: Neumann/free, Dirichlet, Robin,
and, where appropriate, distinct-root alternatives.

Eighth, test stability across time and policy regimes. Parameters such
as $\rho$ may be time-varying, reflecting changing policy regimes rather than
fixed structural constants.

Ninth, check admissibility diagnostics, including short-end behaviour,
loading sign compatibility, and lower-boundary interpretation where
relevant.

Tenth, interpret the results as evidence, not proof.

This sequence is not the only possible empirical implementation. It is a
disciplined starting point.

\section{Summary}\label{summary-7}

This chapter has set out ten principles for taking the IRFR framework to
data.

The central empirical object is the maturity loading $H(\tau)$. It determines
how stochastic shocks are expressed across the curve. In the one-factor
model, the covariance kernel is rank one and the leading maturity
direction is proportional to $H(\tau)$.

The structural parameters have observable interpretations. The decay
scale controls long-maturity decay and, in the Dirichlet repeated-root
case, the hump location. The Robin anchoring parameter satisfies

$\rho=-H'(0)/H(0)$

when H(0) $\neq$ 0, though in practice both H(0) and $H'(0)$ may require
extrapolation. Boundary regimes imply different short-end behaviours and
different maturity-loading shapes.

The empirical task is to estimate these objects, compare boundary
regimes, and test stability across market and policy environments.

The chapter has also sharpened the lower-boundary discipline. Dirichlet
and sign-compatible Robin loadings support lower-boundary
interpretations, but the empirical analyst should distinguish
loading-level sign compatibility, fixed-maturity Feller support, and
full field-level maturity regularity.

The central caution is that empirical resemblance is not proof. Humped
volatility motivates the repeated-root branch, but does not validate it.
Good calibration supports the IRFR interpretation, but does not prove
uniqueness or canonicality.

The main message is:

boundary-selected loadings become empirical hypotheses when their
parameters are estimated from data.

\section{Transition}\label{transition-1}

The next chapter turns from empirical principles to scope, limitations,
and further directions.

The one-factor IRFR framework has been useful because it exposes the
maturity-domain geometry clearly. But empirical application,
multi-factor extension, field-level regularity, and no-arbitrage pricing
each require additional work.

The next chapter therefore states carefully what has been established,
what has been supported conditionally, and what remains open.

\chapter{Scope, Limitations, and Further Directions}\label{chapter-12}

\section{Introduction}\label{introduction-9}

The preceding chapters developed a minimal one-factor IRFR framework for
interest-rate dynamics.

The construction began with the instantaneous rolling forward rate,

and treated the forward curve as a stochastic field over remaining
maturity. The change from calendar maturity to rolling maturity exposed
the transport structure of the curve and identified the short end,

as a genuine boundary of the maturity domain.

The axioms then formalised this geometry. Axioms 1--3 defined the
interior law: local stochastic evolution, rolling-maturity invariance,
and equilibrium moment transport. Axiom 4 introduced boundary closure at
the short end.

The Kolmogorov forward equation was then used to convert the interior
axioms into calculable restrictions. This led to the affine
drift-operator structure, the admissibility condition for maturity
loadings, and the r-branch of admissible modes. The boundary chapters
then showed how natural, Dirichlet, and Robin closure select different
realised maturity loadings from the same underlying boundary-field
structure.

The purpose of the present chapter is not to extend the framework beyond
what has been derived. Its purpose is to state clearly what has been
established, what has been supported by the construction, what remains
conditional, and which questions remain open.

This distinction is important. The monograph has developed a minimal
boundary-field construction. It has not proved a universal theory of all
interest-rate dynamics. Nor has it established empirical validation, a
complete multi-factor theory, or a complete no-arbitrage pricing
framework.

The central claim is more precise:

\emph{within the one-factor IRFR field framework, admissible dynamics
are constrained by rolling maturity and boundary closure.}

This chapter records the scope of that claim.

\section{What has been
established}\label{what-has-been-established}

The first established result is representational.

The IRFR primitive

turns the forward curve into a field over remaining maturity. It is
equivalent to the instantaneous forward curve at the level of curve
values, but exposes a different geometric structure. In this
representation, fixed calendar maturities roll through the maturity
domain according to

This makes deterministic roll-down explicit and identifies

as the short-end boundary.

The second established result is axiomatic.

The four axioms organise the model into an interior law and a boundary
law. Axioms 1--3 govern the admissible dynamics away from the boundary.
Axiom 4 determines how the field is closed at the short end.

The third established result is the affine drift operator.

Under the convention

the KFE and first-moment transport imply the drift-operator form

The actual SDE drift is therefore

This separates mean-reverting displacement from the transport correction
induced by rolling maturity.

The fourth established result is the r-branch.

The branch parameter is

where $a^2$ is the temporal drift--diffusion scale and $\lambda^2$ is the
maturity-decay scale. For r $\neq$ 1, the admissible maturity modes are
distinct exponentials,

and

At

the two modes coalesce. The repeated-root space contains

and

This is the algebraic source of the linear-exponential humped component.

The fifth established result is boundary selection.

The Neumann/free benchmark retains the ordinary exponential short-end
stochastic loading,

Dirichlet closure suppresses short-end stochastic loading and, in the
repeated-root branch, selects the boundary-induced humped mode,

Robin closure gives the partially anchored mixed repeated-root loading,

Thus boundary closure does not merely impose an endpoint condition. It
selects the realised maturity structure of risk.

The sixth established result is the interpretation of volatility shape.

In the one-factor model,

and the instantaneous covariance rate is

Thus the maturity loading $H(\tau)$ is the observable volatility-shape
object. Boundary-selected loadings become empirical hypotheses when
their parameters are estimated from data.

The structural parameters $a^2$, r, and, in the Robin case, $\rho$, are
identifiable through features of the maturity loading, though only under
the modelling assumptions and estimation procedures described in Chapter
11.

\section{What has been supported, but remains
conditional}\label{what-has-been-supported-but-remains-conditional}

Some conclusions are stronger than conjecture, but still depend on
clearly stated technical conditions.

The most important example is lower-boundary admissibility in the
Dirichlet branch.

In the Dirichlet lower-boundary interpretation, write

The Dirichlet boundary imposes

For $\tau>0$, the repeated-root Dirichlet loading is

Under the positive normalisation $C_D\geq0$, this loading is non-negative
on the half-line and vanishes at the short end.

More importantly, under the shifted multiplicative KFE structure, the
Dirichlet branch remains in the positive-state diffusion class. For each
fixed $\tau>0$, the shifted state may be represented in the
form

with positive parameter functions. The stationary KFE then admits a
shifted inverse-gamma density on

Equivalently, in the original variable,

The same Feller boundary classification as in the positive-state
benchmark shows that the lower boundary y = 0 is inaccessible for each
fixed $\tau>0$, provided the process starts in the positive
state domain and the parameter conditions are satisfied.

Thus the lower-boundary interpretation is not merely a static sign
property of the loading. It is supported by the shifted inverse-gamma
stationary density and by maturity-by-maturity Feller inaccessibility of
the lower state boundary. The detailed density and
boundary-classification calculation is given in Appendix B.

The remaining qualification is field-level regularity. To pass from
fixed-$\tau$ positivity to the full maturity-field statement

one also requires sufficient regularity of the field in the maturity
coordinate. Under such regularity,

and

are compatible with preservation of the lower-boundary state domain
across the maturity field.

Robin closure inherits the same local lower-boundary logic when its
loading is sign-compatible. For

with $C_0>0$, non-negativity on the full half-line requires

If this sign condition fails, the simple lower-boundary interpretation
no longer applies globally on the half-line.

\section{What has not been
established}\label{what-has-not-been-established}

The framework has not proved uniqueness.

The monograph has shown that the proposed construction gives an
admissible and structurally disciplined model class. It has not proved
that this class is the unique possible construction satisfying the
axioms, nor that no other coordinate choice or admissibility mechanism
could generate similar structures.

The framework has not empirically validated the model.

The resemblance between boundary-selected loadings and familiar humped
volatility forms motivates calibration and testing. It does not prove
that the IRFR model is empirically correct. Nor does it prove that the r
= 1 branch is the true branch. Empirical validation remains a separate
task.

The framework has not developed a complete multi-factor theory.

The present monograph has deliberately focused on the one-factor case.
The distinction between stochastic dimension and maturity geometry was
central to the exposition. A multi-factor IRFR field may be natural, but
it has not been derived here. In particular, it is not guaranteed that
multi-factor extensions preserve the same admissibility structure or
boundary-selection mechanism without further work.

The framework has not provided a complete no-arbitrage or risk-neutral
pricing theory.

This omission is deliberate rather than merely technical. Conventional
arbitrage-free term-structure modelling is usually built from a
numeraire, often the money-market account, and imposes martingale
conditions on discounted asset prices. The construction therefore begins
with a specified universe of tradable instruments and a strictly positive
tradable numeraire.

The IRFR primitive is different. The object $x(t,	au)$ is a local
rate-field coordinate, not a positive tradable price process. A
unit-notional exposure to an instantaneous IRFR level or increment is a
local rate exposure rather than a self-financing asset. In particular,
such an exposure has no natural positive initial value and cannot itself
serve as a numeraire. This is why the standard risk-free portfolio or
numeraire argument cannot be applied directly to ``IRFR units'' alone.

The affine-closure calculation should therefore be read as a
field-level admissibility condition. It constrains relative stochastic
motion across maturities by requiring local increments to be compatible
with a common loading structure. This is no-arbitrage-like in economic
motivation, but it is not a substitute for a self-financing
no-arbitrage theorem.

A direct risk-neutral formulation would need to add a pricing layer:
first mapping the IRFR field into tradable instruments such as bonds,
FRAs, swaps, futures, or other specified contracts; then choosing a
numeraire and pricing measure; and finally deriving the corresponding
martingale restrictions. The issue is not merely to translate the model
into HJM notation, but to determine how the boundary-field
interpretation interacts with no-arbitrage pricing when the short end is
itself a structural boundary.

For this reason, the monograph treats risk-neutral pricing as an open
problem of a different kind from empirical calibration. The present work
develops the structural admissibility framework: rolling maturity,
affine closure, admissible modes, and boundary-selected loadings. A
no-arbitrage version may be possible, but it requires a separate
treatment of numeraire choice, discounting conventions, tradable
instruments, market price of risk, and the relation between physical and
pricing measures.

These limitations are not defects. They are the boundaries of the
present work.

\section{Status of the r = 1
branch}\label{status-of-the-r-1-branch}

The r = 1 branch has played a central role in the exposition because it
is mathematically distinguished.

At r = 1, the maturity-decay scale and the temporal drift--diffusion
scale coincide:

The two exponential modes

and

coalesce. The repeated-root space then contains the linear-exponential
component

This term is humped. Under Dirichlet closure, it becomes the selected
loading

This structure is close to humped volatility specifications used in
interest-rate modelling. That connection is important, but it should be
read carefully.

The empirical usefulness of humped volatility forms does not prove that
the r = 1 branch is correct. It shows that the branch is economically
interpretable and worth testing. The theoretical result is conditional:

\emph{if the repeated-root branch is relevant, boundary selection gives
a structural route to humped loadings.}

Whether the repeated-root branch is empirically preferred over
distinct-root alternatives is a calibration question.

Thus the r = 1 branch is distinguished within the theory, not
empirically established by resemblance alone.

\section{Status of the boundary-law
interpretation}\label{status-of-the-boundary-law-interpretation}

The boundary-law interpretation is one of the main contributions of the
monograph.

The short end is both a mathematical boundary and the location where
monetary policy acts most directly. This makes boundary closure a
natural place to represent economic regime.

The boundary chapters developed three canonical closures:

\begin{center}\emph{Neumann/free closure: short-end stochastic loading retained;}\end{center}
\begin{center}\emph{Dirichlet closure: short-end stochastic loading suppressed;}\end{center}
\begin{center}\emph{Robin closure: short-end stochastic loading partially retained.}\end{center}
These should be read as structural idealisations. Real monetary-policy
regimes are more complex than any single boundary condition.
Nevertheless, the classification is useful because it separates the
interior admissible dynamics from the economic regime that selects the
realised loading.

The strongest claim is therefore not that every policy regime is
literally a Dirichlet or Robin boundary. The stronger and safer claim
is:

\emph{policy anchoring can be represented as a boundary-selection
mechanism in the IRFR field.}

This is a structural interpretation. Its empirical content depends on
whether estimated volatility loadings behave as the boundary-law
classification predicts.

\section{Minimality and modelling
cost}\label{minimality-and-modelling-cost}

The framework aims to be minimal in a specific sense.

It does not begin by specifying a volatility curve and then deriving the
required arbitrage-free drift. Instead, it begins with a
rolling-maturity field, imposes structural axioms, and asks which drift
and loading structures are admissible.

This changes the modelling cost.

In a flexible reduced-form model, one can introduce a volatility
function because it fits data. In the IRFR framework, a maturity loading
must be justified as part of the admissible field structure and boundary
closure.

This reduces exogenous model intervention, but it also reduces freedom.
Some useful empirical shapes may be outside the minimal model. Some
market features may require additional factors, stochastic volatility,
time-varying boundary laws, or richer state variables.

This is the trade-off.

The minimal model is valuable not because it captures everything, but
because it identifies which features can be derived from a disciplined
interior-law and boundary-law construction.

This minimality also clarifies the relationship with richer nonlinear
term-structure models, including quadratic specifications. Quadratic
models obtain nonlinear yield behaviour by enriching the state
dependence. The IRFR construction asks what nonlinear-looking maturity
structure can be obtained before that step, from rolling-maturity
geometry, admissibility, and boundary selection alone.

\section{Further directions}\label{further-directions}

Several further directions are natural, but they should be regarded as
future work.

A first direction is empirical testing. The calibration principles of
Chapter 11 provide a starting point. Actual validation requires data,
estimation procedures, robustness checks, and comparison against
competing models.

A second direction is to formulate the maturity-regularity conditions
under which fixed-$\tau$ Feller inaccessibility lifts to lower-boundary
preservation of the full maturity field. The Dirichlet branch already
supports lower-boundary admissibility maturity-by-maturity; the
remaining technical task is to state the appropriate function-space
conditions for the whole field.

A third direction is multi-factor extension. A multi-factor IRFR model
would allow several Brownian drivers and several maturity loadings. Such
a model may be necessary for empirical yield-curve dynamics. However,
the admissibility and boundary-selection analysis would need to be
developed again, not merely assumed.

A fourth direction is time-varying boundary parameters. In the Robin
case, the anchoring ratio

could vary through time. This would allow changing policy regimes to
appear as changing boundary laws. The present monograph identifies this
possibility but does not estimate or model such dynamics.

A fifth direction is stochastic boundary laws. Instead of fixed boundary
coefficients, the boundary itself could be driven by a stochastic
process representing policy uncertainty or liquidity conditions. This
may be economically realistic, but it is beyond the minimal framework.

A sixth direction is no-arbitrage pricing. As discussed above, this is
not a routine extension. It requires a careful treatment of the
relationship between the boundary-field structure and the conventional
numeraire-based pricing construction, including the mapping from local
IRFR rate exposures to positive tradable price processes.

These directions are not part of the established result. They are
consequences of the programme and possible areas for future research.

\section{Why the minimal model still
matters}\label{why-the-minimal-model-still-matters}

The limitations just described do not undermine the value of the minimal
model.

A minimal model is useful because it reveals structure that may be
obscured in more flexible specifications. The one-factor IRFR framework
shows how rolling maturity, moment transport, admissible modes, and
boundary closure interact.

It also shows that humped maturity loading can arise as a structural
consequence rather than a modelling input. In the repeated-root
Dirichlet case, the hump is a boundary-selected admissible mode. In the
Robin case, partial anchoring selects a mixed profile.

These are structural results. They do not require the minimal model to
be the final empirical model.

The minimal model therefore plays the role of a benchmark. It gives a
disciplined baseline against which extensions, empirical calibrations,
and alternative specifications can be compared.

The central organising principle remains:

\begin{center}\emph{interior admissibility determines the mode space;}\end{center}
\begin{center}\emph{boundary closure selects the realised loading.}\end{center}
This principle is the main object that any extension should preserve.

\section{Summary}\label{summary-8}

This chapter has stated the scope and limitations of the IRFR
construction.

The monograph has established a minimal one-factor boundary-field
framework. It has shown how the IRFR primitive reveals rolling-maturity
transport, how the axioms divide the theory into an interior law and a
boundary law, how the KFE produces an affine drift operator, how the
r-branch arises, and how boundary closure selects observable loadings.

It has also strengthened the lower-boundary interpretation of the
Dirichlet branch. Under shifted positive-state dynamics, the Dirichlet
branch admits a shifted inverse-gamma stationary density for fixed
maturities $\tau>0$, and the lower boundary is inaccessible
maturity-by-maturity under the Feller conditions. Field-level
preservation requires the corresponding regularity in the maturity
coordinate.

The monograph has not proved uniqueness. It has not empirically
validated the model. It has not developed a full multi-factor framework.
It has not solved the no-arbitrage pricing problem associated with
embedding a boundary-field short end inside a numeraire-based pricing
construction.

These limitations define the boundary of the current work.

The value of the framework lies in the discipline it imposes. It shows
how drift structure, maturity loading, volatility shape, boundary
anchoring, and lower-boundary admissibility can be traced to specific
locations in the theory.

The appropriate final claim is therefore measured:

\emph{the IRFR framework provides a minimal structural route from
rolling maturity and boundary closure to admissible interest-rate
dynamics.}

\section{Transition to the
conclusion}\label{transition-to-the-conclusion}

The final chapter returns to the overall programme of the monograph.

It summarises the construction, restates the main results, identifies
the conjectural claims, and explains why the boundary-law formulation
provides a useful way to think about interest-rate dynamics.

The purpose of the conclusion is not to enlarge the claims, but to close
the argument carefully: what has been derived, what has been supported
by the construction, what has been motivated, and what remains open.

\chapter{Conclusion: A Minimal Boundary-Field Theory of Interest Rates}\label{chapter-13}

\section{Introduction}\label{introduction-10}

This monograph has developed a minimal boundary-field framework for
interest-rate dynamics.

The starting point was the instantaneous rolling forward rate,

\[
x(t,\tau)=F(t,t+\tau),
\]

where \(t\) is calendar time and $\tau$ is remaining maturity. This change of
primitive does not create a new market object. It is equivalent to the
conventional instantaneous forward curve at the level of curve values.
But it exposes a different geometric structure.

In rolling-maturity coordinates, fixed calendar maturities move through
the maturity domain according to

\[
d\tau/dt = -1
\]

The forward curve therefore becomes a field over the half-line

\[
\tau\in[0,\infty),
\]

and the short end

\[
\tau=0
\]

becomes a boundary.

The central idea of the monograph is that this geometry matters. Once
the forward curve is treated as a rolling maturity-domain field, drift,
maturity loading, volatility shape, and policy anchoring can be
organised through a common structure:

interior admissibility determines the mode space;

boundary closure selects the realised loading.

The purpose of this final chapter is to summarise what has been
constructed, state the main claims at the correct strength, and identify
what remains open.

\section{The problem revisited}\label{the-problem-revisited}

Interest-rate models face a persistent tension.

On the one hand, term-structure models must be flexible enough to fit
observed curves, volatilities, and derivatives prices. On the other
hand, excessive flexibility can obscure which features of the model are
structural and which are inserted for calibration.

This tension is especially visible in volatility modelling. Humped
volatility profiles are empirically useful and widely recognisable. But
in many formulations, the maturity shape of volatility is specified
directly. The model may fit observed behaviour, but the shape itself
remains an input.

The IRFR framework asks a different question.

which maturity loadings are admissible under rolling-maturity geometry?

which of those admissible loadings are selected by boundary closure?

This changes the modelling problem. It moves part of the burden from
exogenous specification to structural derivation.

The framework does not eliminate modelling choice. It relocates it.
Assumptions enter through the primitive, the axioms, the admissibility
calculation, and the boundary law. The point is that those assumptions
are made visible.

\section{The IRFR primitive}\label{the-irfr-primitive-1}

The first modelling commitment was to use the instantaneous rolling
forward rate,

\[
x(t,\tau)=F(t,t+\tau)
\]

The conventional representation F(t,T) labels the curve by calendar
maturity. The IRFR representation labels the curve by remaining
maturity.

This makes the maturity coordinate dynamic. A fixed maturity date T
corresponds to

\[
\tau=T-t,
\]

so calendar time transports the maturity point toward the short end.

This simple observation has three consequences.

First, deterministic roll-down becomes explicit. A change in a
fixed-maturity forward rate combines stochastic evolution of the field
with deterministic movement along the maturity axis.

Second, the forward curve becomes a field over maturity, not merely a
collection of unrelated maturity-indexed processes.

Third, the short end becomes a boundary. It is not added later as a
modelling convenience. It is the endpoint of the maturity domain toward
which contracts roll.

This is the geometric foundation of the monograph.

\section{The axioms}\label{the-axioms}

The field representation led to four axioms.

Axiom 1 gave the field local stochastic evolution. Under the convention
used in the monograph,

\[
dx(t,\tau) = -f(x,t,\tau)\,dt + \sqrt{2} g(x,t,\tau)\,dW_t
\]

Here f is the IRFR drift-operator coefficient, g is the diffusion-scale
coefficient, and the actual SDE drift is -f.

Axiom 2 imposed rolling-maturity invariance. The same remaining maturity
should be governed by the same structural law within a given regime,
regardless of the calendar date at which it is observed.

Axiom 3 required equilibrium moment transport. The expected dynamics
must be compatible with deterministic roll-down along the equilibrium
curve $m(\tau)$.

Axiom 4 required boundary closure at

\[
\tau=0
\]

Together, these axioms divided the model into two parts:

Axioms 1-3: interior law;

Axiom 4: boundary law.

This division is the organising principle of the framework.

\section{The affine drift operator}\label{the-affine-drift-operator}

The Kolmogorov forward equation turned the first three axioms into
calculable restrictions.

The first-moment equation implied that the drift operator must close at
the level of the first moment. Requiring compatibility with equilibrium
transport led to the affine drift-operator form

\[
f(x,t,\tau)=a^2{[}x-m(\tau){]}-m'(\tau)
\]

The actual SDE drift is therefore

\[
-f(x,t,\tau) = -a^2{[}x-m(\tau){]} + m'(\tau)
\]

This decomposition has a clear interpretation.

$-a^2[x-m(\tau)]$

is the mean-reverting displacement term in the actual drift.

$m'(\tau)$

is the fixed-$\tau$ transport correction generated by rolling maturity.

Thus the drift structure is not merely chosen for tractability. It
reflects the interaction between stochastic displacement and
deterministic maturity transport.

This was the first major payoff of the IRFR formulation.

\section{Admissible maturity modes}\label{admissible-maturity-modes}

The diffusion-scale coefficient was then written in maturity-loading
form,

\[
g(\tau)=aH(\tau)
\]

The KFE and moment-closure conditions imply that the maturity loading
cannot be arbitrary.

The branch parameter is

\[
r = \lambda^2/a^2,
\]

where $a^2$ is the temporal drift-diffusion scale and $\lambda^2$ is the
maturity-decay scale. For $r\neq1$, the admissible maturity modes are the
distinct exponentials

$e^{-a^2\tau}$ and $e^{-ra^2\tau}$.

At $r=1$, the two modes coalesce. The repeated-root space contains $e^{-a^2\tau}$ and $\tau e^{-a^2\tau}$.

The repeated-root branch is therefore mathematically distinguished
because it generates the linear-exponential component

\[
\tau e^{-a^2\tau}
\]

This component is humped. It vanishes at the short end, rises to an
interior maximum, and decays at long maturity.

The repeated-root branch is structurally important and economically
interpretable. But it is not empirically established by resemblance to
humped volatility alone. It remains a distinguished branch to be tested.

\section{Boundary selection}\label{boundary-selection}

The boundary law determines which admissible loading is realised.

Neumann/free closure provides the non-pinned benchmark. It retains
the ordinary exponential short-end stochastic loading,

$H_N(\tau)=A_Ne^{-a_N^2\tau}$.

The short-end loading is retained:

$H_N(0)=A_N$.

Dirichlet closure represents the idealised pinned-boundary case. It
suppresses short-end stochastic loading and, in the repeated-root
branch, selects the boundary-induced humped mode

$H_D^{(1)}(\tau)=C_D\tau e^{-a_D^2\tau}$.

This loading satisfies

$H_D^{(1)}(0)=0$.

Robin closure represents partial anchoring. It retains residual
short-end stochastic loading while allowing a boundary-induced
repeated-root component:

$H_R(\tau)=(C_0+C_1\tau)e^{-a_R^2\tau}$.

When $C_0\neq0$, the same loading may be written in anchoring-ratio form as

$H_R(\tau)=C_0[1+(a_R^2-\rho)\tau]e^{-a_R^2\tau}$,

where

$\rho=-H_R'(0)/H_R(0)$.

These cases demonstrate the boundary-law principle:

the interior law creates admissible modes;

the boundary law selects the realised combination.

\section{Observable volatility}\label{observable-volatility}

The boundary-selected maturity loading determines observable volatility
shape.

In the one-factor IRFR model,

\[
g(\tau)=aH(\tau)
\]

Define the instantaneous covariance rate by

\[
K(\tau_1,\tau_2) = (1/dt) \operatorname{Cov}(dx(t,\tau_1), dx(t,\tau_2))
\]

Then

\[
K(\tau_1,\tau_2)=2a^2H(\tau_1)H(\tau_2)
\]

This kernel is rank one. Its maturity eigenfunction is proportional to
$H(\tau)$.

Thus the empirical object corresponding to the theory is not merely a
volatility level. It is a maturity loading. The shape of this loading
contains information about decay, anchoring, and boundary regime.

For the Dirichlet repeated-root loading,

$H_D^{(1)}(\tau)=C_D\tau e^{-a_D^2\tau}$,

the maximum of the loading occurs at

\[
\tau_{\max} = 1/a_D^2
\]

For the Robin loading,

$H_R(\tau)=(C_0+C_1\tau)e^{-a_R^2\tau}$,

the anchoring parameter satisfies

$\rho=-H_R'(0)/H_R(0)$

whenever $H_R(0)\neq0$.

These relationships turn boundary-law parameters into empirical
diagnostics. But they do not by themselves validate the model. They
define what should be estimated and tested.

\section{What has been claimed}\label{what-has-been-claimed}

The main claim of the monograph is structural.

Within the one-factor IRFR field framework, rolling-maturity geometry
and boundary closure constrain admissible interest-rate dynamics. Drift
structure, maturity loading, volatility shape, and policy anchoring are
not independent modelling choices. They are connected through the
interior-law and boundary-law construction.

More specifically, the monograph has claimed that:

the IRFR primitive reveals transport and boundary geometry;

the axioms divide the model into interior law and boundary law;

the KFE gives calculable admissibility restrictions;

the affine drift operator follows from first-moment transport;

maturity loadings arise as admissible modes;

boundary closure selects observable loading shapes.

These are the earned claims.

The contribution is therefore not a new arbitrary volatility
parametrisation, but a mechanism by which volatility loadings are
selected from an admissible maturity-mode space.

The framework is also not presented as a replacement for conventional
arbitrage-free representations such as Heath-Jarrow-Morton. Rather, it
gives a structural explanation for why certain maturity loadings may
arise within such representations. Heath-Jarrow-Morton describes the
no-arbitrage representation of forward-rate dynamics; IRFR provides a structural route by which
rolling-maturity geometry and boundary closure may select particular
admissible loadings.

\section{What has not been claimed}\label{what-has-not-been-claimed}

Several stronger claims have deliberately not been made.

The monograph has not proved uniqueness. It has shown an admissible and
disciplined construction, not that no other construction is possible.

It has not empirically validated the r=1 branch. Humped volatility
motivates testing the repeated-root branch, but does not prove it.

It has not provided a complete multi-factor theory. Multi-factor
extensions may be natural, but they require new derivation.

It has not provided a complete no-arbitrage or risk-neutral pricing
implementation. The relationship with Heath-Jarrow-Morton is important,
but the issue is not merely to translate the model into conventional HJM
notation. The IRFR value $x(t,	au)$ is a local rate-field coordinate,
not itself a strictly positive tradable price process; a unit-notional
IRFR exposure cannot directly play the role of a numeraire. A full
pricing formulation would need to address the mapping from IRFR field
exposures to tradable instruments, the numeraire, the money-market
account, market price of risk, discounting conventions, and the way a
boundary-field short end interacts with the zero-return or
instantaneous-rate pricing construct.

The lower-boundary claim is also scoped carefully. The monograph does
not assert a fully developed field-level positivity theorem without
qualification. Rather, the Dirichlet branch supports lower-boundary
admissibility maturity-by-maturity through the shifted inverse-gamma
stationary density and Feller inaccessibility of the lower state
boundary for each fixed $\tau>0$. The remaining technical
qualification is the maturity regularity needed to lift fixed-maturity
positivity to the full field statement.

These limitations define the boundary of the present contribution.

\section{Why the framework matters}\label{why-the-framework-matters}

The value of the IRFR framework lies in its discipline.

It provides a way to connect objects that are often specified
separately:

drift;

volatility loading;

maturity geometry;

short-end policy anchoring.

In the IRFR field, these objects are organised by a common structure.

This does not make the model less conditional. It makes the conditions
clearer.

The framework also gives a structural interpretation of humped maturity
loading. In the Dirichlet repeated-root case, the hump arises as a
boundary-selected admissible mode. In the Robin case, partial anchoring
selects a mixed profile that can retain short-end loading while still
producing curvature or intermediate-maturity concentration.

Thus humped volatility can arise as a structural consequence rather than
a modelling input.

The Dirichlet branch also shows how a lower-boundary regime can be
supported by the shifted positive-state density and Feller boundary
structure, not merely by the sign of the loading.

This is the central modelling payoff.

\section{Open problems}\label{open-problems}

The monograph leaves several open problems.

The first is uniqueness. It remains to determine whether the admissible
construction is unique under the axioms, or whether alternative
admissible model classes exist.

The second is field-level lower-boundary preservation. The
fixed-maturity Dirichlet argument gives shifted inverse-gamma stationary
density and Feller inaccessibility for $\tau>0$. A full
maturity-field theorem requires the corresponding regularity in $\tau$, so
that the fixed-maturity result lifts to the entire half-line.

The third is empirical testing. The calibration principles developed
earlier must be applied to actual market data. Boundary regimes should
be estimated, compared, and tested across policy environments.

The fourth is multi-factor extension. Observed term-structure dynamics
often require more than one factor. Extending the IRFR admissibility and
boundary-law construction to multiple factors is a substantial task.

The fifth is no-arbitrage pricing. The connection to Heath-Jarrow-Morton
must be developed into a full pricing implementation, but doing so
requires confronting the relation between the boundary-field short end,
the money-market numeraire, discounted martingale conditions, and the
conventional zero-return pricing construction.

These problems define a research programme, not a completed theory.

\section{Final statement}\label{final-statement}

The aim of this monograph has been to develop a minimal structural route
from rolling maturity and boundary closure to admissible interest-rate
dynamics.

The construction is deliberately narrow. It is a one-factor
boundary-field framework. It does not claim universality. It does not
claim empirical validation. It does not claim uniqueness.

Its claim is that the rolling-maturity representation reveals a
boundary-field structure that is hidden in the conventional
calendar-maturity description. Once this structure is made explicit, the
short end becomes a genuine boundary, and the modelling problem
separates naturally into an interior admissibility problem and a
boundary-selection problem.

Under the axioms developed here, the interior law determines the
admissible mode space. Boundary closure then selects the realised
maturity loading. In this way, drift, maturity loading, observable
volatility shape, and policy anchoring become linked components of a
single structural construction rather than independent modelling inputs.

The repeated-root branch illustrates the payoff of this construction. It
produces the linear-exponential component

\[
\tau e^{-a^2\tau},
\]

which gives a natural humped maturity profile. Under Dirichlet closure
this component is selected directly. Under Robin closure it appears as
part of a mixed boundary-active loading. Thus familiar humped and
partially humped volatility shapes can be interpreted as
boundary-selected admissible modes.

This does not prove that markets are governed by the repeated-root
branch. It does not prove that any particular boundary regime is
empirically correct. It shows something more limited but still useful: a
class of empirically recognisable volatility loadings can arise from a
minimal boundary-field structure rather than being imposed exogenously.

The resulting theory is therefore best understood as a disciplined
modelling framework. It identifies a primitive, states axioms, derives
admissible dynamics, and interprets observable loading shapes through
boundary closure. Its value lies not in claiming finality, but in making
the structural assumptions visible and testable.

The central conclusion is:

rolling maturity creates a boundary;

the interior law determines admissible modes;

the boundary law selects observable loadings.

The forward curve is therefore not only a curve indexed by maturity
dates. In rolling-maturity coordinates it is a field on a half-line,
transported toward a boundary. Once that boundary is recognised,
volatility shape can be interpreted not merely as a calibration input,
but as the observable footprint of boundary selection.

This is the minimal boundary-field theory of interest rates developed in
this monograph.


\appendix

\chapter{Notation and Conventions}\label{appendix-a}

\section{Purpose of this appendix}
This appendix collects the notation and sign conventions used throughout the monograph.  The main distinctions are between calendar maturity and remaining maturity, between the drift-operator coefficient and the actual SDE drift, and between the diffusion-scale coefficient and the conventional volatility coefficient.

\section{Time and maturity variables}
Calendar time is denoted by
\[
t.
\]
Calendar maturity is denoted by
\[
T.
\]
Remaining maturity is denoted by
\[
\tau.
\]
The relationship is
\[
\tau=T-t,
\qquad
\frac{d\tau}{dt}=-1.
\]

\section{Forward rate and IRFR primitive}
The conventional instantaneous forward rate is
\[
F(t,T).
\]
The instantaneous rolling forward rate is
\[
x(t,\tau)=F(t,t+\tau),
\]
with inverse relation
\[
F(t,T)=x(t,T-t).
\]
Thus the IRFR is not a new market object at the level of curve values. It is the conventional forward curve expressed in rolling-maturity coordinates.

\section{Transport identity}
Differentiating
\[
F(t,T)=x(t,T-t)
\]
along a fixed calendar maturity gives
\[
\frac{d}{dt}F(t,T)
=
\partial_t x(t,\tau)-\partial_\tau x(t,\tau),
\qquad \tau=T-t.
\]
The first term is calendar-time field evolution. The second is deterministic maturity transport.

\section{IRFR stochastic convention}
The one-factor IRFR field is written as
\[
dx(t,\tau)
=
-f(x,t,\tau)\,dt
+
\sqrt{2}\,g(x,t,\tau)\,dW_t.
\]
Here \(f\) is the IRFR drift-operator coefficient. The actual SDE drift is
\[
-f(x,t,\tau).
\]
The diffusion-scale coefficient is \(g\), and the conventional SDE volatility coefficient is
\[
\sqrt{2}\,g(x,t,\tau).
\]
The factor \(\sqrt{2}\) is a normalisation convention.

\section{Kolmogorov forward equation convention}
Under
\[
dx=-f\,dt+\sqrt{2}\,g\,dW_t,
\]
the density \(p(x,t;\tau)\) satisfies
\[
\frac{\partial p}{\partial t}
=
\frac{\partial}{\partial x}(fp)
+
\frac{\partial^2}{\partial x^2}(g^2p).
\]
Equivalently, if
\[
v=2g^2,
\]
then
\[
\frac{\partial p}{\partial t}
=
\frac{\partial}{\partial x}(fp)
+
\frac{1}{2}\frac{\partial^2}{\partial x^2}(vp).
\]
The first-moment equation is
\[
\frac{\partial M_1}{\partial t}
=
-\mathbb{E}[f(x,t,\tau)].
\]

\section{Affine drift operator}
The equilibrium or reference curve is denoted by \(m(\tau)\). The affine drift-operator coefficient derived in the monograph is
\[
f(x,t,\tau)=a^2[x-m(\tau)]-m'(\tau).
\]
Therefore the actual SDE drift is
\[
-f(x,t,\tau)=-a^2[x-m(\tau)]+m'(\tau).
\]
The first term is mean-reverting displacement. The second is the fixed-\(\tau\) transport correction.

\section{Diffusion loading and covariance rate}
In the one-factor loading representation,
\[
g(\tau)=aH(\tau).
\]
The conventional volatility loading is
\[
\sqrt{2}\,aH(\tau).
\]
The instantaneous covariance of increments is
\[
\operatorname{Cov}\bigl(dx(t,\tau_1),dx(t,\tau_2)\bigr)
=
2a^2H(\tau_1)H(\tau_2)\,dt.
\]
The covariance-rate kernel is
\[
K(\tau_1,\tau_2)
=
\frac{1}{dt}\operatorname{Cov}\bigl(dx(t,\tau_1),dx(t,\tau_2)\bigr)
=
2a^2H(\tau_1)H(\tau_2).
\]

\section{Admissible maturity modes}
In the distinct-root case,
\[
H(\tau)=C_1e^{-k_1\tau}+C_2e^{-k_2\tau}.
\]
In the repeated-root case,
\[
H(\tau)=(C_1+C_2\tau)e^{-k\tau}.
\]
Under the central \(r=1\) normalisation,
\[
k=a^2.
\]
The repeated-root branch contains the linear-exponential component
\[
\tau e^{-k\tau},
\]
which is the source of the humped loading selected by Dirichlet closure.

\section{The \texorpdfstring{\(r\)}{r}-branch}
The branch parameter is
\[
r=\frac{\lambda^2}{a^2},
\]
where \(\lambda^2\) is the maturity-decay scale and \(a^2\) is the temporal drift-diffusion scale.  The case \(r\neq1\) gives distinct exponential modes. The case \(r=1\) gives coalescence and the repeated-root loading.

\section{Boundary regimes}
The maturity domain is
\[
\tau\in[0,\infty),
\]
with boundary at \(\tau=0\). The monograph studies three canonical boundary regimes.

The Neumann/free benchmark retains the ordinary exponential short-end stochastic mode:
\[
H_N(\tau)=A_Ne^{-a_N^2\tau}.
\]
The terminology follows the convention fixed in Chapter 6: the operative condition is non-pinning, represented here by the retained exponential loading.

Dirichlet closure suppresses short-end stochastic loading and, in the repeated-root branch, selects
\[
H_D^{(1)}(\tau)=C_D\tau e^{-a_D^2\tau}.
\]
Robin closure represents partial anchoring and selects the mixed repeated-root loading
\[
H_R(\tau)=(C_0+C_1\tau)e^{-a_R^2\tau}.
\]
When \(C_0\neq0\), this can be written as
\[
H_R(\tau)
=
C_0\bigl[1+(a_R^2-\rho)\tau\bigr]e^{-a_R^2\tau},
\]
where
\[
\rho=-\frac{H_R'(0)}{H_R(0)}=a_R^2-\frac{C_1}{C_0}.
\]

\section{Lower-boundary convention}
For the Dirichlet lower-boundary interpretation, write
\[
y(t,\tau)=x(t,\tau)-x_{\min}.
\]
The Dirichlet boundary imposes
\[
y(t,0)=0.
\]
For \(\tau>0\), the shifted positive-state structure supports a shifted inverse-gamma stationary density and Feller inaccessibility of the lower state boundary maturity-by-maturity. Full field-level lower-boundary preservation requires regularity in the maturity coordinate.

\chapter{Algebra of Affine Closure, the \texorpdfstring{\(r\)}{r}-Branch, and Boundary-Selected Loadings}\label{appendix-b}

\section{Purpose of this appendix}
This appendix records the algebra behind Chapters 4--9: the KFE convention, affine drift closure, the \(r\)-branch, the coalescing-root Dirichlet limit, and the boundary-selected loadings.

\section{Stochastic convention and KFE}
For fixed remaining maturity \(\tau\), write
\[
dx(t,\tau)=-f(x,t,\tau)\,dt+\sqrt{2}\,g(x,t,\tau)\,dW_t.
\]
The Kolmogorov forward equation is
\[
\frac{\partial p}{\partial t}
=
\frac{\partial}{\partial x}(fp)
+
\frac{\partial^2}{\partial x^2}(g^2p).
\]

\section{First moment}
Let
\[
M_1(t,\tau)=\mathbb{E}[x(t,\tau)]=\int x p(x,t;\tau)\,dx.
\]
Then
\[
\frac{\partial M_1}{\partial t}
=
\int x\frac{\partial p}{\partial t}\,dx.
\]
Substituting the KFE and integrating by parts gives
\[
\int x\frac{\partial}{\partial x}(fp)\,dx=-\int fp\,dx=-\mathbb{E}[f],
\]
and
\[
\int x\frac{\partial^2}{\partial x^2}(g^2p)\,dx=0.
\]
Therefore
\[
\frac{\partial M_1}{\partial t}=-\mathbb{E}[f(x,t,\tau)].
\]

\section{Affine closure}
First-moment closure requires \(\mathbb{E}[f]\) to depend only on \(M_1\). Thus take
\[
f(x,t,\tau)=A(\tau)x+B(\tau).
\]
Choosing \(A(\tau)=a^2\), imposing mean reversion toward \(m(\tau)\), and incorporating transport gives
\[
B(\tau)=-a^2m(\tau)-m'(\tau).
\]
Thus
\[
f(x,t,\tau)=a^2[x-m(\tau)]-m'(\tau),
\]
and the actual drift is
\[
-f(x,t,\tau)=-a^2[x-m(\tau)]+m'(\tau).
\]

\section{Diffusion loading}
The one-factor diffusion scale is
\[
g(\tau)=aH(\tau).
\]
Hence
\[
\operatorname{Cov}\bigl(dx(t,\tau_1),dx(t,\tau_2)\bigr)
=
2a^2H(\tau_1)H(\tau_2)\,dt,
\]
and
\[
K(\tau_1,\tau_2)=2a^2H(\tau_1)H(\tau_2).
\]

\section{The \texorpdfstring{\(r\)}{r}-branch}
Let \(a^2\) be the temporal drift-diffusion scale and \(\lambda^2\) the maturity-decay scale. Define
\[
r=\frac{\lambda^2}{a^2}>0.
\]
For \(r\neq1\), the admissible modes are
\[
e^{-a^2\tau},
\qquad
 e^{-ra^2\tau}.
\]
Writing \(k=a^2\), these are \(e^{-k\tau}\) and \(e^{-rk\tau}\). At \(r=1\), the roots coalesce and the repeated-root space is spanned by
\[
e^{-k\tau},
\qquad
\tau e^{-k\tau}.
\]

\section{Distinct-root Dirichlet loading}
For \(r_D\neq1\), the distinct-root Dirichlet loading is
\[
H_D^{(r)}(\tau)
=
\frac{r_D}{1+r_D}(x_D-x_\infty)
\left(e^{-r_Da_D^2\tau}-e^{-a_D^2\tau}\right).
\]
The corresponding variance rate is
\[
\frac{\partial}{\partial t}\operatorname{Var}(x^D_{t,t+\tau})
=
2a_D^2\left[H_D^{(r)}(\tau)\right]^2.
\]
Thus
\[
\frac{\partial}{\partial t}\operatorname{Var}(x^D_{t,t+\tau})
=
2a_D^2\left(\frac{r_D}{1+r_D}\right)^2(x_D-x_\infty)^2
\left(e^{-r_Da_D^2\tau}-e^{-a_D^2\tau}\right)^2.
\]

\section{The repeated-root Dirichlet limit}
The case \(r_D=1\) is a coalescing-root case. Consider
\[
\frac{e^{-r_Da_D^2\tau}-e^{-a_D^2\tau}}{1-r_D}.
\]
As \(r_D\to1\),
\[
\frac{e^{-r_Da_D^2\tau}-e^{-a_D^2\tau}}{1-r_D}
\longrightarrow
 a_D^2\tau e^{-a_D^2\tau}.
\]
Therefore the non-trivial repeated-root Dirichlet loading is
\[
H_D^{(1)}(\tau)=\bar C_D\tau e^{-a_D^2\tau}.
\]
The scalar coefficient is a normalisation convention: the limiting factor \(a_D^2\) and any sign inherited from the orientation of the distinct-root basis are absorbed into \(\bar C_D\). The boundary-selected object is the repeated-root mode shape \(\tau e^{-a_D^2\tau}\), not a separate sign convention.

\section{Dirichlet hump and variance rate}
For
\[
H_D^{(1)}(\tau)=C_D\tau e^{-a_D^2\tau},
\]
the variance rate is
\[
\frac{\partial}{\partial t}\operatorname{Var}(x^D_{t,t+\tau})
=
2a_D^2C_D^2\tau^2e^{-2a_D^2\tau}.
\]
Ignoring scale, let
\[
h_D(\tau)=\tau e^{-a_D^2\tau}.
\]
Then
\[
h_D'(\tau)=e^{-a_D^2\tau}(1-a_D^2\tau),
\]
so the maximum occurs at
\[
\tau_{\max}=\frac{1}{a_D^2}.
\]

\section{Shifted inverse-gamma density}
Write the field relative to a lower admissible level:
\[
y(t,\tau)=x(t,\tau)-x_{\min},
\qquad y(t,\tau)\ge0.
\]
For each fixed \(\tau>0\), suppose the shifted Dirichlet dynamics remain in the multiplicative positive-state class
\[
dy_t
=
\kappa_D(\tau)\bigl(\theta_D(\tau)-y_t\bigr)\,dt
+
\sqrt{2}\,\eta_D(\tau)y_t\,dW_t.
\]
The stationary KFE admits the inverse-gamma density
\[
p_D(y;\tau)
=
\frac{\beta_D(\tau)^{\alpha_D(\tau)}}{\Gamma(\alpha_D(\tau))}
 y^{-\alpha_D(\tau)-1}
\exp\left(-\frac{\beta_D(\tau)}{y}\right),
\qquad y>0,
\]
where
\[
\alpha_D(\tau)=1+\frac{\kappa_D(\tau)}{\eta_D^2(\tau)},
\qquad
\beta_D(\tau)=\frac{\kappa_D(\tau)\theta_D(\tau)}{\eta_D^2(\tau)}.
\]
Equivalently, in the original variable,
\[
x>x_{\min}.
\]
This is a multiplicative positive-state diffusion, not a CIR square-root diffusion; the CIR Feller condition does not apply.

\section{Feller inaccessibility}
For fixed \(\tau>0\), the diffusion coefficient is
\[
\sigma_D(y,\tau)=\sqrt{2}\,\eta_D(\tau)y,
\]
so \(\sigma_D(0,\tau)=0\). The drift at the lower boundary is
\[
b_D(0,\tau)=\kappa_D(\tau)\theta_D(\tau)>0.
\]
Thus the lower boundary \(y=0\) is inaccessible for each fixed \(\tau>0\), provided the process starts in the positive state domain and the positive-state parameter conditions hold. To lift this to
\[
y(t,\tau)\ge0
\qquad\text{for all }\tau\ge0,
\]
one also requires sufficient regularity in \(\tau\). At \(\tau=0\), Dirichlet closure imposes
\[
y(t,0)=0.
\]

\section{Neumann/free benchmark}
The Neumann/free benchmark retains the ordinary exponential stochastic mode:
\[
H_N(\tau)=A_Ne^{-a_N^2\tau}.
\]
The short-end loading is
\[
H_N(0)=A_N.
\]
The variance rate is
\[
\frac{\partial}{\partial t}\operatorname{Var}(x^N_{t,t+\tau})
=
2a_N^2A_N^2e^{-2a_N^2\tau}.
\]
Thus stochastic exposure is retained at the short end whenever \(A_N\neq0\). This is the non-pinned benchmark against which Dirichlet and Robin anchoring are compared.

\section{Robin loading and anchoring parameter}
The Robin repeated-root loading is
\[
H_R(\tau)=(C_0+C_1\tau)e^{-a_R^2\tau}.
\]
At the boundary,
\[
H_R(0)=C_0,
\qquad
H_R'(0)=C_1-a_R^2C_0.
\]
The anchoring parameter is
\[
\rho=-\frac{H_R'(0)}{H_R(0)}
=a_R^2-\frac{C_1}{C_0},
\]
whenever \(C_0\neq0\). Hence
\[
H_R(\tau)=C_0\bigl[1+(a_R^2-\rho)\tau\bigr]e^{-a_R^2\tau}.
\]
Writing \(k=a_R^2\), this becomes
\[
H_R(\tau)=C_0[1+(k-\rho)\tau]e^{-k\tau}.
\]

\section{Robin sign compatibility}
For \(C_0>0\), the Robin loading
\[
H_R(\tau)=C_0[1+(a_R^2-\rho)\tau]e^{-a_R^2\tau}
\]
is non-negative on the whole half-line if and only if
\[
\rho\le a_R^2.
\]
If \(\rho>a_R^2\), the loading changes sign at sufficiently long maturity and no longer supports a simple globally non-negative lower-boundary interpretation.

\section{Summary}
The appendix establishes the algebraic spine of the monograph:
\[
dx=-f\,dt+\sqrt{2}g\,dW_t,
\]
\[
f(x,t,\tau)=a^2[x-m(\tau)]-m'(\tau),
\]
\[
g(\tau)=aH(\tau),
\]
\[
r=\frac{\lambda^2}{a^2}.
\]
For \(r=1\), the repeated-root mode space contains
\[
e^{-a^2\tau},
\qquad
\tau e^{-a^2\tau}.
\]
Boundary closure selects the realised loading:
\[
H_N(\tau)=A_Ne^{-a_N^2\tau},
\]
\[
H_D^{(1)}(\tau)=C_D\tau e^{-a_D^2\tau},
\]
\[
H_R(\tau)=(C_0+C_1\tau)e^{-a_R^2\tau}.
\]
The humped Dirichlet loading is therefore not imposed as a primitive volatility input. It is the boundary-selected coalescing-root component of the admissible IRFR mode space.

\end{document}
